% r03_fill_K5a: proof of (K_0), part VII: the inputs of the forward bound (kernel bounds of the hybrid flow, the two stability bounds, the polynomial lower bound at the wall).

\subsubsection*{The inputs of the forward bound}

Throughout $T\in\{T_1,T_2\}$ and $r=r(k)$ is the choice of \Cref{rot:lem:r}; $p(\delta)=b_0(1+\log\delta^{-1})^{\mathrm p}$ is King's small-field function.

\begin{proposition}[kernel bounds of the hybrid flow at a non-zero gauge field]\label{rot:prop:BD}
Let $r\ge2$ and let $A$ be regular in the sense of \cite[Def.~3.2]{K1} on $\Lambda_X$ or on $\Lambda_X/P$ for a torus $P$ fitted by both lattices, with $e_r^\beta\le c_0(a)$.
\begin{enumerate}[label=\textup{(\alph*)},leftmargin=*]
\item \textup{(kernel bounds)} The operators $G_r$, $\Delta^{(r)}$, $C^{(r)}$ of the hybrid flow with the paired contours of \Cref{rot:def:flows} obey the bounds of \cite[Thm.~3.3, (3.5)--(3.8)]{K1} in the periodic form (on the infinite lattice and on tori fitted by both lattices, each in its own quotient distance), with constants depending on $a$, $M$, $\alpha$ only, except the H\"older constant of (3.8), which is $C_\alpha\log^{3\alpha}(1/\eta')$ (\Cref{rot:hyb:rem:Hol}); the differences between two tori, or between a torus and the infinite lattice, with a common fundamental region $F$ that is a union of $(r+1)$-blocks on which the backgrounds agree up to a collar, obey the same bounds with the distance taken in the smaller $d_\wedge$ of the two quotient metrics and the additional factor $e^{-\delta(\dist(\{x,y\},\partial F)+\dist(\supp f,\partial F))}$ for $G_r$ and $e^{-\delta(\dist(z,\partial F)+\dist(w,\partial F))}$ for $\Delta^{(r)}$ and $C^{(r)}$ (\Cref{rot:hyb:lem:domdiff}), in place of the Neumann corrections; the next-step precision $\Delta^{(r)}+a5^{-2}\hat Q^*\hat Q$ obeys $(25\gamma_0^{-1}+a_r^{-1})^{-1}\le\Delta^{(r)}+a5^{-2}\hat Q^*\hat Q\le a_r+a/25$, with $\gamma_0=\gamma_0(a)$ the coercivity constant of \Cref{rot:hyb:lem:coer}, and has an exponentially decaying kernel, and so has its restriction to every subset $\Lambda$ of the coarse lattice, in particular to every union of big blocks, whose inverse $\hat C^{(r)}_\Lambda$ of \cite[(5.53)]{K2} obeys \cite[(5.54)]{K2} and the two-torus clause (\Cref{rot:hyb:prop:CLam}); the fine-slice lines obey \cite[(3.63)--(3.65)]{K1} in the covariant form at regular $A$ (covariant lattice derivatives and transported H\"older quotients, \cite[(2.10)--(2.11)]{Ba3}) and in the plain form at the reduced background of step (5) of the proof of \Cref{rot:thm:A}; and the level-$r$ elements obey $|H^X(x,z)|,|D_AH^X(x,z)|\le Ce^{-\delta_0|x-z|}$ with the H\"older clause, $|B^X(z,w)|=|(A^X)^{-1}(z,w)|\le C(1+|z-w|)^{-1}$ and $|\Gamma^X(x,y)|\le C(1+|x-y|)^{-1}$ with the fine-variable regularity of $H^X$, $B^X$ and $\Gamma^X$ decaying exponentially only at the coarse mass; all constants are uniform in $r$, $k$ and the torus. The same holds at every level $j\le r$ with constants uniform in $j$, and, in the form in which the hybrid lower bound uses it, at the levels $j\in\{0,1\}$. The same statements hold for the standard flow with King's uniform constants (\cite[Thm.~3.3]{K1}, proved in \cite{Ba4}).
\item \textup{(gains)} After the constant-background reduction of \Cref{rot:lem:CB} the difference elements $\Delta\Gamma\defeq\Gamma^X-\Gamma^Y\circ\tau$, $\Delta H$, $\Delta(A^{-1})$ at regular $A$ obey the bounds of \Cref{rot:cor:unif} with the factor $\eta'^{1/2}$; the contour vertices \cite[(3.24)--(3.26)]{K1} at the reduced background differ entrywise by a relative $\eta^{1/2}$ against their size $Ce_rp_r^2$; the constant-background phases \cite[(3.46)]{K1} at matched vertices differ by $\le e\,p_r\lmu^{-1}c_\tau\delta_r^{-1/2}\eta'$; and the per-slice pieces carried by the contour vertex \cite[(3.26)--(3.27)]{K1}, with one end or both ends on contours, obey the slice-$j$ bounds \cite[(3.63)--(3.65)]{K1} in both flows with a relative mismatch $(\eta'/\ell_j)^{1/2}$, and, for a one-ended line with one lattice derivative at its lattice end, with the relative mismatch $(\eta'/\ell_j)^{\theta_\partial}$, $\theta_\partial\defeq\alpha/(2(2+\alpha))\ge\theta_1$ with $\alpha$ the H\"older exponent of \cite[(3.8)]{K1} (\Cref{rot:hyb:lem:DS}).
\end{enumerate}
\end{proposition}
At $A=0$ the proposition is the propagator package (\Cref{rot:prop:fc,rot:lem:H,rot:lem:grad,rot:lem:Gam,rot:lem:slice,rot:cor:unif,rot:cor:C}). At $A\ne0$ the proof, which occupies the rest of this subsection, is the programme of \cite[\S5.3]{K2} written out: coercivity of $-\Delta_A+m_r^2+a_rQ(A)^*Q(A)$ at a regular background by a cover with axis cubes carrying a local gauge for the constant part of the field (the tile properties of \Cref{rot:lem:tile} and the paired contours of \Cref{rot:lem:pair} are what enters: blocks partition the lattice, have bounded diameter, carry contours of bounded length inside the block, and $QQ^*=\mathbf 1$ for every $A$, so that $Q^*Q$ is the projection onto the block-constant mode); Combes--Thomas decay in $L^2$ form \cite{CT73}; the domain-difference lemma for the two-torus clause; Ba{\l}aban's random walk \cite[\S2]{Ba4} with periodic cells fitted by both lattices (\cite[(5.55)--(5.56)]{K2}); the $L^p$ and H\"older bounds of \cite[Lemma~2.2]{Ba4} with the $A=0$ input, the slice clauses stated covariantly; and, for (b), King's constant-background reduction \cite[(3.43)--(3.55)]{K1} applied to both flows on their infinite lattices in the universal-cover form of \Cref{rot:lem:CB}, after which every difference element at matched points is of one of seven types (an $A=0$ line, with the gain $\eta'^{1/2}$ of \Cref{rot:cor:unif}; a contour line, with the kernel difference along the contour plus the flux of the slice kernel through the sliver of \Cref{rot:lem:flux}; a double contour line; a contour vertex; an external lifted field, \Cref{rot:lem:H}; a constant-background phase; a hopping vertex, \Cref{rot:lem:vert}, unchanged), the extension of the lattice kernels off the lattice being the cubic-Hermite cardinal interpolation in the first argument, which decays like the lattice kernel.

% r03_hyb_setup: proof of \Cref{rot:prop:BD}, part I: the hybrid block-spin operators at a gauge field; lengths of the contour pieces; the block-spin algebra for any contour system.
% Input from sec/r03_fill_K5a.tex after the statement of \Cref{rot:prop:BD}. Parts II-VII follow as sec/r03_hyb_coer.tex, r03_hyb_domdiff.tex, r03_hyb_A0.tex, r03_hyb_walk.tex, r03_hyb_ES.tex, r03_hyb_gain.tex.

\paragraph{The hybrid operators at a gauge field.}
Throughout the proof of \Cref{rot:prop:BD} we work in the unit scale of level $r\ge2$ (\cite[(2.20)]{K1}): the physical spacing $\delta_r=5^r\eps_k$ is the unit, the fine spacing is $\eta'=5^{-r}$, $\eta=5\eta'$, and $\ell_j\defeq5^j\eta'$ ($\ell_r=1$). We write $L\defeq5$ for the block factor where a formula is clearer with it, $d=3$, and $f_\alpha=R^{-1}e_\alpha$ for the frame of $\Yk$. The coarse lattice $\Gamma_r=R^{-1}\Z^3$ is common to both flows; $\Gamma_j\defeq R^{-1}(\ell_j\Z^3)$ for $0\le j\le r$, so that $\Gamma_0=\Lambda_Y=R^{-1}(\eta'\Z^3)$ is the fine lattice of the flow $Y$ and $\Gamma_1=R^{-1}(\eta\Z^3)$ is $Y_{k-1}$ rescaled; $\Lambda_X=\eta'\Z^3$ is the fine lattice of the flow $X$ (\Cref{rot:def:match}). Since $5R^{-1}$ is an integer matrix (\Cref{rot:lem:geom}(a)), $\Gamma_1\subset\Lambda_X\cap\Lambda_Y$.

For $1\le j\le r$ and $z\in\Gamma_j$ the \emph{level-$j$ block} of the flow $F$ is
\[
B^F_j(z)\defeq\Lambda_F\cap\big(z+R^{-1}(-\tfrac{\ell_j}2,\tfrac{\ell_j}2)^3\big).
\]
For $F=Y$ this is a cube of $5^{3j}$ points aligned with $(f_\alpha)$; for $F=X$ it is the digitised rotated cube, with no lattice point on its boundary (\Cref{rot:lem:tile}(i)), and $B^X_1(z)=z+\eta'B_0$ is the tile. Both are nested, $B^F_{j+1}(z)=\bigcup_{s\in\{-2,\dots,2\}^3}B^F_j(z+\ell_jR^{-1}s)$, and for $x\in B^F_r(z)$ we write $x_j=x_j(x)$ for the centre of the level-$j$ block containing $x$ ($x_r=z$). Blocks are centred at the sites of the next coarse lattice (\Cref{rot:conv:C}(i)): the root of every block is its centre, in both flows. The matching $\tau:\Lambda_X\to\Lambda_Y$ of \Cref{rot:def:match} maps the tile $B^X_1(x)$ onto the cube $B^Y_1(x)$ for every $x\in\Gamma_1$, with $|\tau x'-x'|\le c_\tau\eta'$, $c_\tau=\sqrt{13}+2\sqrt3<7.1$; points $x''=\tau x'$ are matched. The composite contours are those of \Cref{rot:def:flows}: for the flow $Y$ King's \cite[(2.12)]{K1}, $\Gamma^{Y,(r)}_{z,x}=\Gamma^Y_{z,x_{r-1}}\cup\dots\cup\Gamma^Y_{x_1,x}$, where $\Gamma^Y_{u,v}$ is the L-path that changes the $f_3$-coordinate first, then $f_2$, then $f_1$; for the flow $X$, $\Gamma^{X,(r)}_{z,x'}=\Gamma^X_{z,x_{r-1}}\cup\dots\cup\Gamma^X_{x_2,x_1}\cup\Gamma^X_{x_1,x'}$, where for $i\ge1$ the level-$i$ piece $\Gamma^X_{x_{i+1},x_i}$ is the L-path $\Gamma^Y_{x_{i+1},x_i}$ with every $\eta$-bond $(u,u+\eta f_\alpha)$, $u\in\Gamma_1$, replaced by the fixed staircase $u+\eta'\gamma_\alpha$ (the reversed bond by the reversed staircase), and $\Gamma^X_{x_1,x'}$ is the fixed in-tile path of at most $5$ bonds. The gauge field along a composite contour is read in King's convention \cite[(2.19)]{K1} (\Cref{rot:conv:C}(ii)).

\emph{Backgrounds.} A background is a real bond field $A$ on $\Lambda_F$, identified with a site vector field by $A_{(x,x+\eta'f_\alpha)}=A_\alpha(x)$ (\cite[p.~652]{K1}); $U(A(\Gamma))\defeq\exp[\mathrm ie_rA(\Gamma)]$ with $A(\Gamma)=\sum_{b\in\Gamma}\eta'A_b$ (oriented) and the running charge $e_r=e\,\delta_r^{1/2}$ (\cite[(3.22)]{K1}). \emph{Regular} means \cite[Def.~3.2]{K1}: $|\partial_\mu A_\nu(x)|\le Ce_r^{\beta-1}$ for all $x,\mu,\nu$ and some $\beta>0$. The small-field class at level $r$ is \cite[(3.2)]{K1}: a coarse field $A_r$ on $\Gamma_r$ with $|A_r|\le p_r(\lmu\delta_r)^{-1}$, $|\partial A_r|\le p_r$, $p_r=p(\delta_r)$, lifted to the flow $F$ by \cite[(2.18)]{K1}, $A^{(r),F}\defeq a_rG^{F,\mathrm{ph}}_rQ_F^*A_r$ (the photon fluctuation propagator and the frame-covariant photon average of \Cref{rot:rem:frameav}), plus the Cameron--Martin shift $za_F$ of \Cref{rot:lem:CM}. Both are regular for $\delta_r\le\delta_0$ (\cite[(3.45)]{K1}; \Cref{rot:lem:CM}).

\emph{The pinned operators.} For a flow $F$ and a background $A$ on $\Lambda_F$,
\[
Q_F(A)\phi(z)\defeq\sum_{x\in B^F_r(z)}\eta'^3\,U(A(\Gamma^{F,(r)}_{z,x}))^{-1}\phi(x)\quad(z\in\Gamma_r),\qquad
M_A\defeq-\Delta_A+m_r^2,\qquad G^F_r(A)\defeq\big(M_A+a_rQ_F(A)^*Q_F(A)\big)^{-1},
\]
with $-\Delta_A=D_A^*D_A$ the covariant Laplacian of the fine lattice, $(D_A\phi)(b)=\eta'^{-1}(U(-A_b)\phi(b_+)-\phi(b_-))$, $m_r=m\delta_r$, $a_r=a(1-5^{-2})(1-5^{-2r})^{-1}$, and the pairings $\ip uv=\sum\eta'^3\bar uv$ on $\Lambda_F$ and $\sum_z\bar uv$ on $\Gamma_r$ (unit block volume, $|B^F_r|\eta'^3=1$). This is \cite[(2.11), (2.13)]{K1} in unit scale; the convention $U^{-1}$ makes $Q_F(A)\phi$ transform with the coarse gauge function (\Cref{rot:hyb:lem:alg}). Put
\[
C_A\defeq M_A^{-1},\qquad u^F(A)\defeq C_AQ_F(A)^*,\qquad A^F(A)\defeq\big(a_r^{-1}+Q_F(A)C_AQ_F(A)^*\big)^{-1},\qquad B^F(A)\defeq A^F(A)^{-1},
\]
\[
H^F(A)\defeq a_rG^F_r(A)Q_F(A)^*=u^F(A)A^F(A),\qquad \Gamma^F(A)\defeq C_A-G^F_r(A)=u^F(A)A^F(A)u^F(A)^*=H^FB^FH^{F*}.
\]
By the Woodbury identity these relations are exact for every $A$. $A^F(A)$ is King's effective Laplacian $\Delta^{(r)}(A)=a_r-a_r^2Q_FG_rQ_F^*$ (\cite[(2.14)]{K1}; indeed $a_r(1-a_rQ_FG_rQ_F^*)=a_r(1+a_rQ_FC_AQ_F^*)^{-1}=A^F(A)$), $H^F$ is the lift operator of \cite[(2.18)]{K1} and \Cref{rot:lem:H}, and $C^{(r)}(A)=(A^F(A)+aL^{-2}Q^{(1)}(A)^*Q^{(1)}(A))^{-1}$ is the next fluctuation covariance (\cite[(2.16)]{K1}), $Q^{(1)}$ the one-step average of the coarse lattice. The photon kernels are three copies of the scalar kernels at $A=0$ and mass $\lmu\delta_r$ (\cite[p.~654]{K1}; \Cref{rot:rem:frameav}, \Cref{rot:prop:fc}): no photon kernel sees a background, so the photon sector of \Cref{rot:prop:BD} is the $A=0$ package of \Cref{rot:prop:fc,rot:lem:H,rot:lem:grad,rot:lem:Gam,rot:lem:slice,rot:cor:unif}.

\begin{lemma}[the pieces of the composite contours]\label{rot:hyb:lem:len}
Let $z\in\Gamma_r$, $x'\in B^X_r(z)$, $x''=\tau x'$, and let $x_i$ ($1\le i\le r$) be the common roots.
\begin{enumerate}[label=\textup{(\roman*)},leftmargin=*]
\item For $i\ge1$ the displacement is $x_i-x_{i+1}=\ell_iR^{-1}s$ with $s\in\{-2,\dots,2\}^3$, and every leg of $\Gamma^Y_{x_{i+1},x_i}$ is a concatenation of $|s_\alpha|5^{i-1}$ $\eta$-bonds between points of $\Gamma_1$. The level-$i$ piece has length $\le6\ell_i$ in the flow $Y$ and $\le\frac75\cdot6\ell_i=8.4\,\ell_i$ in the flow $X$; the level-$0$ pieces have length $\le6\eta'$ ($Y$) and $\le5\eta'$ ($X$). The partial contours $\Gamma^{F,(j+1)}_{x_{j+1},x}$ of \cite[(2.19)]{K1} (the pieces of levels $j,j-1,\dots,0$) have length $\le7.5\,\ell_j$ ($Y$) and $\le10.5\,\ell_j$ ($X$); the composite contours have length $\le1.5$ ($Y$) and $\le2.1$ ($X$).
\item Every level-$i$ piece lies inside the open rotated cube $x_{i+1}+R^{-1}(-\frac{\ell_{i+1}}2,\frac{\ell_{i+1}}2)^3$, hence inside $B^F_{i+1}(x_{i+1})$, and the level-$0$ piece inside its first-step cell; so every piece lies in its block and $\Gamma^{F,(r)}_{z,x}\subset B^F_r(z)$.
\item The two composite contours $\Gamma^{X,(r)}_{z,x'}$ and $\Gamma^{Y,(r)}_{z,x''}$ pass through the same points $x_{r-1},\dots,x_1$ and through the same endpoints of $\eta$-bonds; their level-$i$ pieces are within Hausdorff distance $\eta'$ for $i\ge1$, the level-$0$ pieces lie in the ball of radius $\sqrt{13}\eta'$ about $x_1$, and the closed polygon $\Gamma^{X,(r)}_{z,x'}\cup[x',x'']\cup(\Gamma^{Y,(r)}_{z,x''})^{-1}$ bounds a piecewise planar surface $\Sigma(z,x')$, a union of the $\eta$-bond slivers of \Cref{rot:lem:pair} and one polygon inside that ball, of area $\le1.5\eta'+33\eta'^2$. The sub-polygon formed by the partial contours up to level $j+1$ and the matching segment bounds a surface of area $\le1.5\eta'\ell_{j+1}+33\eta'^2$.
\end{enumerate}
\end{lemma}
\begin{proof}
(i) $x_i=\ell_iR^{-1}m_i$ with $m_i\in\Z^3$ the frame coordinates of the block centre and, by nesting, $m_i-5m_{i+1}=:s\in\{-2,\dots,2\}^3$ (the level-$i$ block of $x$ is one of the $125$ sub-blocks of its level-$(i+1)$ block; there are no ties since no lattice point lies on a block face). The L-path from $x_{i+1}$ to $x_i$ has legs $\ell_is_\alpha f_\alpha=5^{i-1}s_\alpha\cdot\eta f_\alpha$, each a run of $|s_\alpha|5^{i-1}$ $\eta$-bonds starting at a point of $\Gamma_1$ (because $x_{i+1}\in\Gamma_{i+1}\subset\Gamma_1$ and $\eta f_\alpha\in\Gamma_1$). Lengths: $\sum_\alpha|s_\alpha|\ell_i\le6\ell_i$; the staircases $\gamma_1,\gamma_2$ have $7$ fine steps for an $\eta$-bond of length $5\eta'$ and $\gamma_3$ has $5$; the in-tile piece has at most $5$ bonds (\Cref{rot:lem:tile}(ii)) and the level-$0$ $Y$-piece at most $6$. Summing, $\sum_{i=1}^{j}8.4\ell_i+5\eta'\le8.4\cdot\frac54\ell_j(1-5^{-j})+5\eta'<10.5\,\ell_j$ and $6\sum_{i=0}^{j}\ell_i\le7.5\,\ell_j$; for $j=r-1$ this gives $2.1$ and $1.5$ (\Cref{rot:lem:pair}). (ii) The endpoints of every $\eta$-bond of $\Gamma^Y_{x_{i+1},x_i}$ are points of $\Gamma_1$ inside the open cube $x_{i+1}+R^{-1}(-\frac{\ell_{i+1}}2,\frac{\ell_{i+1}}2)^3$; in the frame coordinates $(f_\alpha\cdot(\,\cdot\,-x_{i+1}))_\alpha$ they lie in $\eta\Z^3\cap(-\frac{\ell_{i+1}}2,\frac{\ell_{i+1}}2)^3$, and since $\frac{\ell_{i+1}}2=\frac{5^{i+1}}2\eta'$ is an odd multiple of $\frac{\eta'}2$ while the points are multiples of $5\eta'$, each coordinate is at distance $\ge\frac52\eta'$ from the faces. The segment between two such points keeps that distance, and the staircase deviates from the segment by $\le0.6\eta'<\frac52\eta'$ (\Cref{rot:lem:pair}), so it stays in the open cube; its points lie in $\Lambda_X$, hence in $B^X_{i+1}(x_{i+1})$. The in-tile piece lies in the tile (\Cref{rot:lem:tile}(ii)), which is a sub-block. (iii) The Hausdorff distances and the containment of the level-$0$ pieces are \Cref{rot:lem:pair} ($|x'-x_1|\le\sqrt{13}\eta'$, $|x''-x_1|\le2\sqrt3\eta'$, both level-$0$ pieces monotone inside their cells); the two composite contours pass through the same points $x_{r-1},\dots,x_1$ and through the same $\eta$-bond endpoints, so the closed polygon decomposes into the $\eta$-bond slivers (area $\le5\eta'^2$ each, at most $1.5/\eta=0.3/\eta'$ of them, total $\le1.5\eta'$) and the polygon $\Gamma^X_{x_1,x'}\cup[x',x'']\cup(\Gamma^Y_{x_1,x''})^{-1}$ inside the ball of radius $\sqrt{13}\eta'$, of perimeter $\le(5+7.1+6)\eta'$, which the cone from $x_1$ spans with area $\le\frac12\cdot18.1\cdot\sqrt{13}\,\eta'^2<33\eta'^2$. For the partial contours up to level $j+1$ the $Y$-length is $\le7.5\ell_j=1.5\ell_{j+1}$, so there are $\le1.5\ell_{j+1}/\eta$ $\eta$-bonds and the slivers have total area $\le5\eta'^2\cdot1.5\ell_{j+1}/\eta=1.5\eta'\ell_{j+1}$.
\end{proof}

\begin{lemma}[block-spin algebra for any contour system]\label{rot:hyb:lem:alg}
For every background $A$ on $\Lambda_X$ and every system of contours $\Gamma_{z,x}\subset B^X_r(z)$ from the root $z$ to the point $x$:
\begin{enumerate}[label=\textup{(\roman*)},leftmargin=*]
\item $Q_X(A)Q_X(A)^*=\mathbf 1$, and $P_X(A)\defeq Q_X(A)^*Q_X(A)=\sum_zw_zw_z^\dagger$ with $w_z(x)\defeq U(A(\Gamma_{z,x}))\mathbf 1_{B^X_r(z)}(x)$ is an orthogonal projection ($\|w_z\|=1$, orthogonal supports);
\item for a gauge function $\lambda$ on $\Lambda_X$, $Q_X(A+\partial\lambda)(e^{\mathrm ie_r\lambda}\phi)=e^{\mathrm ie_r\lambda|_{\Gamma_r}}Q_X(A)\phi$, hence $G^X_r(A+\partial\lambda)=e^{\mathrm ie_r\lambda}G^X_r(A)e^{-\mathrm ie_r\lambda}$, and likewise for $H^X$, $A^X$, $B^X$, $\Gamma^X$, $C^{(r)}$ with the coarse phases;
\item for a constant field $A\equiv B_0$, $U(B_0(\Gamma_{z,x}))=e^{\mathrm ie_rB_0\cdot(x-z)}$ for every contour, and $G^X_r(B_0)(x,y)=e^{\mathrm ie_rB_0\cdot x}G^X_r(0)(x,y)e^{-\mathrm ie_rB_0\cdot y}$ on $\Lambda_X$ (on a torus: the twisted periodisation of \Cref{rot:hyb:lem:A0}(c));
\item with the contours of \Cref{rot:def:flows}, $Q_X(A)=Q^{\mathrm{rot}}_{r-1}(A)\circ\hat Q(A)$ under the convention \cite[(2.19)]{K1}, and the composite blocks are unions of tiles (\Cref{rot:lem:tile}(iv)).
\end{enumerate}
\end{lemma}
\begin{proof}
(i) is \Cref{rot:lem:tile}(iii),(iv): $|U|=1$ and $|B^X_r(z)|\eta'^3=1$. (ii) $A(\Gamma_{z,x})\mapsto A(\Gamma_{z,x})+\lambda(x)-\lambda(z)$ for any path from $z$ to $x$, so $U(A(\Gamma_{z,x}))^{-1}e^{\mathrm ie_r\lambda(x)}=e^{\mathrm ie_r\lambda(z)}U(A(\Gamma_{z,x}))^{-1}$, and $-\Delta_{A+\partial\lambda}=e^{\mathrm ie_r\lambda}(-\Delta_A)e^{-\mathrm ie_r\lambda}$. (iii) A constant field is $\partial(B_0\cdot x)$ on $\Lambda_X$. (iv) The level-$j$ pieces are the contours of the steps of the flow $Y$ with the bonds re-drawn; iterating the one-step transformations produces exactly the phases of \cite[(2.19)]{K1} along the composite contour (\Cref{rot:conv:C}(ii); \cite[(3.33)--(3.34)]{Ba1}).
\end{proof}

The pairing of the contours is a property of \Cref{rot:def:flows}, not a further hypothesis; a hybrid transformation whose first-step contour were a fine-axis shortest path from the level-$r$ root across the whole block would be a different transformation, for which the contour phases of the two flows differ at order one (\Cref{rot:lem:flux}).

% r03_hyb_coer: proof of \Cref{rot:prop:BD}, part II: coercivity at a regular background, L^2 bounds, Combes--Thomas decay, the decay of C^{(r)}.

\paragraph{Coercivity and decay in $L^2$ form.}
The proof of \Cref{rot:prop:BD}(a) follows \cite[\S2]{Ba4} and \cite[\S5.3]{K2} with two changes: the lower bound on the quadratic form is proved directly at a regular background by a cover with axis cubes (so that no small-domain reduction is needed for the $L^2$ theory), and the local cells of the random walk are periodic cells fitted by both lattices (\cite[(5.55)]{K2}), on which the $A=0$ kernels are the twisted periodisations of the explicit kernels of \Cref{rot:lem:H,rot:lem:slice}. On the tori of the chain the periods lie in $10\Z^3$ in the unit scale of every level $r\le k$ ($P_J=50MJ\,5^{k-r}\Z^3$ and $R^{-1}P_J=10MJ\,5^{k-r}(5R^{-1})\Z^3$, \Cref{rot:lem:adm}(iii)), with sup-systole $s_\infty(P)\ge50$ for $r\le k-1$.

\begin{lemma}[coercivity at a regular background]\label{rot:hyb:lem:coer}
Let $\Omega$ be $\Lambda_X$ or $\Lambda_X/P$ with $P$ fitted by both lattices and $s_\infty(P)\ge10$ in unit scale. There are $c_0(a)>0$ and $\gamma_0(a)>0$ such that for every regular $A$ on $\Omega$ with $e_r^\beta\le c_0$ and every $\phi$,
\[
\ip{\phi}{(-\Delta_A+m_r^2+a_rP_X(A))\phi}\ \ge\ \gamma_0\|\phi\|^2 .
\]
Explicitly $\gamma_0=\frac1{2m_c}\min\{\frac18,\frac{a'}{64(8a'+1)}\}$ with $a'=a(1-5^{-2})\le a_r\le a$ and $m_c$ the multiplicity of the cover below; $m_c=8$ when $P\subset2\Z^3$, as on the tori of the chain.
\end{lemma}
\begin{proof}
\emph{Cover.} Let $\{\Delta_i\}$ be the axis cubes $y_i+[0,4)^3$ with $y_i\in\Lambda_c\defeq2\Z^3+P$ ($P=0$ on $\Lambda_X$); $\Lambda_c$ is $P$-invariant, so the family descends to a finite family on the torus, and a cube does not wrap since every non-zero period has $|\gamma|_\infty\ge10$. Every point, every bond and every block lies in at most $m_c$ cubes, $m_c=8$ if $P\subset2\Z^3$ and $m_c\le64$ in general ($\Lambda_c\subset\Z^3$). Every $\Delta_i$ contains a full block: the covering radius of $\Gamma_r=R^{-1}\Z^3$ is $\frac{\sqrt3}2$, so some $z_i\in\Gamma_r$ has $|z_i-c_i|\le\frac{\sqrt3}2$ ($c_i$ the centre of $\Delta_i$), and $B^X_r(z_i)\subset\{|x-z_i|<\frac{\sqrt3}2\}\subset\{|x-c_i|<\sqrt3\}\subset\Delta_i$ because $\sqrt3<2$. Dropping the non-negative terms of the quadratic form that belong to bonds joining different cubes, to the mass and to the pins of the other blocks, and dividing by the multiplicity,
\[
\ip{\phi}{(-\Delta_A+m_r^2+a_rP_X(A))\phi}\ \ge\ \tfrac1{m_c}\sum_i\Big[\sum_{b\subset\Delta_i}\eta'^3|D_A\phi(b)|^2+a_r|\ip{w_{z_i}}{\phi}|^2\Big]=:\tfrac1{m_c}\sum_iq_i(\phi).
\]
\emph{Local gauge.} On the cube $\Delta_i$ (contractible) write $A=A_0^{(i)}+A'^{(i)}$ with $A_0^{(i)}\defeq A(c_i)$ constant and, by regularity along lattice paths of $\ell^1$-length $\le12$, $|A'^{(i)}|\le12\,Ce_r^{\beta-1}$ on $\Delta_i$. The gauge function $\lambda_i(x)=A_0^{(i)}\cdot x$ removes $A_0^{(i)}$ from both terms of $q_i$ (\Cref{rot:hyb:lem:alg}(ii),(iii), applied to the bonds and the block inside $\Delta_i$), so $q_i(\phi)=q_i^{A'}(\tilde\phi)$ with $\tilde\phi=e^{-\mathrm ie_r\lambda_i}\phi$, $|\tilde\phi|=|\phi|$.
\emph{Perturbation.} $D_{A'}\phi(b)=U(-A'_b)D_0\phi(b)+\eta'^{-1}(U(-A'_b)-1)\phi(b_-)$, hence $\|D_{A'}\phi\|_{\Delta_i}\ge\|D_0\phi\|_{\Delta_i}-e_r|A'|_\infty\|\phi\|_{\Delta_i}$ and $\|D_{A'}\phi\|^2\ge\frac34\|D_0\phi\|^2-3e_r^2|A'|_\infty^2\|\phi\|^2$. For the pin, $w_z(A')-w_z(0)$ has sup norm $\le e_r|A'(\Gamma)|\le2.1\,e_r|A'|_\infty$ (\Cref{rot:hyb:lem:len}(i)), so $\big||\ip{w_z(A')}{\phi}|^2-|\ip{w_z(0)}{\phi}|^2\big|\le(4.2e_r|A'|_\infty+4.41e_r^2|A'|_\infty^2)\|\phi\|_B^2$.
\emph{The free cube.} With $\psi\defeq\tilde\phi-c$ on $\Delta_i$, $c$ the mean, $\sum_{b\subset\Delta_i}\eta'^3|D_0\tilde\phi(b)|^2\ge\lambda_1\|\psi\|^2_{\Delta_i}$ with $\lambda_1=4\eta'^{-2}\sin^2(\pi\eta'/8)\ge\frac14$ (the first Neumann eigenvalue of the cube of $4/\eta'$ points per side; $\sin t\ge\frac2\pi t$), and, since $|B_i|\eta'^3=1$, $|\ip{w_{z_i}}{\tilde\phi}|^2=|c+\ip{\mathbf 1_{B_i}}{\psi}|^2\ge(1-t)|c|^2-(\frac1t-1)\|\psi\|^2_{B_i}$ for $0<t<1$. Choosing $\frac1t-1=\frac{\lambda_1}{2a_r}$,
\[
\lambda_1\|\psi\|^2+a_r|\ip{w_{z_i}}{\tilde\phi}|^2\ \ge\ \tfrac{\lambda_1}2\|\psi\|^2+\tfrac{a_r\lambda_1}{2a_r+\lambda_1}|c|^2\ \ge\ \gamma_1\|\tilde\phi\|^2_{\Delta_i},\qquad
\gamma_1\defeq\min\Big\{\tfrac{\lambda_1}2,\tfrac{a_r\lambda_1}{64(2a_r+\lambda_1)}\Big\},
\]
using $\|\tilde\phi\|^2_{\Delta_i}=\|\psi\|^2+64|c|^2$. \emph{Assembly.} $q_i(\phi)\ge\frac34\gamma_1\|\phi\|^2_{\Delta_i}-(3e_r^2|A'|^2+a_r(4.2e_r|A'|+4.41e_r^2|A'|^2))\|\phi\|^2_{\Delta_i}\ge\frac12\gamma_1\|\phi\|^2_{\Delta_i}$ once $e_r^\beta\le c_0(a)$, and $\sum_i\|\phi\|^2_{\Delta_i}\ge\|\phi\|^2$. With $\lambda_1\ge\frac14$, $\gamma_1\ge\min\{\frac18,\frac{a'}{64(8a'+1)}\}$.
\end{proof}

This is \cite[(5.59)--(5.60)]{K2} (``cover $\square$ by cubes $\{\Delta_i\}$ on the $\eta$-lattice with sides parallel to the coordinate planes \dots\ each cube contains fully at least one block'') written out, with the local gauge step that \cite[p.~345]{K2} leaves to ``$A=B+A'$ with $B$ constant and $A'$ small'' made explicit; it is also \cite[(1.8)]{Ba4} for the hybrid flow. No Poincar\'e inequality on the digitised rotated cube is needed: only axis cubes of the fine lattice enter.

\begin{corollary}[$L^2$ bounds]\label{rot:hyb:cor:L2}
Under the hypotheses of \Cref{rot:hyb:lem:coer}, $T\defeq-\Delta_A+m_r^2+a_rP_X(A)\ge\gamma_0$, so $\|G^X_r(A)\|\le\gamma_0^{-1}$, $\|D_AG^X_r(A)^{1/2}\|\le1$, $\|G^X_r(A)^{1/2}D_A^*\|\le1$, $\|D_AG^X_r(A)D_A^*\|\le1$ and $0\le G^X_r(A)\le\gamma_0^{-1}$ (\cite[Lemma~2.1]{Ba4}, \cite[Lemma~5.7]{K2}). The same holds, on $\Lambda_X$ and on a torus with $s_\infty(P)\ge10L$ in unit scale (true on the tori of the chain for $r\le k-1$), for the operator with the level-$(r+1)$ blocks, $-\Delta_A+m_r^2+a_{r+1}L^{-2}P_{r+1}(A)$ (the pin at scale $L$), with $\gamma_0$ replaced by $\gamma_0L^{-2}$.
\end{corollary}
\begin{proof}
$D_A^*D_A\le T$. For the $(r+1)$-blocks repeat the proof with cubes of side $4L$, corners in $2L\Z^3+P$ (multiplicity $8$ if $P\subset10\Z^3$, as for $P_J$ and $R^{-1}P_J$; periods with $|\gamma|_\infty\ge10L$) and blocks of $L^3$ unit cells (\cite[p.~588]{Ba4}: ``replace $\min\{\pi^2,a_k\}$ by $\min\{\pi^2L^{-2},a_{k+1}L^{-2}\}$'').
\end{proof}

\begin{proposition}[Combes--Thomas decay]\label{rot:hyb:prop:CT}
Under the hypotheses of \Cref{rot:hyb:lem:coer} there are $\theta_0,c_1>0$ depending on $a$ only such that for all $f,f'$ supported in unit cubes,
\[
|\ip f{G^X_r(A)f'}|,\ |\ip f{D_AG^X_r(A)f'}|,\ |\ip f{G^X_r(A)D_A^*f'}|\ \le\ c_1e^{-\theta_0\dist(\supp f,\supp f')}\|f\|_2\|f'\|_2 ,
\]
$\dist$ the quotient distance of $\Omega$. Hence (\cite[(3.6)]{K1}) $|A^X(A)(z,w)|=|a_r\delta_{zw}-a_r^2\ip{w_z}{G^X_r(A)w_w}|\le Ce^{-\theta_0|z-w|}$. The same holds for the operator with the $(r+1)$-pin of \Cref{rot:hyb:cor:L2}, with $\theta_0,c_1$ depending on $a$ and $L$.
\end{proposition}
\begin{proof}
For a weight $\psi$ with $|\psi(x)-\psi(x')|\le|x-x'|$ and $0\le t\le\theta_0$ put $T_\psi\defeq e^{t\psi}Te^{-t\psi}$. On $\Lambda_X$ one may take the linear weight $\psi=\hat n\cdot x$; on a torus $\vartheta\cdot x$ is not a function, and $\psi=\min\{\dist(\cdot,\supp f'),R\}$ is used, with the quotient distance; the estimates below use only the Lipschitz bound on bonds and inside blocks. The hopping part changes by a form-bounded perturbation: $e^{t\psi}D_Ae^{-t\psi}=D_A+B_t$, with $B_t$ the bond multiplication by $\eta'^{-1}(e^{-t(\psi(b_+)-\psi(b_-))}-1)$ composed with the covariant shift (bond functions weighted at $b_-$), of norm $\|B_t\|\le\sqrt3\,te^{t\eta'}$ (every site is the end $b_+$ of three bonds); with $e^{t\psi}D_A^*e^{-t\psi}=(D_A+B_{-t})^*$ this gives $|\ip\phi{(T_\psi-T)\phi'}|\le Ct(\|D_A\phi\|\|\phi'\|+\|\phi\|\|D_A\phi'\|)+Ct^2\|\phi\|\|\phi'\|$ for $t\eta'\le1$. The pin changes by $\sum_zw_z(x)\overline{w_z(y)}(e^{t(\psi(x)-\psi(y))}-1)$ with $|\psi(x)-\psi(y)|\le|x-y|\le\sqrt3$ inside a block, an operator of norm $\le2\sqrt3\,te^{\sqrt3t}$ (one rank-one piece per block, orthogonal supports). Hence $\|T^{-1/2}(T_\psi-T)T^{-1/2}\|\le Ct\gamma_0^{-1/2}+Ct^2\gamma_0^{-1}\le\frac12$ for $t\le\theta_0(a)$, using $\|D_AT^{-1/2}\|\le1$ (\Cref{rot:hyb:cor:L2}); so $T_\psi=T^{1/2}(1+T^{-1/2}(T_\psi-T)T^{-1/2})T^{1/2}$ is invertible with $\|T_\psi^{-1}\|\le2\gamma_0^{-1}$. Then $\ip f{G^X_rf'}=\ip{e^{-t\psi}f}{T_\psi^{-1}e^{t\psi}f'}$, since $G^X_r=e^{-t\psi}T_\psi^{-1}e^{t\psi}$; choose $t=\theta_0$ and $\psi=\min\{\dist(\cdot,\supp f'),R\}$ with $R\ge\dist(\supp f,\supp f')$ (or, on $\Lambda_X$, $\psi=\hat n\cdot x$ with $\hat n$ the unit vector from $\supp f'$ to $\supp f$), and use $\sup_{\supp f}e^{-t\psi}\sup_{\supp f'}e^{t\psi}=e^{-\theta_0\dist(\supp f,\supp f')}$ for the distance weight ($\psi=0$ on $\supp f'$), $\le e^{-\theta_0(\dist-2\sqrt3)}$ for the linear one. The derivative versions use $\|D_AT_\psi^{-1}\|,\|T_\psi^{-1}D_A^*\|\le C$ (the same algebra). For $A^X$: $w_z,w_w$ have unit norm and supports in blocks of diameter $\le\sqrt3$. The same proof with the $(r+1)$-pin (blocks of diameter $\le\sqrt3L$, \Cref{rot:hyb:cor:L2}) gives the bounds for that operator.
\end{proof}

\begin{corollary}[{\cite[(3.6)]{K1} for $C^{(r)}$}]\label{rot:hyb:cor:Ck}
$|C^{(r),X}(A)(z,w)|\le Ce^{-\delta_0d(z,w)}$ in the quotient distance of the domain. For two domains $\Omega_1$, $\Omega_2$ as in \Cref{rot:hyb:lem:domdiff} (a common fundamental region $F$, a union of $(r+1)$-blocks, on which the backgrounds agree up to a collar of width $W_0$) and $y,y'\in\Gamma_r\cap F$,
\[
|C^{(r)}(\Omega_1,A_1;y,y')-C^{(r)}(\Omega_2,A_2;y,y')|\le Ce^{-\delta(d_\wedge(y,y')+\dist(y,\partial F)+\dist(y',\partial F))},
\]
$d_\wedge$ the smaller of the two quotient distances; the same for $\Delta^{(r)}=A^X$.
\end{corollary}
\begin{proof}
\cite[\S3]{Ba4} writes $C^{(r)}(\Omega,A;y,y')$ through the Gaussian integral (3.1), factorises it over the next-level blocks after the gauge transformation ``$\psi(y)\to U(A(\Gamma_{y,z}))\psi(y)$'' (p.~587), and arrives at the identity (3.8): $C^{(r)}$ is a bounded combination of $\ip{q_r(y)}{G_r(\Omega,\Lambda,A)q_r(y')}$ and $\ip{q_{r+1}(y)}{\cdot\,q_{r+1}(y')}$, with $q_r(y)$, $q_{r+1}(y)$ the block modes (3.7) of unit and $L^{-d/2}$ norm, where $G_r(\Omega,\Lambda,A)$ is the fluctuation propagator pinned on the $(r+1)$-blocks (3.5). The two inputs of this algebra are that the blocks partition the lattice and that every contour runs from the block root to the point, both true for the hybrid flow (\Cref{rot:hyb:lem:alg}(i),(iv), \Cref{rot:lem:tile}); the decay of $G_r(\Omega,\Lambda,A)$ in the $L^2$-weighted form is \Cref{rot:hyb:prop:CT} for the $(r+1)$-pin. This is what \cite[pp.~345--346]{K2} says: ``To prove Proposition 5.6, we use the method presented in Chap.~3 of [Ba4]. This method requires only $L^2$-bounds on the operators $G_k(\Omega,A)$.'' The two-domain bound is not \cite[(3.9)--(3.10)]{Ba4}, which change the pin region inside one domain: on each domain (3.6) gives $C^{(r)}=V^*G_r(\Omega,\Omega^{(r)},A)V-(\beta/a_r)P+a_r^{-1}$ with $G_r(\Omega,\Omega^{(r)},A)$ the $(r+1)$-pinned operator (3.5), $V=Q_r^*(1-\beta P)$, $\beta=a_{r+1}L^{-2}/a_r$ ($Q_r=Q_X(A)$, $Q=Q^{(1)}(A)$, $P=Q^*Q$), whose local terms and block functions agree on the common region, so that $\delta C^{(r)}(y,y')=\ip{v(y)}{\delta G\,v(y')}$ there; $\delta G$ for the $(r+1)$-pin obeys the weighted bound \eqref{rot:hyb:eq:B1} of the geometric resolvent identity $\ip{f'}{(G_2-G_1)f}=\ip{f'}{G_1[H_2,\chi]G_2f}$ with \Cref{rot:hyb:prop:CT} on each side (\Cref{rot:hyb:lem:domdiff}).
\end{proof}

% r03_hyb_domdiff: proof of \Cref{rot:prop:BD}, part III: the domain-difference lemma (two tori, or a torus against the infinite lattice) and the restricted precision inverses on unions of big blocks.

\begin{lemma}[domain difference]\label{rot:hyb:lem:domdiff}
Let $\Omega_1=\Lambda_X/P_1$ be a torus fitted by both lattices with $s_\infty(P_1)\ge10L$ in unit scale, and $F$ a fundamental region of $P_1$ that is a union of $(r+1)$-blocks. Let $\Omega_2$ be $\Lambda_X$, or a torus $\Lambda_X/P_2$ fitted by both lattices into which $F$ embeds (the projection is injective on $F$); $F$ is identified with its images. Let $A_i$ be regular on $\Omega_i$ with $e_r^\beta\le c_0(a)$, and let $A_1=A_2$ on every bond with both ends in $F_0\defeq\{x\in F:\dist(x,\partial F)\ge W_0\}$, for a fixed $W_0\ge0$. Put $H_i\defeq-\Delta_{A_i}+m_r^2+a_{r+1}L^{-2}P_{r+1}(A_i)$ on $\Omega_i$ (the $(r+1)$-pin, $P_{r+1}=Q_{r+1}^*Q_{r+1}$ with the level-$(r+1)$ blocks) and $G_i\defeq H_i^{-1}$. Let $d_i$ be the quotient metric of $\Omega_i$ and $d_\wedge\defeq\min\{d_1,d_2\}$ on $F$; distances to $\partial F$ and to $F_0^c$ are taken in $\Lambda_X$, for a lift of $F$. There are $C,\delta>0$ depending on $a,L,W_0$ only such that for $f,f'$ supported in unit cubes in $F$
\begin{equation}\label{rot:hyb:eq:B1}
|\ip{f'}{(G_1-G_2)f}|\le C\exp\!\big[-\delta\big(d_\wedge(\supp f,\supp f')+\dist(\supp f,\partial F)+\dist(\supp f',\partial F)\big)\big]\|f\|\|f'\|,
\end{equation}
and, for $y,y'\in\Gamma_r\cap F$,
\begin{equation}\label{rot:hyb:eq:B2}
|C^{(r)}(\Omega_1,A_1;y,y')-C^{(r)}(\Omega_2,A_2;y,y')|\le C\exp\!\big[-\delta\big(d_\wedge(y,y')+\dist(y,\partial F)+\dist(y',\partial F)\big)\big].
\end{equation}
The same bound \eqref{rot:hyb:eq:B2} holds for $\Delta^{(r)}=A^X$, from \eqref{rot:hyb:eq:B1} at the $r$-pin. In particular $\delta C^{(r)}$ and $\delta\Delta^{(r)}$ carry the factor $e^{-\delta\dist(\{y,y'\},\partial F)}$ of \cite[Thm.~3.3]{K1}, in the metric $d_\wedge$.
\end{lemma}

The metric is forced: for $y$ near one face of $F$ and $y'$ near the opposite face, $d_1(y,y')=O(1)$, so $C^{(r)}(\Omega_1,A_1;y,y')$ is of order one, while $C^{(r)}(\Omega_2,A_2;y,y')$ is small when $\Omega_2=\Lambda_X$; no bound of the form \eqref{rot:hyb:eq:B2} with the distance $|y-y'|$ of the lift in place of $d_\wedge$ can hold. The kernel bounds on each domain separately are in the quotient distance of that domain (\Cref{rot:hyb:prop:CT}).

\begin{proof}
\emph{Step 0 (decay on each domain).} By \Cref{rot:hyb:prop:CT} at the $(r+1)$-pin, with the distance weight on a torus and the linear weight on $\Lambda_X$, for $f,f'$ supported in unit cubes and $i=1,2$,
$|\ip f{G_if'}|,\ |\ip f{D_{A_i}G_if'}|,\ |\ip f{G_iD_{A_i}^*f'}|\le c_1e^{-\theta_0d_i(\supp f,\supp f')}\|f\|\|f'\|$, $\theta_0=\theta_0(a,L)$.

\emph{Step 1 (locality).} Put $r_1\defeq1+\sqrt3L$ and $S\defeq\{x\in F_0:\dist(x,F_0^c)\ge r_1\}$. For $u$ supported in $S$, $H_1u=H_2u$ (both supported in $F_0$, identified). Indeed $(H_iu)(x)$ involves $u$ at the bond neighbours of $x$ and on the $(r+1)$-block of $x$, the bond variables of the bonds meeting $\supp u$ (both ends in $F_0$), and the block functions $w_z$ of the blocks meeting $S$; such a block lies in $F$ ($F$ is a union of $(r+1)$-blocks) within $\sqrt3L$ of $S$, hence in $F_0$, where $A_1=A_2$; its contours lie inside it (\Cref{rot:hyb:lem:alg}), so the contour phases of $w_z$ agree.

\emph{Step 2 (geometric resolvent identity).} Let $\chi(x)\defeq\min\{1,\max\{0,\dist(x,F_0^c)-r_1\}\}$ (Lipschitz constant $1$), so $\supp\chi\subset S$ and $\chi=1$ on $S''\defeq\{x:\dist(x,F_0^c)\ge r_1+1\}$. Let $f,f'$ be supported in $S''$. Then $\chi G_2f$ is supported in $S$, so by Step~1 $H_1(\chi G_2f)=H_2(\chi G_2f)=\chi f+[H_2,\chi]G_2f=f+[H_2,\chi]G_2f$. The function $[H_2,\chi]G_2f$ vanishes at points farther than $r_1$ from the set where $\chi$ is not locally constant, so it is supported in $F_0$ and is a function on $\Omega_1$. Applying $G_1$ and pairing with $f'$ ($\chi f'=f'$, $\chi$ real),
\begin{equation}\label{rot:hyb:eq:B3}
\ip{f'}{G_2f}-\ip{f'}{G_1f}=\ip{f'}{G_1[H_2,\chi]G_2f}.
\end{equation}

\emph{Step 3 (the commutator).} With $(D_A\phi)(b)=\eta'^{-1}(U(-A_b)\phi(b_+)-\phi(b_-))$, $D_A(\chi u)(b)=\chi(b_-)(D_Au)(b)+(Mu)(b)$, $(Mu)(b)\defeq\eta'^{-1}(\chi(b_+)-\chi(b_-))U(-A_b)u(b_+)$; taking adjoints, $\chi D_A^*-D_A^*\chi_-=M^*$, hence $[-\Delta_A,\chi]=D_A^*M-M^*D_A$. Since $|\chi(b_+)-\chi(b_-)|\le\eta'$ and every site is the end $b_+$ of three bonds, $\|M\|\le\sqrt3$. For the pin, $P_{r+1}=\sum_zw_zw_z^*$ with $\|w_z\|=1$ and $\supp w_z\subset B_{r+1}(z)$ (\Cref{rot:hyb:lem:alg}), so $[P_{r+1},\chi]=\sum_z\big(w_zw_z^*(\chi-\chi(z))-(\chi-\chi(z))w_zw_z^*\big)$, each term acting in $\ell^2(B_{r+1}(z))$ with norm $\le2\sup_{B_{r+1}(z)}|\chi-\chi(z)|\le\sqrt3L$, and zero unless $\chi$ is non-constant on the block; $m_r^2$ commutes with $\chi$. All these terms live on bonds and blocks meeting the slab $\{r_1\le\dist(\cdot,F_0^c)\le r_1+1\}$, hence on the set $Y\defeq\{x:\tfrac12\le\dist(x,F_0^c)\le r_1+2+\sqrt3L\}\subset F_0$ (the pieces lie within $\sqrt3L+\eta'$ of the slab and $r_1-\sqrt3L-\eta'>\tfrac12$), on whose bonds $D_{A_1}=D_{A_2}$.

\emph{Step 4 (bound of \eqref{rot:hyb:eq:B3}).} Write $\ip{f'}{G_1D^*MG_2f}=\ip{DG_1f'}{MG_2f}$, $\ip{f'}{G_1M^*DG_2f}=\ip{MG_1f'}{DG_2f}$ and the pin terms as $\ip{G_1f'}{[w_zw_z^*,\chi]G_2f}$. Split $Y$ into unit cubes $Q$ (every bond and block meets at most $(\sqrt3L+2)^3$ of them) and apply the Cauchy--Schwarz inequality cube by cube, with $\|\mathbf 1_QDG_2f\|=\sup\{|\ip h{DG_2f}|:h\in\ell^2(Q),\|h\|=1\}$ and the same without $D$, and Step~0 on $\Omega_1$ for the left factor and on $\Omega_2$ for the right factor. Since $\dist(x,\partial F)\le\dist(x,F_0^c)+W_0$, $Y\subset K\defeq\{x\in F:\dist(x,\partial F)\le W_0+r_2\}$ with $r_2\defeq2r_1+2$, and
\begin{equation}\label{rot:hyb:eq:B4}
|\ip{f'}{(G_2-G_1)f}|\le C_L\sum_{Q\subset K}e^{-\theta_0(d_1(\supp f',Q)+d_2(Q,\supp f))}\|f\|\|f'\|.
\end{equation}
\emph{Step 5 (geometry).} For $x\in F$ and $q\in F$, a shortest lattice path in $\Omega_i$ from $x$ to $q$, lifted from $x$, either stays in the lift of $F$ (the projection is injective on $F$), and then has length $\ge|x-q|\ge\dist(x,\partial F)-\dist(q,\partial F)$, or it leaves it, and then has length $\ge\dist(x,\partial F)$. Hence $d_i(x,q)\ge\dist(x,\partial F)-\dist(q,\partial F)$, and for $q\in K$, $d_i(x,q)\ge\dist(x,\partial F)-W_0-r_2-1$ (the $1$ for the unit cubes). Splitting $\theta_0(d_1+d_2)=\frac{\theta_0}2d_1+\frac{\theta_0}2(d_1+d_2)$ in \eqref{rot:hyb:eq:B4},
\[
\theta_0\big(d_1(\supp f',Q)+d_2(Q,\supp f)\big)\ge\tfrac{\theta_0}2d_1(\supp f',Q)+\tfrac{\theta_0}2\big(\dist(\supp f,\partial F)+\dist(\supp f',\partial F)\big)-\theta_0(W_0+r_2+3),
\]
and $\sum_Qe^{-(\theta_0/2)d_1(\supp f',Q)}\le C(\theta_0)$, since the unit cubes of a fundamental region at $d_1$-distance in $[n,n+1)$ from a unit cube number at most $C(n+2)^3$ (choose for each the lift nearest to a fixed lift of $\supp f'$). So, for $f,f'$ supported in $S''$,
\[
\text{(ii)}\qquad|\ip{f'}{(G_1-G_2)f}|\le C'e^{-(\theta_0/2)(\dist(\supp f,\partial F)+\dist(\supp f',\partial F))}\|f\|\|f'\|.
\]
Since $\dist(x,F_0^c)\ge\dist(x,\partial F)-W_0$, a unit cube not contained in $S''$ lies within $r_3\defeq W_0+r_1+3$ of $\partial F$. If $\supp f$ is such a cube, then by the first inequality of this step $d_\wedge(\supp f,\supp f')\ge\dist(\supp f',\partial F)-r_3-2$, so $\frac{\theta_0}2(\dist(\supp f,\partial F)+\dist(\supp f',\partial F))\le\frac{\theta_0}2d_\wedge(\supp f,\supp f')+\theta_0(r_3+1)$, and (ii) follows from (i) below with $C'$ replaced by $\max\{C',2c_1e^{\theta_0(r_3+1)}\}$; the same if $\supp f'$ is such a cube. So (ii) holds for all $f,f'$ supported in unit cubes in $F$.

\emph{Step 6 (proof of \eqref{rot:hyb:eq:B1}).} By Step~0, (i) $|\ip{f'}{(G_1-G_2)f}|\le|\ip{f'}{G_1f}|+|\ip{f'}{G_2f}|\le2c_1e^{-\theta_0d_\wedge(\supp f,\supp f')}\|f\|\|f'\|$. Then $\min\{\text{(i)},\text{(ii)}\}\le\text{(i)}^{1/2}\text{(ii)}^{1/2}$, whose exponent is $\frac{\theta_0}2d_\wedge+\frac{\theta_0}4(\dist(\supp f,\partial F)+\dist(\supp f',\partial F))$; this is \eqref{rot:hyb:eq:B1} with $\delta=\theta_0/4$ and $C=(2c_1C')^{1/2}$.

\emph{Step 7 (proof of \eqref{rot:hyb:eq:B2}).} Take $\Lambda=\Omega^{(r)}$ in \cite[(3.5)]{Ba4}; then $B^r(\Lambda)=\Omega$ and $G_r(\Omega_i,\Omega_i^{(r)},A_i)=G_i$. In the conventions above ($Q_rQ_r^*=\mathbf 1$, $QQ^*=\mathbf 1$, $P=Q^*Q$, $P_{r+1}=Q_r^*PQ_r$, with $Q_r=Q_X(A)$ and $Q=Q^{(1)}(A)$), \cite[(3.6)]{Ba4} reads
\[
C^{(r)}(\Omega_i,A_i)=V_i^*G_iV_i-\tfrac{\beta}{a_r}P(A_i)+a_r^{-1},\qquad V_i\defeq Q_r^*(A_i)\big(1-\beta P(A_i)\big),\qquad\beta\defeq\frac{a_{r+1}L^{-2}}{a_r},
\]
with $a_{r+1}=aa_r/(aL^{-2}+a_r)$; it is obtained by integrating the coarse variable out of the Gaussian of \cite[(3.1)]{Ba4}, is an algebraic identity between bounded operators, and holds on $\Lambda_X$ as on a torus. The vector $v_i(y)\defeq V_i\delta_y$ is supported in the $(r+1)$-block of $y$, is built from contours inside that block (\Cref{rot:hyb:lem:alg}), and has $\|v_i(y)\|\le1+\beta\le2$.
\emph{Case 1:} the $(r+1)$-blocks of $y$ and $y'$ lie in $F_0$. Then $v_1(y)=v_2(y)=:v(y)$, the same for $y'$, and $P(A_1;y,y')=P(A_2;y,y')$, since $A_1=A_2$ there; so $\delta C^{(r)}(y,y')=\ip{v(y)}{(G_1-G_2)v(y')}$. Splitting $v(y)$, $v(y')$ into at most $(\sqrt3L+2)^3$ unit-cube pieces and applying \eqref{rot:hyb:eq:B1} to each pair, $|\delta C^{(r)}(y,y')|\le4(\sqrt3L+2)^6e^{6\delta\sqrt3L}C\,e^{-\delta(d_\wedge(y,y')+\dist(y,\partial F)+\dist(y',\partial F))}$.
\emph{Case 2:} a $(r+1)$-block of $y$ or $y'$ meets $F\setminus F_0$, say that of $y$, so $\dist(y,\partial F)\le r_4\defeq W_0+\sqrt3L$. Then $|\delta C^{(r)}|\le|C^{(r)}(\Omega_1)|+|C^{(r)}(\Omega_2)|\le2C_0e^{-\delta_0d_\wedge(y,y')}$ by the single-domain decay of \Cref{rot:hyb:cor:Ck} on each domain (in its quotient distance), and $d_\wedge(y,y')\ge\dist(y',\partial F)-r_4$ by Step~5, so $e^{-\delta_0d_\wedge}\le e^{\delta_0r_4}e^{-(\delta_0/2)(d_\wedge(y,y')+\dist(y,\partial F)+\dist(y',\partial F))}$. Both cases give \eqref{rot:hyb:eq:B2} with $\min\{\delta,\delta_0/2\}$ in place of $\delta$.
For $\Delta^{(r)}=a_r-a_r^2Q_rG^X_rQ_r^*$ (\cite[(2.14)]{K1}, the $r$-pin) the same argument applies with $w_z$ in place of $v(y)$, with \eqref{rot:hyb:eq:B1} at the $r$-pin (Steps 0--6 with $r_1=1+\sqrt3$, \Cref{rot:hyb:lem:coer}).
\end{proof}

\begin{proposition}[{\cite[Prop.~5.6]{K2} for the hybrid flow}]\label{rot:hyb:prop:CLam}
Let $X\defeq A^X(A)+aL^{-2}P(A)$ on the coarse lattice $\Gamma_r$ of the domain, $P=Q^{(1)}(A)^*Q^{(1)}(A)$. Then $(L^2\gamma_0^{-1}+a_r^{-1})^{-1}\le X\le a_r+aL^{-2}$ and $|X(y,y')|\le Ce^{-\theta d(y,y')}$; for every subset $\Lambda$ of the coarse lattice, in particular every union of big blocks, the same bounds hold for the compression $X|_\Lambda$, and $\hat C^{(r)}_\Lambda\defeq(X|_\Lambda)^{-1}$ (\cite[(5.53)]{K2}) obeys $|\hat C^{(r)}_\Lambda(y,y')|\le Ce^{-t_*d(y,y')}$ (\cite[(5.54)]{K2}), $|\hat C^{(r)}_\Lambda(y,y')-C^{(r)}(y,y')|\le Ce^{-(t_*/6)(d(y,y')+d(y,\Lambda^c)+d(y',\Lambda^c))}$, and, under the hypotheses of \Cref{rot:hyb:lem:domdiff} with $\Lambda\subset F$, the two-domain bound \eqref{rot:hyb:eq:B2} in the metric $d_\wedge$. The constants depend on $a$, $L$ (and $W_0$) only.
\end{proposition}
\begin{proof}
Upper bound: $A^X=(a_r^{-1}+Q_rC_AQ_r^*)^{-1}\le a_r$ and $P\le1$. Lower bound: \cite[(3.6)]{Ba4} in the form of Step~7 above gives $X^{-1}=C^{(r)}=V^*G_{r+1}V-(\beta/a_r)P+a_r^{-1}\le\|V\|^2\|G_{r+1}\|+a_r^{-1}$ with $\|V\|\le1$ and $\|G_{r+1}\|\le L^2\gamma_0^{-1}$ (\Cref{rot:hyb:cor:L2}, the $(r+1)$-pin). Decay of $X$: \Cref{rot:hyb:prop:CT} for $A^X$ and the finite range of $P$ (its kernel vanishes unless $y,y'$ lie in one $L$-block, of diameter $<\sqrt3L$). Compression to $\Lambda$ keeps both bounds and the entries. Then, for a $1$-Lipschitz $\psi$ on $\Lambda$, $e^{t\psi}X|_\Lambda e^{-t\psi}-X|_\Lambda$ has entries bounded by $Ct\,d(y,y')e^{-(\theta/2)d(y,y')}$ for $t\le\theta/2$, hence norm $\le Ct$ (Schur test), so $e^{t\psi}X|_\Lambda e^{-t\psi}$ is invertible with inverse of norm $\le2(L^2\gamma_0^{-1}+a_r^{-1})$ for $t\le t_*$, $t_*$ depending on $a,L$ only; with $\psi=\min\{d(\cdot,y'),d(y,y')\}$ this gives $|\hat C^{(r)}_\Lambda(y,y')|\le Ce^{-t_*d(y,y')}$ (the entry is bounded by the operator norm on the counting measure of $\Gamma_r$). The difference with $C^{(r)}$ is $\hat C^{(r)}_\Lambda X_{\Lambda\Lambda^c}C^{(r)}_{\Lambda^c\Lambda}$ (the block identity for the inverse of a compression), whose decay follows from the three factors' decay and the geometric mean as in Step~6; the two-domain bound is the resolvent identity on $\ell^2(\Lambda)$ with \eqref{rot:hyb:eq:B2} for $\Delta^{(r)}$ ($\delta P=0$ off the collar), Step~5 and Step~6's interpolation.
\end{proof}

% r03_hyb_A0: proof of \Cref{rot:prop:BD}, part IV: the A=0 input (the hybrid kernels on Lambda_X and on periodic cells) and the pointwise and Hoelder bounds at a regular background on a cell.

\paragraph{The $A=0$ input and the bounds on a cell.}

\begin{lemma}[$A=0$ bounds for the hybrid flow]\label{rot:hyb:lem:A0}
On $\Lambda_X$, with $S^X_j$ the slice kernels of \Cref{rot:def:match} ($G^X_r(0)=\sum_{j<r}S^X_j$ exactly):
\begin{enumerate}[label=\textup{(\alph*)},leftmargin=*]
\item $|S^X_j(x,y)|\le C\ell_j^{-1}e^{-\delta|x-y|/\ell_j}$, $|\nabla_xS^X_j(x,y)|\le C\ell_j^{-2}e^{-\delta|x-y|/\ell_j}$, and for $0<\alpha<1$
\[
|\nabla_xS^X_j(x,y)-\nabla_xS^X_j(x',y)|\le C_\alpha\log^{3\alpha}(2+\ell_j/\eta')\,\ell_j^{-2-\alpha}|x-x'|^\alpha e^{-\delta\min(|x-y|,|x'-y|)/\ell_j}
\]
(the logarithm is the slice's own, $\le\log^{3\alpha}(1/\eta')$);
\item consequently $\|G^X_r(0)f\|_\infty+\|\nabla G^X_r(0)f\|_\infty\le C\|f\|_\infty$ with $C$ uniform in $r$, $[\nabla G^X_r(0)f]_\alpha\le C_\alpha\log^{3\alpha}(1/\eta')\|f\|_\infty$, and $\|G^X_r(0)f\|_q\le C\|f\|_p$ for $0\le\frac1p-\frac1q<\frac23$, with the exponential localisation $e^{-\delta_0\dist(x,\supp f)}$ (\cite[(2.16)--(2.17)]{Ba4});
\item on a periodic cell $\square=\Lambda_X/P$ fitted by both lattices, of side $\ge2M$, and for a constant field $B_0$,
\[
G^X_r(\square,B_0)(x,y)=e^{\mathrm ie_rB_0\cdot(x-y)}\sum_{\gamma\in P}e^{\mathrm ie_rB_0\cdot\gamma}G^X_r(0)(x,y+\gamma)
\]
(twisted periodisation), and (b) holds on $\square$ for $G^X_r(\square,B_0)$ with constants independent of $B_0$ and $M$, the H\"older constant carrying the same $\log^{3\alpha}(1/\eta')$.
\end{enumerate}
\end{lemma}
\begin{proof}
(a) In the unit scale of level $j$, $S^X_j=H^X_jC^X_jH^{X*}_j$ (\Cref{rot:def:match}). The decay $|H^X_j(x,z)|\le Ce^{-\delta|x-z|}$ and the gradient bound $|\nabla H^X_j|\le C_\nabla$ are \Cref{rot:lem:H,rot:lem:grad} (the decay is the Combes--Thomas decay of the coercive coarse operators, \cite{CT73}); the decay of the coarse kernel $C^X_j=(A^X_j+bQ^*Q)^{-1}$ is the one used in the proof of \Cref{rot:lem:slice} (its symbol is analytic and bounded on a strip because $A^X_j+bQ^*Q\ge c$). \emph{The H\"older derivative.} The second difference of $H^X$ in $x$ multiplies the alias-$l$ summand of $M^X$ in the representation of \Cref{rot:lem:H} by $d(q+l)^2$ with $|d(Q)|\le|Q|$; since $\Delta'_X(Q)\ge c|Q|^2$ on each alias cell and $\sum_lu_X(q+l)^2=1$, the alias sum is $\le C\sum_l|u_X(q+l)|\le C$, and integrating against $|u_{r-1}|$ over $V_\eta$ costs $J_0\le C\log^3(1/\eta)$ (\Cref{rot:lem:CK0}). So $|\nabla\nabla H^X|\le\Lambda_\eta\defeq C\log^3(1/\eta)$ and, by interpolation with $|\nabla H^X|\le C_\nabla$,
\[
|\nabla H^X(x,z)-\nabla H^X(x',z)|\le\min\{2C_\nabla,\Lambda_\eta|x-x'|\}\le(2C_\nabla)^{1-\alpha}\Lambda_\eta^{\alpha}|x-x'|^\alpha\le C_\alpha\log^{3\alpha}(1/\eta)|x-x'|^\alpha
\]
for every $\alpha<1$; the geometric mean with the decay gives the localised form. Composing with $C^X_j$ and $H^{X*}_j$ (exponentially decaying kernels) and rescaling by \cite[(2.20)]{K1} (scalar kernels scale as $\ell^{-1}$ in $d=3$, derivatives as $\ell^{-2}$) gives (a); at level $j$ the logarithm is $\log(1/\eta_j)$ with $\eta_j=5^{1-j}$, i.e.\ $\le\log(2+\ell_j/\eta')$. \emph{Summation.} $\|S^X_jf\|_\infty\le C\ell_j^{-1}\ell_j^3\|f\|_\infty$, $\|\nabla S^X_jf\|_\infty\le C\ell_j\|f\|_\infty$, $[\nabla S^X_jf]_\alpha\le C_\alpha\log^{3\alpha}(2+\ell_j/\eta')\,\ell_j^{1-\alpha}\|f\|_\infty$; summing over $j<r$ converges for $\alpha<1$, the H\"older term to $\le C_\alpha\sum_{j<r}5^{-(r-j)(1-\alpha)}(1+j\ln5)^{3\alpha}\le C_\alpha\log^{3\alpha}(1/\eta')$ (this is \cite[(2.38)--(2.39)]{Ba4} with the slice kernels in place of $G_jQ_j^*C^{(j)}Q_jG_j$). The $L^p\to L^q$ bounds follow from the kernel bound $|G^X_r(0)(x,y)|\le C\min(\eta'^{-1},|x-y|^{-1})e^{-\delta|x-y|}$ by Young's inequality and the Riesz--Thorin theorem (\cite[p.~583]{Ba4}). (c) Both sides are the inverse of $-\Delta_{B_0}+m_r^2+a_rP_X(B_0)$ on $P$-periodic functions: $P\subset\Gamma_r$ makes the hybrid block system $P$-periodic, and $e^{\mathrm ie_rB_0\cdot x}$ conjugates the operator to the $A=0$ operator on functions with the twisted periodicity $\psi(x+\gamma)=e^{-\mathrm ie_rB_0\cdot\gamma}\psi(x)$, whose kernel is the twisted image sum (\cite[(B.3)]{K2} with phases). The images with $\gamma\ne0$ contribute $\le Ce^{-\delta_0(2M-\sqrt3)}$ to every bound in (b), uniformly in $B_0$.
\end{proof}

\begin{lemma}[{\cite[Lemma~2.2]{Ba4}, \cite[Lemma~5.8]{K2}, for the hybrid flow}]\label{rot:hyb:lem:HOL}
Let $\square$ be a periodic cell fitted by both lattices of side $2M$, and $\tilde A=A_0+A'$ regular on $\square$, $A_0$ constant, with $|A'|_\infty,|\partial A'|_\infty\le c'e_r^{\beta-1}$ and $A'$ supported away from the boundary of the fundamental region. Then for $e_r^\beta\le c_0(a,M)$, with $c_0,c_1,c_2$ uniform in $r$,
\[
\|G^X_r(\square,\tilde A)f\|_{1,0}\le c_1\|f\|_\infty,\qquad\|G^X_r(\square,\tilde A)f\|_{1,\alpha}\le c_1C_\alpha\log^{3\alpha}(1/\eta')\|f\|_\infty,
\]
\[
\|G^X_r(\square,\tilde A)f\|_q,\ \|D_{\tilde A}G^X_rf\|_q,\ \|G^X_rD_{\tilde A}^*f\|_q\le c_2\|f\|_p
\]
for $1\le p,q\le\infty$ with $\frac1p-\frac1q<\frac23$, where $\|\cdot\|_{1,\alpha}$ is the covariant H\"older norm \cite[(2.14)]{Ba4} with the field $\tilde A$ and $\|\cdot\|_{1,0}$ the sup norm plus the sup of the covariant derivative.
\end{lemma}
\begin{proof}
\cite[(2.24)--(2.26), (2.31)--(2.33)]{Ba4}: $G_r(\square,\tilde A)=\sum_{n\ge0}G_r(\square,A_0)[V_rG_r(\square,A_0)]^n$ with $V_r=V_r(A',A_0)$ the vertex operator of \cite[(3.16)]{Ba1} ($=$ \cite[(3.53)]{K1}); $V_r$ is a first-order differential operator with coefficients $O(e_r|A'|,e_r|\partial A'|)$ (\cite[(2.32)]{Ba4}, using that $A'$ vanishes near the boundary of the fundamental region) plus the pin part $a_r[Q(A_0)^*F_{2,r}(A',A_0)+F_{2,r}^*Q(A_0)+F_{2,r}^*F_{2,r}]$ with $(F_{2,r}(A',A_0)\phi)(y)=\sum_{x\in B(y)}\eta'^3(U(A'(\Gamma_{y,x}))-1)U(A_0(\Gamma_{y,x}))^{-1}\phi(x)$, $\|F_{2,r}\|_{\infty\to\infty}\le\sup|U(A'(\Gamma))-1|\le2.1\,e_r|A'|_\infty$ by \Cref{rot:hyb:lem:len}(i). Hence $\|V_rg\|_\infty\le C(M)e_r^\beta\|g\|_{1,0}$. By \Cref{rot:hyb:lem:A0}(b),(c), $\|G_r(\square,A_0)g\|_{1,0}\le c\|g\|_\infty$ with $c$ uniform in $r$ and $A_0$, and $\|G_r(\square,A_0)g\|_{1,\alpha}\le cC_\alpha\log^{3\alpha}(1/\eta')\|g\|_\infty$ (the constant background is a twisted periodisation of the $A=0$ kernel, and the H\"older norm (2.14) with field $A_0$ is the plain H\"older norm after the local gauge $e^{\mathrm ie_rA_0\cdot x}$ on each unit cube). \emph{Order of summation.} Put $g_0\defeq f$, $g_n\defeq V_rG_r(\square,A_0)g_{n-1}$. Then $\|g_n\|_\infty\le C(M)e_r^\beta\|G_r(\square,A_0)g_{n-1}\|_{1,0}\le(C(M)c\,e_r^\beta)^n\|f\|_\infty$, so for $C(M)c\,e_r^\beta\le\frac12$, a condition uniform in $r$, the sum $g\defeq\sum_ng_n$ converges in the sup norm with $\|g\|_\infty\le2\|f\|_\infty$, and $G_r(\square,\tilde A)f=\sum_nG_r(\square,A_0)g_n=G_r(\square,A_0)g$ (the cell is finite, every operator is bounded, and the sums may be exchanged). Hence $\|G_r(\square,\tilde A)f\|_{1,0}\le2c\|f\|_\infty$ and $\|G_r(\square,\tilde A)f\|_{1,\alpha}\le2cC_\alpha\log^{3\alpha}(1/\eta')\|f\|_\infty$: the H\"older constant enters once, at the end, and never the smallness condition. (Summing in $\|\cdot\|_{1,\alpha}$, as \cite[(2.33)]{Ba4} may with its uniform constant, would turn the condition into $e_r^\beta\log^{3\alpha}(1/\eta')\le c_0$, which depends on $r$.) The norms (2.14) with the fields $A_0$ and $\tilde A$ differ on the cell by at most $Ce_r^\beta\|\cdot\|_{1,0}$ ($D_{\tilde A}=U(A')D_{A_0}+F_{1,r}(A')$ and $|U(A'(\Gamma_{y,x}))-1|\le Ce_r|A'|_\infty|x-y|$ for the bounded-length paths of (2.14)), which changes $c_1$ by a factor $1+O(e_r^\beta)$. The $L^p\to L^q$ bounds are summed in the same order, with $D_{\tilde A}=U(A')D_{A_0}+F_{1,r}(A')$.
\end{proof}

\begin{remark}[the H\"older constant]\label{rot:hyb:rem:Hol}
In \Cref{rot:prop:BD}(a) the H\"older clause \cite[(3.8)]{K1} of $G^X_r(A)$, and the H\"older clauses of $H^X(A)$ and $\Gamma^X(A)$, hold with the constant $C_\alpha\log^{3\alpha}(1/\eta')$, $0<\alpha<1$ (\Cref{rot:hyb:lem:A0}, \Cref{rot:hyb:lem:HOL}, and the random walk below); on slice $j$ the constant is $C_\alpha\log^{3\alpha}(2+\ell_j/\eta')$ (\Cref{rot:hyb:lem:A0}(a)). The sup, decay and one-derivative bounds (3.6)--(3.7) and the $\|\cdot\|_{1,0}$ bounds have constants uniform in $r$, and so does every smallness condition. \cite[(3.8)]{K1} asks for some $0<\alpha<1$ with some constant (\cite[p.~656]{K1}). Every use of the clause in this section is either as a gain, where $\eta'^\alpha\log^{3\alpha}(1/\eta')\le C_{\alpha,\alpha'}\eta'^{\alpha'}$ and, slice by slice, $(\eta'/\ell_j)^\alpha\log^{3\alpha}(2+\ell_j/\eta')\le C_{\alpha,\alpha'}(\eta'/\ell_j)^{\alpha'}$ for $\alpha'<\alpha$ (\Cref{rot:thm:H1S}, whose exponent $\theta_1$ is fixed by $\alpha$ alone), or inside a per-flow bound that already allows a factor $(1+k)^p$ (\Cref{rot:cor:OF}), where $\log^{3\alpha}(1/\eta')=(r\ln5)^{3\alpha}\le C(1+k)^{3\alpha}$.
\end{remark}

% r03_hyb_walk: proof of \Cref{rot:prop:BD}, part V: the random walk with periodic cells, the boundary clause, the remaining kernels, the slice clauses, the levels 0 and 1; this completes part (a).

\begin{proof}[Proof of \Cref{rot:prop:BD}(a)]
\emph{Random walk.} Let $\Omega$ be $\Lambda_X$ or a torus fitted by both lattices with $P\subset M\Z^3$, $A$ regular on $\Omega$. For $j\in\Z^3$ let $\square_j$ be the periodic cell of side $2M$ ``centred'' at $Mj$ and fitted by both lattices (which is possible iff $25\mid2M$, as the tori of the chain allow; \cite[p.~345]{K2}: ``blocks $\{\square_j\}$ which are small scale versions of $D_J$, that is they can be fitted by both the usual lattice and the rotated lattice''), $h_j$ the partition of unity of \cite[\S2]{Ba4} ($h_j=1$ on the $M$-cube around $Mj$, supported in the $\frac43M$-cube, $\sum_jh_j^2=1$), and $\tilde A_j=A_0^{(j)}+\theta_jA'^{(j)}$ with $A_0^{(j)}=A(Mj)$ and the cutoff $\theta_j$ of \cite[p.~575]{Ba4}, so that $\tilde A_j=A$ on $\supp h_j$, is regular with constant $c'=O(Mc)$, and is constant near the boundary of the fundamental region of $\square_j$. With $G_0\defeq\sum_jh_jG^X_r(\square_j,\tilde A_j)h_j$ one computes, exactly as \cite[(2.3)--(2.9)]{Ba4} and \cite[(5.8)--(5.12)]{K2},
\[
(-\Delta_A+m_r^2+a_rP_X(A))G_0=\mathbf 1-\cR,\qquad \cR=\sum_jK_jG^X_r(\square_j,\tilde A_j)h_j,
\]
\[
(K_j\phi)(x)=\mathbf 1_{\square_j}(x)\Big[\sum_{b\in\mathrm{st}(x)}(\partial h_j)(b)(D_A\phi)(b)+(\Delta h_j)(x)\phi(x)-(R_r(A,\partial h_j)\phi)(x)\Big],
\]
where $R_r(A,\partial h)\phi(x)=\sum_{x'\in B(z(x))}\eta'^3U(A(\Gamma_{z,x}))U(A(\Gamma_{z,x'}))^{-1}(\partial h)(\Gamma_{x,x'})\phi(x')$ collects the commutator of the pin with the multiplication by $h_j$ (\cite[(2.8)]{Ba4}); $(\partial h)(\Gamma_{x,x'})=h(x')-h(x)$ for any path, so $|R_r(A,\partial h_j)\phi(x)|\le\sup|\partial h_j|\cdot2.1\cdot\sup_B|\phi|$ by \Cref{rot:hyb:lem:len}(i). The three identities used --- the product rule for $-\Delta_A$ (\cite[(2.5)]{Ba4}), the commutator of the pin with a smooth function (\cite[(2.8)]{Ba4}), and the agreement of the cell operator with the global one on $\supp h_j$ (\cite[(2.6)]{Ba4}; for a periodic cell this holds because $\supp h_j$ stays at distance $\ge\sqrt3$ from the boundary of the fundamental region, so no block and no bond of $\supp h_j$ wraps) --- hold for any block system with contours from the root to the point. Hence $G^X_r(\Omega,A)=\sum_{n\ge0}G_0\cR^n$ in $L^2$ (\Cref{rot:hyb:cor:L2} gives $\|\cR\|_{2\to2}\le O(M^{-1})$ as in \cite{Ba4}, since $\sup|\partial h_j|=O(M^{-1})$ and $\sup|\Delta h_j|=O(M^{-2})$), i.e.\ the path representation \cite[(2.13)]{Ba4}, \cite[(5.56)]{K2}.

\emph{Bounds.} \cite[(2.18)--(2.22)]{Ba4} verbatim: for $x,x'$ in $\supp h_{\omega_0}$ and $f$ supported in $\square_{\omega_n}$, each factor $\|K_{\omega_i}G^X_r(\square_{\omega_i},\tilde A_{\omega_i})h_{\omega_i}\|_{\infty\to\infty}$ is $\le O(M^{-1})c_1$ by \Cref{rot:hyb:lem:HOL} and the bound on $K_j$ ($c_1$ the uniform $\|\cdot\|_{1,0}$ constant of \Cref{rot:hyb:lem:HOL}: $K_j$ takes one covariant derivative of its argument, so the choice of $M$ below is uniform in $r$), the first factor is bounded in $\|\cdot\|_{1,\alpha}$ by \Cref{rot:hyb:lem:HOL} (the only place where $C_\alpha\log^{3\alpha}(1/\eta')$ enters, and only in the clause (3.8)), the number of paths of length $n$ is $\le2^d(3^d)^{n-1}2^d$, and $n\ge M^{-1}\dist(\{x,x'\},\supp f)-2$; for $M$ fixed with $3^dc_1O(1)M^{-1}\le e^{-1}$ the sum converges and gives \cite[(3.7)--(3.8)]{K1} with $\delta_0=(2M)^{-1}$. The tail with $n>n_0$ is handled with the $L^p$ bounds (2.17) as in \cite[(2.21)]{Ba4}. The bounds (3.6) for $\Delta^{(r)}$ and $C^{(r)}$ are \Cref{rot:hyb:prop:CT} and \Cref{rot:hyb:cor:Ck}; the bounds for the next-step precision $\Delta^{(r)}+aL^{-2}\hat Q^*\hat Q$ and its restricted inverses $\hat C^{(r)}_\Lambda$ are \Cref{rot:hyb:prop:CLam}.

\emph{Two-torus clause.} For two domains $\Omega_1$, $\Omega_2$ sharing a fundamental region $F$ as in \Cref{rot:hyb:lem:domdiff} (or a torus against $\Lambda_X$), the representations coincide term by term for the paths whose cells do not meet the boundary of $F$, and the difference is bounded, as on \cite[p.~579]{Ba4}, by the same sum restricted to paths meeting the boundary; such a path of $n$ cells also joins $x$ to $\supp f$ in its own domain, so $\sqrt3Mn\ge\max\{d_i(\{x,x'\},\supp f),\dist(\{x,x'\},\partial F)+\dist(\supp f,\partial F)\}-c$ for $i=1,2$, and the average of the two exponents gives the factor $e^{-(2\sqrt3M)^{-1}(d_\wedge(\{x,x'\},\supp f)+\dist(\{x,x'\},\partial F)+\dist(\supp f,\partial F))}$ for the difference of the $G^X_r$ clauses (3.7)--(3.8), with $d_\wedge\le d_i$ the smaller of the two quotient distances. For $\Delta^{(r)}$, $C^{(r)}$ and $\hat C^{(r)}_\Lambda$ the two-torus clause is \Cref{rot:hyb:lem:domdiff} and \Cref{rot:hyb:prop:CLam}. The metric $d_\wedge$ cannot be replaced by the distance of the lift (the remark after \Cref{rot:hyb:lem:domdiff}); the uses of the clause in this section (\Cref{rot:lem:CB}(N), the kernel sweep of \Cref{rot:thm:Asrc}, and \Cref{rot:lem:TC}(IV)) need only the decay on each torus in its own quotient distance together with the boundary factor, which is what is proved.

\emph{The remaining kernels.} $H^X(A)=a_rG^X_r(A)Q_X(A)^*$: apply (3.7)--(3.8) to $f=w_z$ ($\|w_z\|_\infty=1$, support in a block of diameter $\le\sqrt3$). $B^X(A)=a_r^{-1}+Q_X(A)C_AQ_X(A)^*$: by the diamagnetic inequality $|C_A(x,y)|\le C_0(x,y)=(-\Delta+m_r^2)^{-1}(x,y)$ entrywise ($-\Delta_A+m_r^2=c-T_A$ with $|T_A|=T_0$ entrywise and $(c-T_0)^{-1}=\sum c^{-n-1}T_0^n\ge0$), so $|B^X(z,w)|\le a_r^{-1}\delta_{zw}+\sum_{x\in B(z),y\in B(w)}\eta'^6C_0(x-y)\le C(1+|z-w|)^{-1}$ ($d=3$). $\Gamma^X(A)=H^XB^XH^{X*}$: the fine-variable regularity is that of $H^X$, and $|\Gamma^X(x,y)|\le\sum_{z,w}Ce^{-\delta_0|x-z|}(1+|z-w|)^{-1}e^{-\delta_0|y-w|}\le C(1+|x-y|)^{-1}$. $B^X(A)$ and $\Gamma^X(A)$ are not unit-decaying: $\Gamma^X(x,y)\approx C_A(x,y)$ for $|x-y|\gg1$, at the coarse mass $m_r=m\delta_r\to0$; \cite[Thm.~3.3]{K1} never claims otherwise (it concerns $G_r$, $\Delta^{(r)}$, $C^{(r)}$). Where a renormalised family carries a line $\Gamma^F$ or $(A^F)^{-1}$ at the top scale, either the coarse-slice decomposition of \Cref{rot:cor:unif}(b) is used, or, in the localised form \cite[(3.55)]{K1}, all vertex positions lie in the region $F_v$ of \Cref{rot:lem:CB}, of diameter $\le Cp_r$, where a line bounded by $C(1+|x-y|)^{-1}$ costs at most a factor $Cp_r^3$, absorbed in the field-norm polynomial of \Cref{rot:prop:cp}.

\emph{The slice clauses.} Nothing in the argument above uses the level: the same proof at level $j\le r$, in the unit scale of level $j$, gives (3.6)--(3.8) for the step-$j$ operators $G_j$, $\Delta^{(j)}$, $C^{(j)}$, $H_j$ of the hybrid flow with constants uniform in $j$ (the first step, with the tiles as blocks, is covered by \Cref{rot:hyb:lem:coer} and \Cref{rot:hyb:lem:A0} with $B^X_1$ in place of $B^X_r$; the steps $j\ge2$ are the standard rotated steps acting on composite blocks, \Cref{rot:hyb:lem:alg}(iv)). King's derivation of the slice bounds \cite[(3.63)--(3.65)]{K1} from Theorem~3.3 (\cite[Prop.~3.7]{K1}: the slice $S_j=H_jC_jH_j^*$ at a regular background, the one- and two-ended contour clauses summed along contours of bounded length) therefore applies verbatim to the hybrid flow, with the covariant lattice derivatives and transported H\"older quotients of \cite[(2.10)--(2.11)]{Ba3} at regular $A$; after the constant-background reduction of \Cref{rot:lem:CB} the compared lines are the $A=0$ slices of \Cref{rot:hyb:lem:A0}, so there the plain form holds. For the contour clauses the contour lengths are those of \Cref{rot:hyb:lem:len}(i).

\emph{Levels $0$ and $1$.} In the form in which the hybrid lower bound (\Cref{rot:prop:H4}) consumes (a) at its first steps (precision bounds, precision decay, sup-decay of $C^{(0)}$ and $C^{(1)}$), (a) holds also at $j\in\{0,1\}$ by the same proofs, which use $r\ge2$ nowhere: at $j=1$ the level-$1$ block is the tile, inside the ball of radius $\frac{\sqrt3}2$ about its centre, with an in-tile contour of length $\le1$ in the unit scale of level $1$, so \Cref{rot:hyb:lem:coer}, \Cref{rot:hyb:cor:L2}, \Cref{rot:hyb:prop:CT} and \Cref{rot:hyb:lem:domdiff} read verbatim (the first Neumann eigenvalue of the cube is then $100\sin^2(\pi/40)>\frac14$), and the fine spacing $\eta'=\frac15$ is a constant; at $j=0$ the precision $-\Delta_A+m_0^2+aL^{-2}\hat Q^*\hat Q$ has finite range and its lower bound is the $(r+1)$-pin clause of \Cref{rot:hyb:cor:L2} at $r=0$.

\emph{The Neumann variant.} \cite[Thm.~3.3]{K1} is printed for Neumann boundary conditions on a rectangular parallelepiped of large blocks. For the hybrid flow such a domain is a digitised rotated parallelepiped, on which neither the reflection representation \cite[(2.42)]{Ba4} nor the axis-cube cover near the boundary is available; the Neumann variant is not claimed and not used: \cite[\S\S4--5]{K2} decouple with periodic cells in the random walk, the localisation of \Cref{rot:lem:CB} runs on the infinite lattice in the universal-cover form, and the only Neumann object of the chain, the boundary factor of $\delta G_r$, is replaced by the two-torus clause above and the same-domain clause \Cref{rot:lem:CB}(N).

\emph{The standard flow.} For the flow $Y$ every statement of (a) is \cite[Thm.~3.3]{K1}, proved in \cite{Ba4} on $\eta'\Z^3$ in the frame of $\Yk$, with King's uniform constants; its two-torus clause and restricted inverses are \cite[Prop.~5.6]{K2} and \Cref{rot:hyb:lem:domdiff}, whose proof uses nothing specific to the hybrid blocks. This completes the proof of (a).
\end{proof}

% r03_hyb_ES: proof of \Cref{rot:prop:BD}, part VI: the continuum extension of the Y kernels (cubic-Hermite cardinal interpolation), its decay, its multiplier form at the X sites, and the decaying forms at the X sites.

\paragraph{The continuum extension of the $Y$ kernels.}
The comparisons of \Cref{rot:hyb:lem:CV,rot:hyb:lem:CLN} below compare a sum over $X$ bonds with a sum over $Y$ bonds through \Cref{rot:hyb:lem:CL}, which needs a $C^1$ field on the block that equals the $Y$ kernel on the $Y$ lattice. The Fourier formula of \Cref{rot:lem:H} evaluated off the lattice (the sharp-zone extension) is not admissible for this: the symbol does not vanish on the faces of the zone, so the extension and its gradient have algebraic tails of order $\eta'^3/|r-z|$ and $\eta'^2/|r-z|$ off the lattice, not summable over a torus in $d=3$. The field used is the cubic-Hermite cardinal interpolation of the lattice kernels: it is $C^1$, equals the lattice kernel on $\Lambda_Y$, has exponentially decaying value and gradient by the lattice bounds alone, and obeys the multiplier comparison of \Cref{rot:lem:H} at the $X$ sites with an explicit constant.

\begin{definition}[cubic-Hermite cardinal extension]\label{rot:hyb:def:ES}
$S(u)\defeq3u^2-2u^3$, $\psi(s)\defeq S(1-|s|)\mathbf 1_{|s|<1}=(1-|s|)^2(1+2|s|)\mathbf 1_{|s|<1}$, $\Psi(v)\defeq\prod_\alpha\psi(v\cdot f_\alpha/\eta')$ (the frame of $\Yk$). For $F$ on $\Lambda_Y=R^{-1}(\eta'\Z^3)$ put $(I_SF)(r)\defeq\sum_{y\in\Lambda_Y}\Psi(r-y)F(y)$, $r\in\R^3$. The extensions are $E^S(r,z)\defeq I_S[H^Y(\cdot,z)](r)$ for the level-$r$ lift (scalar and photon), $E^S_j(r,z)\defeq I_S[H^Y_j(\cdot,z)](r)$ for the level-$j$ lift, $\tilde S^S_j(r,y)\defeq\sum_{z,z'}E^S_j(r,z)C^Y_j(z,z')\overline{H^Y_j(y,z')}=I_S[S^Y_j(\cdot,y)](r)$ for $y\in\Lambda_Y$ ($S^Y_j=H^Y_jC^Y_jH^{Y*}_j$, \Cref{rot:def:match}), likewise for the photon slices, and $E^S\varphi\defeq I_S[H^Y\varphi]$ for a coarse field $\varphi$ (componentwise in any frame: $I_S$ is scalar and linear). The extension acts in the first argument only; second arguments stay on $\Lambda_Y$.
\end{definition}

\begin{lemma}[cell identity]\label{rot:hyb:lem:ESH}
Write the frame coordinates of $r$ as $f_\alpha\cdot r=\eta'(n_\alpha+u_\alpha)$, $n\in\Z^3$, $u\in[0,1)^3$, with cell corners $y_\sigma\defeq R^{-1}\eta'(n+\sigma)$, $\sigma\in\{0,1\}^3$, and weights $\omega_\sigma(u)\defeq\prod_\alpha\omega_{\sigma_\alpha}(u_\alpha)$, $\omega_0\defeq1-S$, $\omega_1\defeq S$ (so $\omega_\sigma\ge0$, $\sum_\sigma\omega_\sigma=1$). Then
\textup{(i)} $I_SF(r)=\sum_\sigma\omega_\sigma(u)F(y_\sigma)$;
\textup{(ii)} $\partial_{f_1}I_SF(r)=S'(u_1)\sum_{\sigma'\in\{0,1\}^2}\omega_{\sigma'}(u_2,u_3)\,(\partial_1F)(y_{(0,\sigma')})$ with the lattice forward difference $\partial_1F(y)\defeq[F(y+\eta'f_1)-F(y)]/\eta'$, and likewise for $f_2,f_3$;
\textup{(iii)} $|I_SF(r)|\le\max_\sigma|F(y_\sigma)|$ and $|\partial_{f_\alpha}I_SF(r)|\le\frac32\max|\partial_\alpha F|$ over the four cell edges parallel to $f_\alpha$, so $|\nabla I_SF(r)|\le\frac{3\sqrt3}2\max_{\mathrm{edges}}|\partial F|$;
\textup{(iv)} $I_SF\in C^1(\R^3)$ and $I_SF=F$ on $\Lambda_Y$.
\end{lemma}
\begin{proof}
Only the corners with $|(r-y)\cdot f_\alpha|<\eta'$ for all $\alpha$ contribute, i.e.\ the $y_\sigma$; for them $\psi((r-y_\sigma)\cdot f_\alpha/\eta')=\psi(u_\alpha-\sigma_\alpha)$, which is $S(1-u_\alpha)=1-S(u_\alpha)$ for $\sigma_\alpha=0$ and $S(u_\alpha)$ for $\sigma_\alpha=1$ ($S(u)+S(1-u)=1$): this is (i). Differentiating (i) in $r\cdot f_1=\eta'(n_1+u_1)$ gives $\eta'^{-1}S'(u_1)\sum_{\sigma'}\omega_{\sigma'}[F(y_{(1,\sigma')})-F(y_{(0,\sigma')})]$, which is (ii). (iii) follows from $\omega\ge0$, $\sum\omega=1$, $0\le S'(u)=6u(1-u)\le\frac32$. (iv) $\psi'(s)=-6s(1-|s|)$ on $|s|<1$ vanishes at $s=0,\pm1$, so $\psi\in C^1(\R)$, and the sum is locally finite; $\psi(m)=\delta_{m0}$ for $m\in\Z$.
\end{proof}

\begin{lemma}[decay of the extensions]\label{rot:hyb:lem:ESB}
The lattice inputs are \cite[(3.6), (3.7), (3.63)]{K1} for the standard flow: $|H^Y(y,z)|,|\partial H^Y(y,z)|\le Ce^{-\delta_0|y-z|}$ ((3.7) with $f=Q^*\delta_z$, supported in the block of $z$, $\|f\|_\infty\le1$; for the photon, ``the same bounds for the vector field operators'', \cite[p.~656]{K1}); at level $j$, in the unit scale of level $j$, $|H^Y_j|,|\partial H^Y_j|\le Ce^{-\delta_0|\cdot|}$ and $|C^Y_j(z,z')|\le Ce^{-\delta_0|z-z'|}$. Then for $0<\delta<\delta_0$, uniformly in $j,r,\eta'$, the torus and $m_j\ge0$, for all $r\in\R^3$, $z$, and $y\in\Lambda_Y$, in the unit scale of level $r$:
\begin{enumerate}[label=\textup{(\alph*)},leftmargin=*]
\item $|E^S(r,z)|,\ \frac{2}{3\sqrt3}|\nabla_rE^S(r,z)|\le Ce^{\sqrt3\delta_0\eta'}e^{-\delta_0|r-z|}$, and the same for $E^S_j$ in the unit scale of level $j$;
\item $|\tilde S^S_j(r,y)|\le C\ell_j^{-1}e^{-\delta|r-y|/\ell_j}$, $|\nabla_r\tilde S^S_j(r,y)|,\ |\partial_y\tilde S^S_j(r,y)|\le C\ell_j^{-2}e^{-\delta|r-y|/\ell_j}$ and $|\nabla_r\partial_y\tilde S^S_j(r,y)|\le C\ell_j^{-3}e^{-\delta|r-y|/\ell_j}$ ($\partial_y$ a lattice difference in the second argument);
\item $E^S(\cdot,z)$, $\tilde S^S_j(\cdot,y)$, $E^S\varphi$ are $C^1$ and equal $H^Y(\cdot,z)$, $S^Y_j(\cdot,y)$, $H^Y\varphi$ on $\Lambda_Y$.
\end{enumerate}
\end{lemma}
\begin{proof}
(a) \Cref{rot:hyb:lem:ESH}(iii) with $F=H^Y(\cdot,z)$ and $|y_\sigma-r|\le\sqrt3\eta'$. (b) In the unit scale of level $j$ put $F(y')\defeq S^Y_j(y',y)$, resp.\ $\partial_yS^Y_j(y',y)$ ($\partial_y$ acts on the second argument and commutes with $I_S$). With $S^Y_j=H^Y_jC^Y_jH^{Y*}_j$, $F$ and its lattice gradient in $y'$ are sums $\sum_{z,z'}$ of products of three kernels with at most one lattice difference each, so with $e^{-\delta_0(|y'-z|+|z-z'|+|z'-y|)}\le e^{-\delta|y'-y|}e^{-(\delta_0-\delta)(|y'-z|+|z-z'|)}$ they are $\le C^3K_\delta e^{-\delta|y'-y|}$, $K_\delta\defeq(\sum_{v\in\Z^3}e^{-(\delta_0-\delta)(|v|-\sqrt3)})^2$. \Cref{rot:hyb:lem:ESH}(iii) at the spacing $\eta'/\ell_j\le\frac15$ adds $e^{\sqrt3\delta/5}$ and the factor $\frac{3\sqrt3}2$ for the gradient. Rescale by \cite[(2.20)]{K1}: kernel $\ell_j^{-1}$, each derivative $\ell_j^{-1}$. (c) \Cref{rot:hyb:lem:ESH}(iv). Only lattice inputs with at most one derivative occur; no same-argument second derivative of a block-spin kernel is used.
\end{proof}

\begin{lemma}[multiplier form at the $X$ sites]\label{rot:hyb:lem:ADMS}
Unit scale of level $r$; $x\in\Gamma_1$, $s\in\R^3$, $|s|\le\sqrt{13}\eta'$. In the notation of \Cref{rot:lem:H},
\[
E^S(x+s,z)=(2\pi)^{-3}\!\int_{V_\eta}\!dq\,e^{\mathrm iq\cdot(x-z)}u_{r-1}(q)M^Y_S(q;s)A^Y(q),\qquad M^Y_S(q;s)\defeq\sum_{l\in L^Y_2}\frac{w_l(q;s)\,u_Y(q+l)}{\Delta'_Y(q+l)},\qquad
w_l(q;s)\defeq\sum_{L\in\Lambda_Y^*}\chi(q+l+L)e^{\mathrm i(q+l+L)\cdot s},
\]
$\chi(Q)\defeq\prod_\alpha\hat\psi(\eta'Q\cdot f_\alpha)$, $\hat\psi(t)=\int\psi(\sigma)e^{-\mathrm it\sigma}d\sigma=12(2-2\cos t-t\sin t)/t^4$, $\Lambda_Y^*=\frac{2\pi}{\eta'}R^{-1}\Z^3$; and on $V_\eta$: $|w_l|\le W^3\le1.20$ for $l\ne0$, $|w_0(q;s)-e^{\mathrm iq\cdot s}|\le0.106\,\eta|q|$, $|M^Y_S(q;s)-\Delta_R^{-1}|\le207\,\eta|q|^{-1}$. Consequently, for $x'=x+s\in\Lambda_X$ in the tile of $x$, $|H^X(x',z)-E^S(x',z)|\le C^S_EJ_1\eta$ with $C^S_E\defeq380a+0.2a(C_*/20+8.3a)$ ($C_*$ of \Cref{rot:lem:H}), and the same at level $j$ in the unit scale of level $j$.
\end{lemma}
\begin{proof}
Poisson summation over $\Lambda_Y$ ($\hat\Psi(Q)=\eta'^3\chi(Q)$): $\sum_y\Psi(r-y)e^{\mathrm iQ\cdot y}=\sum_L\chi(Q+L)e^{\mathrm i(Q+L)\cdot r}$. Inserted in $H^Y(y,z)=(2\pi)^{-3}\int_{V_{\eta'}}e^{\mathrm iQ\cdot(y-z)}h(Q)dQ$ ($h$ periodic under $\Lambda_Y^*$, $e^{\mathrm iL\cdot z}=1$) this gives $E^S(r,z)=(2\pi)^{-3}\int_{\R^3}e^{\mathrm iQ\cdot(r-z)}\chi(Q)h(Q)dQ$; splitting $\R^3$ into the translates $V_\eta+l+L$ and using $e^{\mathrm i(l+L)\cdot x}=1$ gives the formula. Bounds: $\psi\ge0$ gives $|\hat\psi|\le\hat\psi(0)=1$; $\hat\psi(2\pi n)=0$ for $n\ne0$; for $|\tau|\le\pi/5$ and $n\ne0$, $|2-2\cos\tau-(2\pi n+\tau)\sin\tau|\le\tau^4/12+2\pi|n||\tau|\le1.004\cdot2\pi|n||\tau|$ and $|2\pi n+\tau|\ge0.9\cdot2\pi|n|$, so $|\hat\psi(2\pi n+\tau)|\le0.074|\tau|/|n|^3$ and $T(\tau)\defeq\sum_{n\ne0}|\hat\psi(\tau+2\pi n)|\le2\zeta(3)\cdot0.074|\tau|\le0.18|\tau|$ (below the cruder $0.213|\tau|$ is used); $1-\hat\psi(t)\le t^2/15$ on $[0,\pi]$ (the series $1-t^2/15+t^4/560-\dots$); $W\defeq\sup_t\sum_n|\hat\psi(t+2\pi n)|\le1.06$ (a one-dimensional bound: the sum is $1$ at the integers and its oscillation on a period is controlled by the bound on $T$). For $l\ne0$: $|w_l|\le\prod_\alpha\sum_n|\hat\psi(t_\alpha+2\pi n)|\le W^3$, $t_\alpha\defeq\eta'(q+l)\cdot f_\alpha$. For $l=0$ ($|t_\alpha|\le\pi/5$ on $V_\eta$, $t\defeq\eta'q$): $|w_0-e^{\mathrm iq\cdot s}|\le(1-\chi(q))+[\prod_\alpha(1+T(t_\alpha))-1]\le|t|^2/15+(e^{0.213\sum_\alpha|t_\alpha|}-1)\le(0.073+0.454)|t|\le0.106\,\eta|q|$ (convexity on $\sum|t_\alpha|\le3\pi/5$). Main term: $|w_0u_Y-1|\le0.106\eta|q|+|q||s|+(1-u_Y)\le(0.106+0.722+0.218)\eta|q|$, over $\Delta'_Y\ge\frac4{\pi^2}|q|^2$: $\le2.59\eta|q|^{-1}$; $|\Delta'^{-1}_Y-\Delta_R^{-1}|\le1.12\eta|q|^{-1}$ (proof of \Cref{rot:lem:H}); the $124$ aliases: $\le124\,W^3\eta^2/4\le37.2\eta^2\le203\eta|q|^{-1}$ ($|q|\le\sqrt3\pi/\eta$). So $|M^Y_S-\Delta_R^{-1}|\le207\eta|q|^{-1}$, and with $|M^X-\Delta_R^{-1}|\le172.1\eta|q|^{-1}$ at the same offset $s$ (proof of \Cref{rot:lem:H}), $|M^X-M^Y_S|\le380\eta|q|^{-1}$. Then the multiplier step of \Cref{rot:lem:H} verbatim: $|M^XA^X-M^Y_SA^Y|\le380a\eta|q|^{-1}+0.2a\eta^2A^Y|M^Y_S|\le C^S_E\eta|q|^{-1}$; integrate against $|u_{r-1}|$.
\end{proof}

\begin{corollary}[decaying forms at the $X$ sites]\label{rot:hyb:cor:DF}
Unit scale of level $r$; $x'\in\Lambda_X$ in the tile of $x\in\Gamma_1$, $x''=\tau x'$; $y'\in\Lambda_X$, $y''=\tau y'$.
\begin{enumerate}[label=\textup{(\alph*)},leftmargin=*]
\item (lifts) $|H^X(x',z)-E^S(x',z)|\le C(C^S_EJ_1\eta)^{1/2}e^{-\delta|x-z|/2}$ and $|H^Y(x'',z)-E^S(x',z)|\le C(375a\,J_1\eta)^{1/2}e^{-\delta|x-z|/2}$; the same for the photon lifts and for $E^S_j$ at level $j$ (unit scale of level $j$, $\eta\to\eta_j=5^{1-j}$).
\item (slices) for $2\le j<r$: $|S^X_j(x',y')-\tilde S^S_j(x',y'')|\le C(\eta'/\ell_j)^{1/2}\ell_j^{-1}e^{-\delta|x'-y'|/\ell_j}$.
\item (slices, the derivative at the lattice end) for $2\le j<r$: $|\partial_{y''}S^Y_j(x'',y'')-\partial_{y''}\tilde S^S_j(x',y'')|\le C(\eta'/\ell_j)^{1/2}\ell_j^{-2}e^{-\delta|x'-y'|/\ell_j}$.
\end{enumerate}
For $j\in\{0,1\}$ each term is bounded separately by \Cref{rot:hyb:lem:ESB} and the per-flow bounds of \Cref{rot:cor:unif}(a); there is no gain ($\ell_j\le5\eta'$).
\end{corollary}
\begin{proof}
(a) \Cref{rot:hyb:lem:ADMS} gives the sup bound $C^S_EJ_1\eta$; the second form is the same computation with $M^Y(q)$ at the cube offset $r_Y=x''-x$ in place of $M^X$: $|w_0u_Y-e^{\mathrm iq\cdot r_Y}u_Y|\le(0.106+c_\tau/5)\eta|q|$ and $\sum_{l\ne0}|w_l-e^{\mathrm i(q+l)\cdot r_Y}|\le124(W^3+1)$, giving $\le375a\,J_1\eta$. The decay is \Cref{rot:hyb:lem:ESB}(a) for $E^S$ and the separate bounds $|H^F(x^F,z)|\le Ce^{-\delta|x^F-z|}$ of \Cref{rot:lem:H}; take the geometric mean ($|x'-x|\le\sqrt{13}\eta'$). (b) In the unit scale of level $j$,
\[
S^X_j(x',y')-\tilde S^S_j(x',y'')=\sum_{z,z'}\big[(H^X_j-E^S_j)(x',z)\,C^X_j(z,z')\overline{H^X_j(y',z')}+E^S_j(x',z)(C^X_j-C^Y_j)(z,z')\overline{H^X_j(y',z')}+E^S_j(x',z)C^Y_j(z,z')\overline{(H^X_j(y',z')-H^Y_j(y'',z'))}\big]:
\]
(a), the bound $|C^X_j-C^Y_j|\le C\eta_je^{-\delta|z-z'|/2}$ of the proof of \Cref{rot:lem:slice}, \Cref{rot:hyb:lem:ESB}(a) and the decaying form of \Cref{rot:lem:H} at matched lattice points; sum the exponentials and rescale ($\eta_j=5\eta'/\ell_j$). This is \Cref{rot:lem:slice} with $E^S_j$ in place of $H^Y_j\circ\tau$ in the first argument. (c) The difference is $\sum_{z,z'}(H^Y_j(x'',z)-E^S_j(x',z))C^Y_j(z,z')\overline{\partial H^Y_j(y'',z')}$: the second form of (a), \cite[(3.6)]{K1}, and \cite[(3.7)]{K1} for the lattice gradient at $y''$; one more factor $\ell_j^{-1}$ on rescaling. (Alternatively, by \Cref{rot:hyb:lem:ESB}(c) the difference is $\cA(x'')-\cA(x')$ for $\cA\defeq\partial_{y''}\tilde S^S_j(\cdot,y'')$, so it is $\le c_\tau\eta'\sup|\nabla\cA|\le C(\eta'/\ell_j)\ell_j^{-2}e^{-\delta|x'-y'|/\ell_j}$ by \Cref{rot:hyb:lem:ESB}(b).)
\end{proof}

% r03_hyb_dslice: the slice derivative at matched points (the input of the one-derivative clause of the contour lines).

\begin{lemma}[the slice derivative at matched points]\label{rot:hyb:lem:DS}
Let $\alpha\in(0,1)$ be the H\"older exponent of \cite[(3.8)]{K1} (as in \Cref{rot:thm:H1S}) and put
\[
\theta_\partial\defeq\frac{\alpha}{2(2+\alpha)},\qquad\text{so that}\qquad\theta_1=\frac{\alpha^2}{4(1+\alpha)^2}\le\theta_\partial\le\tfrac14 .
\]
For $2\le j<r$, in the unit scale of level $r$, let $\nabla^X_yS^X_j(x',y')\defeq\sum_\mu\partial^X_\mu S^X_j(x',y')\,e_\mu$ be the vector of forward lattice differences of the $X$ slice in its lattice end $y'$, and $\nabla^Y_yS^Y_j(x'',y'')\defeq\sum_\nu\partial^Y_\nu S^Y_j(x'',y'')\,f_\nu$ the same for the $Y$ slice along the $Y$ frame $f_\nu=R^{-1}e_\nu$, both read as vectors in the common frame (\Cref{rot:lem:vert}(iii)). Then for matched $(x',y')$, $(x'',y'')=(\tau x',\tau y')$,
\[
\big|\nabla^X_yS^X_j(x',y')-\nabla^Y_yS^Y_j(x'',y'')\big|\ \le\ C\Big(\frac{\eta'}{\ell_j}\Big)^{\theta_\partial}\ell_j^{-2}e^{-\delta|x'-y'|/\ell_j},
\]
with $C=C(\alpha)$ uniform in $j$, $r$ and the torus; the same for the photon slices. For $j\in\{0,1\}$ each gradient is bounded separately by $C\ell_j^{-2}e^{-\delta|x'-y'|/\ell_j}$ (\cite[(3.63)]{K1} for $Y$, \Cref{rot:hyb:lem:A0}(a) for $X$) and $(\eta'/\ell_j)^{\theta_\partial}\ge5^{-1/4}$: there is no gain, and the factor is paid by a degree loss, as in \Cref{rot:cor:unif}(a).
\end{lemma}
\begin{proof}
Write $D_j\defeq S^X_j-S^Y_j\circ\tau$ as in \Cref{rot:cor:unif}, so that $|D_j(p',q')|\le C_D(\eta'/\ell_j)^{1/2}\ell_j^{-1}e^{-\delta|p'-q'|/\ell_j}$ for every pair of matched points (\Cref{rot:lem:slice}; $j\ge2$). The per-flow inputs are the gradient bound $|\nabla_yS^F_j(x,y)|\le C\ell_j^{-2}e^{-\delta|x-y|/\ell_j}$ and the H\"older clause of the gradient in the lattice end,
\[
\big|\nabla_yS^F_j(x,y)-\nabla_yS^F_j(x,q)\big|\le C'\ell_j^{-2}\Big(\frac{|y-q|}{\ell_j}\Big)^{\alpha'}e^{-\delta\min(|x-y|,|x-q|)/\ell_j}\qquad(|y-q|\le2\ell_j),\qquad\alpha'\defeq\alpha/2,
\]
for both flows. For $F=Y$ this is \cite[(3.63)]{K1} (\cite[Prop.~3.7]{K1}, from \cite[Thm.~3.3, (3.8)]{K1} by scaling) with King's constant, and $\alpha'<\alpha$ costs nothing on $|y-q|\le2\ell_j$. For $F=X$ it is \Cref{rot:hyb:lem:A0}(a) read in the lattice end ($S^X_j=H^X_jC^X_jH^{X*}_j$ is symmetric), whose factor $C_\alpha\log^{3\alpha}(2+\ell_j/\eta')\,(|y-q|/\ell_j)^{\alpha}$ is $\le C'(|y-q|/\ell_j)^{\alpha'}$: with $t\defeq|y-q|/\ell_j\in[\eta'/\ell_j,2]$ one has $\log(2+\ell_j/\eta')\le\log(2+1/t)$, and $t^{\alpha/2}\log^{3\alpha}(2+1/t)$ is bounded on $(0,2]$ (the device of \Cref{rot:hyb:rem:Hol}).

Fix a unit vector $e_\mu$ of the $X$ frame and a multiple $\rho=N\eta'$ of $\eta'$ with $\eta'\le\rho\le\ell_j$, to be chosen, and put $p'\defeq y'+\rho e_\mu\in\Lambda_X$, $p''\defeq\tau p'$. All points below lie within $2\ell_j$ of $y'$, so the decay factors of the inputs differ from $e^{-\delta|x'-y'|/\ell_j}$ by constants, which are not written.

\emph{The $X$ side.} Telescoping the forward differences along the segment, $S^X_j(x',p')-S^X_j(x',y')=\sum_{i=0}^{N-1}\eta'\,\partial^X_\mu S^X_j(x',y'+i\eta'e_\mu)$, and replacing each difference by the one at $y'$ with the H\"older clause,
\[
\Big|\rho^{-1}\big[S^X_j(x',p')-S^X_j(x',y')\big]-e_\mu\cdot\nabla^X_yS^X_j(x',y')\Big|\le C\ell_j^{-2}\Big(\frac{\rho}{\ell_j}\Big)^{\alpha'}e^{-\delta|x'-y'|/\ell_j}.
\]

\emph{The $Y$ side.} $p''-y''$ is a vector of the $Y$ lattice, $p''-y''=\sum_\nu n_\nu\eta'f_\nu$ with $n_\nu\eta'=(p''-y'')\cdot f_\nu$, and $p''-y''=\rho e_\mu+w$ with $w=(\tau p'-p')-(\tau y'-y')$, $|w|\le2c_\tau\eta'$ (\Cref{rot:def:match}). Join $y''$ to $p''$ by the lattice path that takes the $|n_1|$ steps along $\pm f_1$ first, then those along $\pm f_2$, then those along $\pm f_3$; it has at most $\sqrt3(\rho+2c_\tau\eta')/\eta'+3$ steps, and every point of it is within $\rho+2c_\tau\eta'+\eta'$ of $y''$. Telescoping along the path (a step along $-f_\nu$ from $q$ is $-\eta'\,\partial^Y_\nu S^Y_j(x'',q-\eta'f_\nu)$) and replacing each difference by the one at $y''$ with the H\"older clause,
\[
S^Y_j(x'',p'')-S^Y_j(x'',y'')=(p''-y'')\cdot\nabla^Y_yS^Y_j(x'',y'')+\mathrm{Rem},\qquad|\mathrm{Rem}|\le C\rho\,\ell_j^{-2}\Big(\frac{\rho+2c_\tau\eta'}{\ell_j}\Big)^{\alpha'}e^{-\delta|x'-y'|/\ell_j},
\]
and $(p''-y'')\cdot\nabla^Y_yS^Y_j=\rho\,e_\mu\cdot\nabla^Y_yS^Y_j+w\cdot\nabla^Y_yS^Y_j$ with $|w\cdot\nabla^Y_yS^Y_j(x'',y'')|\le2c_\tau\eta'\,C\ell_j^{-2}e^{-\delta|x'-y'|/\ell_j}$. Hence, with $(a+b)^{\alpha'}\le a^{\alpha'}+b^{\alpha'}$,
\[
\Big|\rho^{-1}\big[S^Y_j(x'',p'')-S^Y_j(x'',y'')\big]-e_\mu\cdot\nabla^Y_yS^Y_j(x'',y'')\Big|\le C\ell_j^{-2}\Big[\Big(\frac{\rho}{\ell_j}\Big)^{\alpha'}+\Big(\frac{\eta'}{\ell_j}\Big)^{\alpha'}+\frac{\eta'}{\rho}\Big]e^{-\delta|x'-y'|/\ell_j}.
\]

\emph{The comparison.} The pairs $(x',y'),(x'',y'')$ and $(x',p'),(x'',p'')$ are matched, so
\[
\rho^{-1}\Big|\big[S^X_j(x',p')-S^X_j(x',y')\big]-\big[S^Y_j(x'',p'')-S^Y_j(x'',y'')\big]\Big|\le\rho^{-1}\big(|D_j(x',p')|+|D_j(x',y')|\big)\le C\Big(\frac{\eta'}{\ell_j}\Big)^{1/2}\frac{\ell_j}{\rho}\,\ell_j^{-2}e^{-\delta|x'-y'|/\ell_j}.
\]
Adding the three displays, with $s\defeq\rho/\ell_j\in[\eta'/\ell_j,1]$,
\[
\big|e_\mu\cdot\big(\nabla^X_yS^X_j(x',y')-\nabla^Y_yS^Y_j(x'',y'')\big)\big|\le C\ell_j^{-2}e^{-\delta|x'-y'|/\ell_j}\Big[\Big(\frac{\eta'}{\ell_j}\Big)^{1/2}\frac1s+s^{\alpha'}+\Big(\frac{\eta'}{\ell_j}\Big)^{\alpha'}+\frac{\eta'}{\ell_j}\,\frac1s\Big].
\]
Choose $s$ as the smallest multiple of $\eta'/\ell_j$ that is $\ge(\eta'/\ell_j)^{1/(2(1+\alpha'))}$; since $1/(2(1+\alpha'))<1$ it lies in $[\eta'/\ell_j,1]$ and is $\le2(\eta'/\ell_j)^{1/(2(1+\alpha'))}$. The first two terms are then $\le C(\eta'/\ell_j)^{\alpha'/(2(1+\alpha'))}$, and so are the last two, because $\alpha'\ge\alpha'/(2(1+\alpha'))$ and $1-\frac1{2(1+\alpha')}=\frac{1+2\alpha'}{2(1+\alpha')}\ge\frac{\alpha'}{2(1+\alpha')}$. With $\alpha'=\alpha/2$, $\frac{\alpha'}{2(1+\alpha')}=\frac{\alpha}{2(2+\alpha)}=\theta_\partial$. Summing over the three $e_\mu$ gives the vector bound; the photon slices are the scalar ones times the identity. The inequalities $\theta_1\le\theta_\partial\le\frac14$ are $\alpha(2+\alpha)\le2(1+\alpha)^2$ and $2\alpha\le2+\alpha$.
\end{proof}

% r03_hyb_gain: proof of \Cref{rot:prop:BD}, part VII: the gain of every difference element at a non-zero gauge field (the lift of a constant field and the phases, the contour-line comparison, contour vertices, contour lines, double contour lines, the list); this completes part (b).

\paragraph{The difference elements at a non-zero gauge field.}
After the constant-background reduction of \Cref{rot:lem:CB}, applied to both flows on their infinite lattices with the torus lines as exact periodisations, every propagator that is compared between the flows is an $A=0$ line $G_r(\Lambda_F,0)=\sum_{j<r}S^F_j$, exactly the kernels of \Cref{rot:prop:fc,rot:lem:H,rot:lem:slice,rot:cor:unif}, with all vertices in the region $F_v$, and the background enters only through the explicit coarse fields $A_r(w_1)\cdots A_r(w_q)$ of \cite[(3.55)]{K1} (the same for both flows), the constant-background phases \cite[(3.46)]{K1}, the contour vertices \cite[(3.24)--(3.26)]{K1} with their contour lines \cite[(3.27)]{K1}, the external lifted fields $\tilde A^{(r),F}$, and the per-flow remainders, which are bounded by \Cref{rot:prop:BD}(a) at the regular fields of \Cref{rot:lem:CB}(R). The elements that are not $A=0$ lines are compared here.

\begin{lemma}[the lift of a constant field; the phases at matched vertices]\label{rot:hyb:lem:bph}
The lift \cite[(2.18)]{K1} of a constant coarse field $B$ to either flow is the constant fine field $\vr_rB$, $\vr_r\defeq a_r/(a_r+\lmu^2\delta_r^2)\le1$, the same number in both flows (\cite[(3.43)]{K1}). For matched $v'\in\Lambda_X$, $v''=\tau v'$ and the common root $y\in\Gamma_r$,
\[
\big|U(\vr_rB(\Gamma^X_{y,v'}))-U(\vr_rB(\Gamma^Y_{y,v''}))\big|\le e_r|B|\,c_\tau\eta'\le e\,p_r\lmu^{-1}c_\tau\,\delta_r^{-1/2}\eta'
\]
for $B$ in the small-field class.
\end{lemma}
\begin{proof}
$Q_F^*\mathbf 1_{\Gamma_r}=\mathbf 1_{\Lambda_F}$ for the frame-covariant photon average (\Cref{rot:lem:tile}(iii): $(\hat Q^*\psi)(x)=\sum_\alpha\psi_\alpha(y(x))f_\alpha$, which for the frame components $\psi_\alpha=f_\alpha\cdot B$ of a constant $B$ is $B$), and $(-\Delta+\lmu^2\delta_r^2)\mathbf 1=\lmu^2\delta_r^2\mathbf 1$, so $(M+a_rQ_F^*Q_F)\mathbf 1=(\lmu^2\delta_r^2+a_r)\mathbf 1$ and $a_rG^{F,\mathrm{ph}}_rQ_F^*B=\vr_rB$ for both flows. Both phases are $e^{\mathrm ie_r\vr_rB\cdot(v-y)}$ with $v=v'$, $v''$ (\Cref{rot:hyb:lem:alg}(iii)); $|v'-v''|\le c_\tau\eta'$, $\vr_r\le1$, $|B|\le p_r(\lmu\delta_r)^{-1}$ (\cite[(3.2)]{K1}), $e_r=e\delta_r^{1/2}$.
\end{proof}

\begin{lemma}[contour-line comparison]\label{rot:hyb:lem:CL}
Let $\cA$ be a $C^1$ vector field on a neighbourhood of the level-$r$ block $z+R^{-1}[-\frac12,\frac12]^3$ with $|\cA|\le P_0$, $|\nabla\cA|\le P_1$ there. Let $a^X_b\defeq\cA(x_b)\cdot e_\mu$ for the $\Lambda_X$-bond $b=(x_b,x_b+\eta'e_\mu)$ and $a^Y_{b''}\defeq\cA(x_{b''})\cdot f_\alpha$ for the $\Lambda_Y$-bond $b''=(x_{b''},x_{b''}+\eta'f_\alpha)$ (site sampling, King's identification of bond and site fields). Then for $x'\in B^X_r(z)$, $x''=\tau x'$,
\[
\big|a^X(\Gamma^{X,(r)}_{z,x'})-a^Y(\Gamma^{Y,(r)}_{z,x''})\big|\ \le\ 10\,(P_0+P_1)\,\eta' .
\]
The same holds for the partial contours $\Gamma^{(j+1)}_{x_{j+1},x}$ of \cite[(2.19)]{K1}, with $P_0,P_1$ taken on the level-$(j+1)$ block, the sliver area $\le1.5\eta'\ell_{j+1}+33\eta'^2$ in the flux term and the lengths $\le1.5\ell_{j+1}$ ($Y$), $\le2.1\ell_{j+1}$ ($X$) in the Riemann term: the bound is $\le(7.1P_0+(6.6\ell_{j+1}+66\eta')P_1)\eta'\le C(P_0+P_1\ell_{j+1})\eta'$.
\end{lemma}
\begin{proof}
\emph{Riemann step.} For a bond $b$ of length $\eta'$ sampled at its initial point, $|\eta'\,\cA(x_b)\cdot\hat b-\int_b\cA\cdot dl|\le\eta'\cdot\eta'P_1$. Summing over the $\le2.1/\eta'$ bonds of $\Gamma^X$ and the $\le1.5/\eta'$ bonds of $\Gamma^Y$: $|a^X(\Gamma^X)-\int_{\Gamma^X}\cA|+|a^Y(\Gamma^Y)-\int_{\Gamma^Y}\cA|\le3.6\,P_1\eta'$. \emph{Stokes step.} $\int_{\Gamma^X}\cA-\int_{\Gamma^Y}\cA=\oint_{\partial\Sigma}\cA-\int_{[x',x'']}\cA=\int_\Sigma\curl\cA\cdot dS-\int_{[x',x'']}\cA$ with the surface $\Sigma=\Sigma(z,x')$ of \Cref{rot:hyb:lem:len}(iii) (piecewise planar, so Stokes' theorem applies piece by piece and the interior edges cancel); $|\curl\cA|\le2P_1$, $\mathrm{area}(\Sigma)\le1.5\eta'+33\eta'^2\le3\eta'$ for $r\ge2$, and $|[x',x'']|\le c_\tau\eta'$. Total $\le(3.6+6)P_1\eta'+7.1P_0\eta'\le10(P_0+P_1)\eta'$. For partial contours the same decomposition applies to the sub-polygon up to level $j+1$ (\Cref{rot:hyb:lem:len}(i),(iii)).
\end{proof}

\begin{lemma}[contour vertices at the lifted background]\label{rot:hyb:lem:CV}
Let $A'^F=\vartheta\tilde A^{(r),F}+za_F$ be the reduced background of the flow $F$ (\Cref{rot:lem:CB}). Then for matched $(y,x')$, $(y,x'')$,
\[
\big|U(A'^X(\Gamma^{X,(r)}_{y,x'}))-U(A'^Y(\Gamma^{Y,(r)}_{y,x''}))\big|\ \le\ e_r\big|A'^X(\Gamma^X_{y,x'})-A'^Y(\Gamma^Y_{y,x''})\big|\ \le\ C\,e_r\,p_r^2\,\eta^{1/2},
\]
and hence the kernels of $F^X_{2,r}(A'^X,0)$ and of $F^Y_{2,r}(A'^Y,0)\circ\tau$ (\cite[(3.51)]{K1}, the contour vertices (3.24)--(3.26)) differ entrywise by $\le Ce_rp_r^2\eta^{1/2}$ against their own size $\le Ce_rp_r^2$: a relative gain $\eta^{1/2}$, and an absolute gain $\eta^{1/2}e\,p_r^2\delta_r^{1/2}$. The products (3.24)--(3.26)$\times$(3.26) at the same $y$ inherit the bound.
\end{lemma}
\begin{proof}
Let $\cA\defeq\vartheta\,E^S_{\mathrm{ph}}\tilde A_r+za^{\mathrm c}$, where $E^S_{\mathrm{ph}}\tilde A_r\defeq I_S[H^{Y,\mathrm{ph}}\tilde A_r]$ is the cubic-Hermite extension of the photon lift of $\tilde A_r$ along the flow $Y$ (\Cref{rot:hyb:def:ES}; componentwise in any frame) and $a^{\mathrm c}$ the smooth Cameron--Martin form (\Cref{rot:lem:CM}). $\cA$ is $C^1$ (\Cref{rot:hyb:lem:ESH}(iv)). By \Cref{rot:hyb:lem:ESH}(iii) and \Cref{rot:lem:CB}(R), $|E^S_{\mathrm{ph}}\tilde A_r(r)|\le\max_\sigma|\tilde A^{(r),Y}(y_\sigma)|\le Cp_r(|r-x_0|+\sqrt3\eta')$ and $|\nabla E^S_{\mathrm{ph}}\tilde A_r(r)|\le\frac{3\sqrt3}2\max_{\mathrm{edges}}|\partial\tilde A^{(r),Y}|\le Cp_r$; with \Cref{rot:lem:CM} for $a^{\mathrm c}$, $P_0,P_1\le Cp_r^2$ on the vertex region (of diameter $\le Cp_r$). Then $a^Y_{b''}=\cA(x_{b''})\cdot f_\alpha$ exactly (\Cref{rot:hyb:lem:ESB}(c); \Cref{rot:lem:CM} for $a_Y$), while
\[
|a^X_b-\cA(x_b)\cdot e_\mu|\le\vartheta\sum_z|H^{X,\mathrm{ph}}(x_b,z)-E^S_{\mathrm{ph}}(x_b,z)||\tilde A_r(z)|+|a_{X,\mu}(x_b)-a^{\mathrm c}\cdot e_\mu|\le C(C^S_EJ_1\eta)^{1/2}p_r^2+C\eta'^2
\]
by \Cref{rot:hyb:cor:DF}(a) at the $X$ site $x_b$ (the decaying form; $|\tilde A_r(z)|\le p_r(|z-x_0|+1)$, the contributions from far away as in \Cref{rot:lem:CB}(R)) and \Cref{rot:lem:CM}. Summing the lift mismatch with weight $\eta'$ over the $\le2.1/\eta'$ bonds of $\Gamma^X$ gives $\le C\eta^{1/2}p_r^2$, and \Cref{rot:hyb:lem:CL} gives $\le10(P_0+P_1)\eta'\le Cp_r^2\eta'$ for the rest. Multiply by $e_r$ for the phases ($|U(a)-U(a')|\le e_r|a-a'|$), and by $\eta'^3$ per summand of (3.51).
\end{proof}

\begin{lemma}[{the contour lines \cite[(3.27)]{K1}}]\label{rot:hyb:lem:CLN}
For $0\le j<r$ let $S^{F,\mathrm{ph}}_j$ be the stage-$j$ photon slice kernel (an $A=0$ kernel: the scalar slice $S^F_j$ times the identity, frame-free), and let $\cL^F_j(x;y)\defeq\sum_{b_1\in\Gamma^{F,(j+1)}_{x_{j+1},x}}\eta'\,S^F_j(x_{b_1},y)\,\hat b_1\in\R^3$ ($\hat b_1$ the direction of $b_1$, $y\in\Lambda_F$) be King's contour line (3.27) as a vector: its component along the direction of the photon bond $b=(y,y+\eta'\hat e)$ is $G_{(j)}(\Gamma^{(j+1)}_{x_{j+1},x},b)$ of \cite[(3.27)]{K1}. Lines are compared as vectors (\Cref{rot:lem:vert}(iii)). Then for matched $(x',y')$, $(x'',y'')=(\tau x',\tau y')$,
\[
\big|\cL^X_j(x';y')-\cL^Y_j(x'';y'')\big|\ \le\ C\Big(\frac{\eta'}{\ell_j}\Big)^{\theta}e^{-\delta\dist(\Gamma^{(j+1)},y')/\ell_j},\qquad\theta=\tfrac12,
\]
against the unweakened bound $Ce^{-\delta\dist/\ell_j}$ of each line: a slice-$j$ line weakened by the degree $\theta$ in the sense of \Cref{rot:thm:H1S}, with the prefactor $\eta'^\theta$. With one lattice derivative at the lattice end $y$ (the gradients of the two lines read as vectors in the common frame, \Cref{rot:hyb:lem:DS}) the same holds with one more factor $\ell_j^{-1}$ and, in place of $(\eta'/\ell_j)^{1/2}$, the gain $(\eta'/\ell_j)^{\theta_\partial}$ of \Cref{rot:hyb:lem:DS}, $\theta_\partial=\alpha/(2(2+\alpha))\ge\theta_1$: a slice-$j$ line with one derivative, weakened by the degree $\theta_\partial$.
\end{lemma}
\begin{proof}
Keep the second argument of the $Y$ kernel at the lattice point $y''$ throughout (no extension in the second argument is used) and write
\[
\cL^X_j(x';y')-\cL^Y_j(x'';y'')=\sum_{b_1\in\Gamma^X}\eta'\big[S^X_j(x_{b_1},y')-\tilde S^S_j(x_{b_1},y'')\big]\hat b_1
+\Big[\sum_{b_1\in\Gamma^X}\eta'\tilde S^S_j(x_{b_1},y'')\hat b_1-\sum_{b''_1\in\Gamma^Y}\eta'S^Y_j(x_{b''_1},y'')\hat b''_1\Big].
\]
\emph{First sum} ($2\le j<r$): \Cref{rot:hyb:cor:DF}(b) at the $X$ sites $x_{b_1}$, per bond $\le C(\eta'/\ell_j)^{1/2}\ell_j^{-1}e^{-\delta|x_{b_1}-y'|/\ell_j}$; with weight $\eta'$ over the $\le10.5\ell_j/\eta'$ bonds of the partial $X$ contour (\Cref{rot:hyb:lem:len}(i)): $\le C(\eta'/\ell_j)^{1/2}e^{-\delta\dist(\Gamma^X,y')/\ell_j}$. \emph{Bracket}: for a unit vector $u$, $u\cdot[\,\cdot\,]=a^X(\Gamma^X)-a^Y(\Gamma^Y)$ for the field $\cA_u(r)\defeq\tilde S^S_j(r,y'')\,u$, which is $C^1$ (\Cref{rot:hyb:lem:ESB}(c)), whose $Y$ samples are exactly $S^Y_j(x_{b''},y'')\,u\cdot f_\alpha$, and which has $P_0\le C\ell_j^{-1}e^{-\delta\dist(\Gamma,y'')/\ell_j}$, $P_1\le C\ell_j^{-2}e^{-\delta\dist(\Gamma,y'')/\ell_j}$ on the level-$(j+1)$ block (\Cref{rot:hyb:lem:ESB}(b); the block's diameter $\sqrt3\ell_{j+1}$ and $|y'-y''|\le c_\tau\eta'$ change only the constant). \Cref{rot:hyb:lem:CL} for partial contours gives $\le C(P_0+P_1\ell_{j+1})\eta'\le C(\eta'/\ell_j)e^{-\delta\dist/\ell_j}$ (flux, end segment and Riemann terms together: the kernel difference along the contour plus the flux of the slice kernel through the sliver, as \Cref{rot:cor:unif}(c) states); $u$ over an orthonormal basis gives the vector bound. The matching $y'\leftrightarrow y''$ of the second argument sits inside \Cref{rot:hyb:cor:DF}(b). \emph{The derivative version} ($2\le j<r$). With the lattice derivative on the lattice end, split the first sum at the matched $Y$ point: $\sum_{b_1\in\Gamma^X}\eta'\big[\nabla^X_yS^X_j(x_{b_1},y')-\nabla^Y_yS^Y_j(\tau x_{b_1},y'')\big]\hat b_1+\sum_{b_1\in\Gamma^X}\eta'\big[\nabla^Y_yS^Y_j(\tau x_{b_1},y'')-\nabla^Y_y\tilde S^S_j(x_{b_1},y'')\big]\hat b_1$ (gradients as vectors, the matrix $[\,\cdot\,]\hat b_1$ contracted on the line's index). The first sum is \Cref{rot:hyb:lem:DS} at the matched pairs $(x_{b_1},y')$, $(\tau x_{b_1},y'')$, per bond $\le C(\eta'/\ell_j)^{\theta_\partial}\ell_j^{-2}e^{-\delta|x_{b_1}-y'|/\ell_j}$, hence $\le C(\eta'/\ell_j)^{\theta_\partial}\ell_j^{-1}e^{-\delta\dist(\Gamma^X,y')/\ell_j}$ after the sum over the $\le10.5\ell_j/\eta'$ bonds with weight $\eta'$; the second sum is \Cref{rot:hyb:cor:DF}(c) per bond, $\le C(\eta'/\ell_j)^{1/2}\ell_j^{-1}e^{-\delta\dist/\ell_j}$ after the same sum; and the bracket is \Cref{rot:hyb:lem:CL} for the $C^1$ field $\partial_{y''}\tilde S^S_j(\cdot,y'')\,u$ (\Cref{rot:hyb:lem:ESB}(b),(c): $P_0\le C\ell_j^{-2}e^{-\delta\dist/\ell_j}$, $P_1\le C\ell_j^{-3}e^{-\delta\dist/\ell_j}$), $\le C(\eta'/\ell_j)\ell_j^{-1}e^{-\delta\dist/\ell_j}$. The weakest gain is the first sum's, $(\eta'/\ell_j)^{\theta_\partial}$. \emph{$j\in\{0,1\}$}: each line, also with the derivative, is bounded separately by $Ce^{-\delta\dist/\ell_j}$ (times $\ell_j^{-1}$ with the derivative) (\cite[(3.64)]{K1} for $Y$; \Cref{rot:hyb:lem:A0}(a) summed along $\le10.5\ell_j$ for $X$), and $(\eta'/\ell_j)^{1/2}\ge5^{-1/2}$: there is no gain and the factor $\eta'^\theta$ is paid by a degree loss, exactly as in \Cref{rot:cor:unif}(a).
\end{proof}

\begin{lemma}[double contour lines]\label{rot:hyb:lem:DCL}
For $0\le j<r$ and $x_1,x_2\in\Lambda_F$ let
\[
\cD^F_j(x_1;x_2)\defeq\sum_{b_1\in\Gamma^{F,(j+1)}_{x_{1,j+1},x_1}}\sum_{b_2\in\Gamma^{F,(j+1)}_{x_{2,j+1},x_2}}\eta'^2S^F_j(x_{b_1},x_{b_2})\,\hat b_1\cdot\hat b_2,
\]
the line joining two contour factors of \cite[(3.26)]{K1} (at one vertex with $m\ge2$, at the two factors of a product (3.26)$\times$(3.26), or at two vertices). Then $|\cD^F_j(x_1;x_2)|\le C\ell_je^{-\delta|x_{1,j+1}-x_{2,j+1}|/\ell_j}$ for both flows, and at matched points $|\cD^X_j(x_1';x_2')-\cD^Y_j(x_1'';x_2'')|\le C(\eta'/\ell_j)^{1/2}\ell_je^{-\delta|x_{1,j+1}-x_{2,j+1}|/\ell_j}$. In Ba{\l}aban's degree $D(H)=\sum_ld_l+\sum_va_v+3(|V(H)|-1)$ (\cite[(2.1)--(2.2)]{Ba3}, with the two summation points of a block product as vertices, as in \Cref{rot:cor:OF}) a double line has $d_{l_0}=+1$, and every connected all-scale-summed subgraph $H$ containing it has $D(H)\ge1$: the double line is never the binding constraint of the weakened-element bound.
\end{lemma}
\begin{proof}
$\cD^F_j(x_1;x_2)=\sum_{b_2}\eta'\hat b_2\cdot\cL^F_j(x_1;x_{b_2})$. The size bound: lengths $\le10.5\ell_j$ ($X$), $\le7.5\ell_j$ ($Y$) (\Cref{rot:hyb:lem:len}(i)), \Cref{rot:hyb:lem:A0}(a), \cite[(3.63)]{K1}; since both contours lie in open rotated cubes of side $5\ell_j$ about their roots, $e^{-\delta|x_{b_1}-x_{b_2}|/\ell_j}\le e^{5\sqrt3\delta}e^{-\delta|x_{1,j+1}-x_{2,j+1}|/\ell_j}$. For $j\ge2$ put $\cV(r)\defeq\sum_{b_1''}\eta'\,I_S[S^Y_j(x_{b_1''},\cdot)](r)\,\hat b_1''$ (extension in the second argument); then
\[
\cD^X-\cD^Y=\sum_{b_2}\eta'\hat b_2\cdot[\cL^X_j(x_1';x_{b_2})-\cL^Y_j(x_1'';\tau x_{b_2})]+\sum_{b_2}\eta'\hat b_2\cdot[\cV(\tau x_{b_2})-\cV(x_{b_2})]+[a^X(\Gamma^X_2)-a^Y(\Gamma^Y_2)]_{\cA=\cV}
\]
exactly (\Cref{rot:lem:vert}(i); $\tau x_{b_2}\in\Lambda_Y$); the three terms are bounded by \Cref{rot:hyb:lem:CLN} summed along $\Gamma^X_2$, by $c_\tau\eta'\sup|\nabla\cV|$ with $|\nabla\cV|\le C\ell_j^{-1}$ (\Cref{rot:hyb:lem:ESH}(iii), \Cref{rot:hyb:lem:ESB}(b)), and by \Cref{rot:hyb:lem:CL} for partial contours with $P_0\le C$, $P_1\le C\ell_j^{-1}$. For $j\le1$ use the size bounds and $(\eta'/\ell_j)^{1/2}\ge5^{-1/2}$. Degree: the size $C\ell_j$ against $\ell_j^{-1}$ for a slice line ($d_l=-1$) and $\ell_j^0$ for a one-ended contour line ($d_l=0$) gives $d_{l_0}=+1$ (\Cref{rot:cor:OF}: $d_l=-1+\#(\text{contour ends})$). For a connected all-scale-summed $H$ containing $l_0$: if $H\setminus l_0$ is connected (or $l_0$ is a self-loop), $D(H)=D(H\setminus l_0)+1\ge1$, because $D\ge0$ for every all-scale-summed subgraph that does not contain $l_0$ (the per-flow renormalised structure, \Cref{rot:cor:OF}); if $l_0$ is a bridge between $H_1$ and $H_2$, $D(H)=D(H_1)+D(H_2)+1+3\ge4$. So $D(H)-\theta\ge\frac12$ for $\theta=\frac12$.
\end{proof}

\begin{proposition}[every difference element at $A\ne0$]\label{rot:hyb:prop:list}
In the setting of \Cref{rot:prop:BD}(b), after \Cref{rot:lem:CB}, order the elements of a renormalised family consistently across the two flows and telescope as in \Cref{rot:prop:cp}. Every difference element $e^X_i-e^Y_i$ is of one of the following types, with the stated bound:
\begin{enumerate}[label=\textup{(\roman*)},leftmargin=*]
\item an $A=0$ line: a fine slice $S^F_j$, $j<r$ (\Cref{rot:cor:unif}(a): $C_D\eta'^{\theta}\ell_j^{-1-\theta}e^{-\delta|x'-y'|/\ell_j}$, $\theta=\frac12$), a coarse slice or $\Gamma^F_n$ (\Cref{rot:cor:unif}(b): $C\eta^{1/2}\ell_j^{-1}e^{-\delta|x'-y'|/\ell_j}$), the coarse multiplier (\Cref{rot:prop:fc}: $0.2a\eta^2$ relative), on $\Lambda_F$;
\item a contour line (3.27): \Cref{rot:hyb:lem:CLN}, gain $(\eta'/\ell_j)^{1/2}$; with one lattice derivative at its lattice end, gain $(\eta'/\ell_j)^{\theta_\partial}$, $\theta_\partial=\alpha/(2(2+\alpha))$ (\Cref{rot:hyb:lem:DS});
\item a contour vertex (3.24)--(3.26): \Cref{rot:hyb:lem:CV}, entrywise $\le Ce_rp_r^2\eta^{1/2}$, relative gain $\eta^{1/2}$;
\item an external lifted field $\tilde A^{(r),F}$, $H^F\phi_r$ or its covariant derivative: \Cref{rot:lem:H,rot:lem:grad}, $\le C\eta^{1/2}$ times the field norm;
\item a constant-background phase (3.46): \Cref{rot:hyb:lem:bph}, $\le e\,p_r\lmu^{-1}c_\tau\delta_r^{-1/2}\eta'$;
\item a hopping vertex of $V_r(A',0)$ or a vertex function (couplings, partition of unity, the Cameron--Martin form): \Cref{rot:lem:vert}(ii),(iii), unchanged;
\item a double contour line, joining two contour factors of (3.26): \Cref{rot:hyb:lem:DCL}, size $\le C\ell_j$, relative gain $(\eta'/\ell_j)^{1/2}$ on slice $j$.
\end{enumerate}
Each weakened element therefore obeys the bound of its type weakened by the degree $\theta=\frac12$ (by $\theta_\partial\ge\theta_1$ for a contour line carrying a lattice derivative, which is all \Cref{rot:thm:H1S} asks) with a constant $\le C\,p_r^{2}\,\delta_r^{-1/2}$ (the worst of (iii) and (v)); the per-flow remainders (\Cref{rot:lem:CB}(2)) contribute $(Ce_rp_r^2)^{n+1}+e^{-\delta'p_r}$ per graph.
\end{proposition}
\begin{proof}
\Cref{rot:lem:CB} lists what carries the background after the reduction; the compared lines are $A=0$ lines, the contour lines, one-ended (3.27) and double (both ends on contours), are the only lines whose shape depends on the flow, and the only background-dependent vertices are the contour vertices (3.24)--(3.26) (\cite[p.~658]{K1}: ``The remaining vertices occur when the renormalization transformation of the scalar field is expanded about $A'=0$'') and the phases (3.46). The external coarse fields $A_r(w_i)$ of (3.55) are common. The bounds are the cited lemmas.
\end{proof}

\begin{proof}[Proof of \Cref{rot:prop:BD}(b)]
The difference elements $\Delta\Gamma$, $\Delta H$, $\Delta(A^{-1})$ at regular $A$ are, after the reduction, the $A=0$ differences of \Cref{rot:cor:unif} (types (i) and (iv) of \Cref{rot:hyb:prop:list}), with the factor $\eta'^{1/2}$. The contour vertices at the reduced background differ entrywise by a relative $\eta^{1/2}$ against their size $Ce_rp_r^2$ (\Cref{rot:hyb:lem:CV}); the constant-background phases at matched vertices differ by $\le e\,p_r\lmu^{-1}c_\tau\delta_r^{-1/2}\eta'$ (\Cref{rot:hyb:lem:bph}); the per-slice pieces carried by the contour vertex, with one end on a contour (\cite[(3.27)]{K1}) or both ends on contours, obey the slice-$j$ bounds \cite[(3.63)--(3.65)]{K1} in both flows ((a) for the hybrid flow, \cite[Prop.~3.7]{K1} for the standard flow) with the relative mismatch $(\eta'/\ell_j)^{1/2}$ (\Cref{rot:hyb:lem:CLN,rot:hyb:lem:DCL}), the one-ended lines also with one lattice derivative at the lattice end, there with the relative mismatch $(\eta'/\ell_j)^{\theta_\partial}$ of \Cref{rot:hyb:lem:DS} (\Cref{rot:hyb:lem:CLN}; a double line has no lattice end). A constant background is pure gauge on the infinite lattice and gives the same phase, up to the matching segment, for every contour (\Cref{rot:hyb:lem:bph}); what would break (b) are mis-paired contours or mis-rooted blocks, both excluded by \Cref{rot:def:flows} and \Cref{rot:conv:C}(i). The gains are not uniform in the size of the constant part of the background: they carry the small-field bound \cite[(3.2)]{K1}, as King's own reduction does.
\end{proof}


\begin{proposition}[King's upper bound for the rotated lattice]\label{rot:prop:H3}
For $\Yk$ on $T$ and every real smooth periodic source $s=(zG,h)$, with the free-boundary convention of \Cref{rot:def:tori} and the centred blocks of \Cref{rot:conv:C},
\[
\Zun_{\Yk,T}(s)\le e^{C\delta_r^\sigma|T|}\int\chi_r\,e^{-S^{(r)}_T}+e^{-p(\delta_r)^2+C|T|}\,\eps^{-q},
\]
$C=C(s,|T|)$, $q=q(s)<\infty$, with $q=0$ when $h=0$.
\end{proposition}
\begin{proof}
On $T_2=R^{-1}\DJ$, which is aligned in $Y$'s frame, the passive rotation $R$ (\Cref{rot:lem:walledalg}(b)) maps $(\Yk,T_2)$ to $(\Xk,\DJ)$ with its standard frame, so the statement is King's Theorem~3.1 \cite[(3.4)]{K1}, with the additive large-field term in the form in which King uses it (\cite[(3.11)]{K1}, \cite[(4.7)]{K2}), on King's torus at block factor $5$: this is \Cref{thm:a06:BALA}(iii) with \Cref{hyp:k01:BALint}, read with the torus-convention counterterms replaced by the free-boundary ones at a constant-shifted source (\Cref{rot:prop:CT}(b)) and with corner-rooted blocks replaced by centred ones (\Cref{rot:conv:C}(i)); King's constant depends on $g$, $h$ and $|T|$ and is uniform under a shift of the scalar source by a bounded constant, which is the uniformity of \cite[Thm.~3.1]{K1} in the mass over a compact mass set (\cite[p.~655]{K1}). On $T_1=\DJ$ the frame torus of $\Yk$ is $R\DJ$, skew in $Y$'s frame at every scale; this instance is \Cref{rot:prop:SK}. For a walled source ($\int h\ne0$) the constant of the large-field term inherits \cite[(2.23)]{K1}, which fails as printed when $\int h\ne0$ (\Cref{rot:rem:log}); what the inductive proof supplies is a bounded $h$-linear $c$-number per step, so the large-field term is read with the polynomial envelope factor $\eps^{-q}$, which \Cref{rot:prop:Kforward} absorbs.
\end{proof}

% r03_r1_setting: the cumulant bound for the step Gaussian, part I: the setting, the lattice facts, the constants, and the kernel and decoupling lemmas. Input from sec/r03_fill_K5a.tex before Theorem rot:thm:R1.

\subsubsection*{The cumulant bound for the step Gaussian}

The fluctuation integral of every block-spin step of both flows has the form $\int\chi_b\,e^{V}d\mu_P$: $\mu_P$ is a centred Gaussian whose precision is coercive and has exponentially decaying off-diagonal entries, $V$ is a real polynomial with exponentially decaying coefficients, and $\chi_b$ is a product over sites of indicators of a fixed convex body scaled by the threshold $b$. \Cref{rot:thm:R1} below is the order-$\bar n$ cumulant expansion of such an integral, with an error proportional to the volume of the support of $V$ and, in the lower bound, a normalisation term proportional to $|\Lambda|$. The printed cluster-expansion lemma of Benfatto and coauthors, which King cites for this step, is stated for the nearest-neighbour Markov field with a cutoff growing with the distance to the support of the interaction; its proof uses the Markov property and the nearest-neighbour formulas for conditional Gaussians in four places. The proof below keeps its architecture (cells with corridors, four shifted pavements, the cutoff moved from the corridors into the cells) and replaces those four uses by the decoupling and conditioning lemmas of this part, which hold for any coercive precision with exponentially decaying entries.

\paragraph{The setting.} Let $\mathbb{L}$ be one of: the lattice $\Z^3$; a cubic torus $(\Z/N\Z)^3$; a skew torus $\Z^3/mG_\pm$, $m\in\N$, where $G_+\defeq\langle(40,-30,0),(30,40,0),(0,0,50)\rangle$ and $G_-$ is its mirror image under $x_2\mapsto-x_2$. These are the frame period lattices of the skew tori of \Cref{rot:prop:SK} in the units of a step lattice: $\Gamma'=R^{\mathrm T}(50MJ\,\Z^3)=MJ\,G_+$ for $R\DJ$ in $Y$'s frame and $RP_J=MJ\,G_-$ for $R\DJ$ in $X$'s frame, so $m=MJ$ divided by the spacing of the step, an integer because every block lattice contains $P_J$ (\Cref{rot:lem:adm}(iii)). Distances are quotient Euclidean distances, $|x-y|\defeq\min_g|\tilde x-\tilde y-g|$ over the period lattice (a metric), and $\dist(x,S)\defeq\min_{y\in S}|x-y|$. $\Lambda\subset\mathbb{L}$ is finite, $z=(z_x)_{x\in\Lambda}$ with $z_x\in\R^n$, and a complex Hermitian operator on $\C^\Lambda$ is read in its real form (blocks $\bigl(\begin{smallmatrix}\Re&-\Im\\\Im&\Re\end{smallmatrix}\bigr)$), which preserves the spectrum and the block operator norms. The hypotheses are:
\begin{enumerate}[label=\textup{(R\arabic*)},leftmargin=*]
\item $\mu_P\defeq N(0,P^{-1})$ with $P$ real symmetric on $(\R^n)^\Lambda$ and $P\ge\gamma_0>0$ (no upper bound is needed);
\item $\|P(x,y)\|\le c_Pe^{-\delta_P|x-y|}$ for $x\ne y$ (block operator norms); the diagonal blocks are arbitrary;
\item $V=\sum_{Y\subset J}V_Y$ with $J\subset\Lambda$, where $V_Y$ is a real polynomial in $(z_y)_{y\in Y}$ of degree at most $D$, without constant term, each of whose monomials involves every site of $Y$ (its \emph{support}); $\ell(Y)$ is the length of a shortest connected graph on a vertex set containing $Y$ (edges between points of $\mathbb{L}$, of length their distance); $|V_Y|_1$ is the sum of the moduli of the coefficients; for $\theta\ge0$ and for any polynomial $F=\sum_YF_Y$ sorted by supports in this way,
\[
\|F\|_\theta\defeq\sup_{x\in\mathbb{L}}\sum_{Y\ni x}e^{\theta\ell(Y)}|F_Y|_1,\qquad\mathsf A\defeq\|V\|_\varsigma\ \text{ for a fixed rate }\varsigma>0,\qquad V_R\defeq\sum_{Y\subset R}V_Y;
\]
\item $\chi_\beta(z)\defeq\prod_{x\in\Lambda}\mathbf 1\{z_x\in\beta K\}$ with $K\subset\R^n$ a closed symmetric convex body, $\bar B(1)\subset K\subset\bar B(c_K)$: the Euclidean ball ($c_K=1$) or Ba{\l}aban's two-component cutoff $K=\{|z^{A}|\le1,\ |z^{\phi}|\le1\}$, $c_K=\sqrt2$, which restricts the photon and the scalar fluctuation at each site separately.
\end{enumerate}
$\kappa_\mu(F_1,\dots,F_k)$ denotes the joint cumulant under a law $\mu$ and $\kappa_\mu(F;k)\defeq\kappa_\mu(F,\dots,F)$ the $k$-th cumulant, the truncated expectation $\langle F^k\rangle^T$. For $X\subset\Lambda$, $P_{XX}$ is the principal block, $G_X\defeq(P_{XX})^{-1}$ and $C\defeq P^{-1}=G_\Lambda$; for $Y\defeq\Lambda\setminus X$ the conditional law of $z_X$ given $z_Y$ under $\mu_P$ is $N(-G_XP_{XY}z_Y,\,G_X)$, by completing the square in the exponent. $\chi^X_\beta\defeq\prod_{x\in X}\mathbf 1\{z_x\in\beta K\}$.

\begin{lemma}[lattice facts]\label{rot:lem:R1lat}
\textup{(i)} $x\mapsto|x-y_0|$ and $x\mapsto\dist(x,S)$ are $1$-Lipschitz. \textup{(ii)} For every $f\ge0$ and $x\in\mathbb{L}$, $\sum_{y\in\mathbb{L}}f(|x-y|)\le\sum_{z\in\Z^3}f(|z|)$. \textup{(iii)} For $x,y\in Y$, $|x-y|\le\ell(Y)$; a shortest connected graph for $Y$ joins any two of its vertices by a path of length at most $\ell(Y)$; and $\ell(Y')\le\ell(Y)$ for $Y'\subset Y$. \textup{(iv)} For $\theta\le\varsigma$ and $r\ge0$ the part of $V$ made of the monomials of length $\ell(Y)\ge r$ has $\|\cdot\|_\theta\le\mathsf Ae^{-(\varsigma-\theta)r}$; in particular $\sum_{Y\ni x,\ \ell(Y)\ge r}|V_Y|_1\le\mathsf Ae^{-\varsigma r}$.
\end{lemma}
\begin{proof}
(i) is the triangle inequality of the quotient metric. (ii) Fix a lift $\tilde x$ and for each $y$ choose the lift $\tilde y$ nearest to $\tilde x$, so that $f(|x-y|)=f(|\tilde x-\tilde y|)$; distinct $y$ have distinct lifts, so the sum is a sub-sum of $\sum_{\tilde z\in\Z^3}f(|\tilde x-\tilde z|)$. (iii) A path in the graph is no longer than the graph, and by the triangle inequality the distance of its endpoints is at most its length; a graph for $Y$ is a graph for $Y'$. (iv) is the definition of the norm, since $e^{\theta\ell}=e^{\varsigma\ell}e^{-(\varsigma-\theta)\ell}$.
\end{proof}

\paragraph{Constants.} Put $\tilde c(\theta)\defeq\sum_{z\in\Z^3}e^{-\theta|z|}$, $\mathsf K\defeq\sum_{z\in\Z^3}|z|e^{-\delta_P|z|/2}$, $c_d\defeq\tilde c(\delta_P/2)$, and
\begin{equation}\label{rot:eq:R1const}
\mu_1\defeq\min\Big\{\frac{\delta_P}2,\frac{\gamma_0}{2c_P\mathsf K}\Big\},\quad g_0\defeq\frac2{\gamma_0},\quad\mathsf M_1\defeq\frac{2c_P}{\gamma_0}\tilde c(\mu_1/2)^2,\quad c_{\mathrm D}\defeq\frac{4c_P}{\gamma_0^2}\tilde c(\mu_1/2)^2,\quad\gamma_K\defeq\min\Big\{\frac12,\frac1{2c_K\mathsf M_1}\Big\},\quad\bar\mu\defeq\min\{\mu_1,\varsigma\}.
\end{equation}
By \Cref{rot:lem:R1lat}(ii) every sum over $\mathbb{L}$ of a function of the distance is bounded by the corresponding sum over $\Z^3$, so these constants serve on every $\mathbb{L}$. They depend on $(\gamma_0,c_P,\delta_P,c_K,\varsigma)$ only, and no lower bound on $c_P$ is used anywhere ($c_P\to0$ gives $\mu_1=\delta_P/2$, $\mathsf M_1\to0$, $\gamma_K=\frac12$).

\begin{lemma}[Combes--Thomas bound]\label{rot:lem:R1ct}
For $X\subset\Lambda$ and $r\ge0$, $\|(P_{XX}+r)^{-1}(x,y)\|\le2(\gamma_0+r)^{-1}e^{-\mu_1|x-y|}$. In particular $\|G_X(x,y)\|\le g_0e^{-\mu_1|x-y|}$ for every $X$, and $\|C(x,y)\|\le g_0e^{-\mu_1|x-y|}$.
\end{lemma}
\begin{proof}
Fix $y_0\in X$ and $f\defeq|\cdot-y_0|$, and conjugate: $P^f\defeq e^{\mu_1f}P_{XX}e^{-\mu_1f}$. Then $(P^f-P_{XX})(x,y)=P(x,y)\big(e^{\mu_1(f(x)-f(y))}-1\big)$ vanishes on the diagonal and, by (R2), \Cref{rot:lem:R1lat}(i) and $e^t-1\le te^t$ for $t\ge0$, has block norm at most $c_Pe^{-\delta_P|x-y|}\mu_1|x-y|e^{\mu_1|x-y|}\le c_P\mu_1|x-y|e^{-\delta_P|x-y|/2}$, since $\mu_1\le\delta_P/2$; by Schur's test and \Cref{rot:lem:R1lat}(ii) its row and column sums are at most $c_P\mu_1\mathsf K\le\gamma_0/2$, so $\|P^f-P_{XX}\|\le\gamma_0/2$. Hence $\Re\langle u,(P^f+r)u\rangle\ge(\gamma_0/2+r)\|u\|^2$ for every $u$ ($P_{XX}\ge\gamma_0$ as a principal block), so $\|(P^f+r)^{-1}\|\le(\gamma_0/2+r)^{-1}\le2(\gamma_0+r)^{-1}$, and $(P_{XX}+r)^{-1}(x,y_0)=e^{-\mu_1f(x)}(P^f+r)^{-1}(x,y_0)e^{\mu_1f(y_0)}$ has norm at most $2(\gamma_0+r)^{-1}e^{-\mu_1|x-y_0|}$.
\end{proof}

\begin{lemma}[conditioning]\label{rot:lem:R1sch}
Let $X\subset\Lambda$ and $Y\defeq\Lambda\setminus X$. Then $0<G_X\le C_{XX}\le\gamma_0^{-1}$, so every component variance of $N(\cdot,G_X)$ and of $\mu_P$ is at most $\gamma_0^{-1}$. The conditional mean $h\defeq-G_XP_{XY}z_Y$ obeys
\[
|h(x)|\le\mathsf M_1e^{-(\mu_1/2)\dist(x,Y)}\|z_Y\|_\infty\qquad(x\in X),
\]
in particular $\|h\|_\infty\le\beta/2$ whenever $\|z_Y\|_\infty\le c_K\gamma_K\beta$. For $x,y\in X$,
\[
\|(C-G_X)(x,y)\|\le c_{\mathrm D}\exp\Big(-\frac{\mu_1}8\big(\dist(x,Y)+\dist(y,Y)\big)-\frac{\mu_1}4|x-y|\Big).
\]
\end{lemma}
\begin{proof}
$C_{XX}=(P_{XX}-P_{XY}P_{YY}^{-1}P_{YX})^{-1}\ge P_{XX}^{-1}=G_X$, because the subtracted operator is non-negative and inversion reverses the order of positive operators; and $C_{XX}\le\|C\|\le\gamma_0^{-1}$. For the mean, \Cref{rot:lem:R1ct} and (R2) (the pairs $(u,v)\in X\times Y$ are off-diagonal) give $|h(x)|\le\sum_{u\in X,v\in Y}g_0e^{-\mu_1|x-u|}c_Pe^{-\delta_P|u-v|}|z_v|$; as $\delta_P\ge2\mu_1$ and $|x-u|+|u-v|\ge|x-v|\ge\dist(x,Y)$, the exponent is at most $-\frac{\mu_1}2\dist(x,Y)-\frac{\mu_1}2(|x-u|+|u-v|)$, and the double sum of the last factor is at most $\tilde c(\mu_1/2)^2$ by \Cref{rot:lem:R1lat}(ii). The particular case is $\mathsf M_1c_K\gamma_K\le\frac12$. For the covariance difference, row $X$ of $PC=\mathbf 1$ reads $P_{XX}C_{XX}+P_{XY}C_{YX}=\mathbf 1$, so $C_{XX}-G_X=-G_XP_{XY}C_{YX}$, with kernel $-\sum_{u\in X,v\in Y}G_X(x,u)P(u,v)C(v,y)$, bounded by $g_0^2c_P\sum_{u,v}e^{-\mu_1L}$ with $L\defeq|x-u|+|u-v|+|v-y|$ (using $\delta_P\ge\mu_1$). Since $L\ge|x-y|$, $L\ge|x-v|\ge\dist(x,Y)$ and $L\ge|v-y|\ge\dist(y,Y)$, write $e^{-\mu_1L}=e^{-\mu_1L/4}e^{-\mu_1L/4}e^{-\mu_1L/2}\le e^{-\frac{\mu_1}8(\dist(x,Y)+\dist(y,Y))}e^{-\frac{\mu_1}4|x-y|}e^{-\frac{\mu_1}2(|x-u|+|u-v|)}$ and sum the last factor over $u,v$.
\end{proof}

\begin{lemma}[decoupling across corridors]\label{rot:lem:R1dec}
Let $\Lambda=\Gamma\sqcup D$ with $D=\bigsqcup_iD_i$ and $\dist(D_i,D_j)\ge w_s$ for $i\ne j$. Let $\mu_P(\cdot\,|\,z_\Gamma)$ be the conditional law of $z_D$ given $z_\Gamma$ and $\nu_{\mathrm{loc}}(\cdot\,|\,z_\Gamma)\defeq\bigotimes_iN(-G_{D_i}P_{D_i\Gamma}z_\Gamma,\,G_{D_i})$ its decoupled version; put $\omega\defeq c_Pc_de^{-\delta_Pw_s/2}$ and suppose $\omega\le\gamma_0/2$. On $\{\|z_D\|_\infty\le\beta,\ \|z_\Gamma\|_\infty\le\beta\}$,
\[
\Big|\log\frac{d\mu_P(\cdot\,|\,z_\Gamma)}{d\nu_{\mathrm{loc}}(\cdot\,|\,z_\Gamma)}\Big|\le\omega_{\mathrm{tot}}\defeq\omega|D|\Big[\frac{\beta^2}2+\frac n{\gamma_0}+\frac{(c_P\tilde c(\delta_P)\beta)^2}{2\gamma_0^2}\Big],
\]
hence $e^{-\omega_{\mathrm{tot}}}\int F\,d\nu_{\mathrm{loc}}\le\int F\,d\mu_P(\cdot\,|\,z_\Gamma)\le e^{\omega_{\mathrm{tot}}}\int F\,d\nu_{\mathrm{loc}}$ for every $F\ge0$ supported in that set. Each factor of $\nu_{\mathrm{loc}}$ is a law of the type of \Cref{rot:lem:R1sch}, with $X=D_i$, $Y=\Lambda\setminus D_i$ and $z_{D_j}=0$ for $j\ne i$.
\end{lemma}
\begin{proof}
Put $u\defeq P_{D\Gamma}z_\Gamma$, $B\defeq\bigoplus_iP_{D_iD_i}$ and $O\defeq P_{DD}-B$ (the blocks $P_{D_iD_j}$, $i\ne j$). The two laws are $N(-P_{DD}^{-1}u,P_{DD}^{-1})$ and $N(-B^{-1}u,B^{-1})$, since $(B^{-1}u)_{D_i}=G_{D_i}P_{D_i\Gamma}z_\Gamma$. The logarithm of the density of $N(-A^{-1}u,A^{-1})$ is $\frac12\log\det(A/2\pi)-\frac12\langle z,Az\rangle-\langle z,u\rangle-\frac12\langle u,A^{-1}u\rangle$; the linear terms agree for $A=P_{DD}$ and $A=B$, so the log-ratio is $\frac12\log\det(P_{DD}B^{-1})-\frac12\langle z_D,Oz_D\rangle-\frac12\langle u,(P_{DD}^{-1}-B^{-1})u\rangle$. Every entry of $O$ is an off-diagonal entry of $P$ between points at distance at least $w_s$, so by (R2) and \Cref{rot:lem:R1lat}(ii) its row and column sums are at most $c_Pe^{-\delta_Pw_s/2}\tilde c(\delta_P/2)=\omega$, and $\|O\|\le\omega$. First term: $|\frac12\langle z_D,Oz_D\rangle|\le\frac12\beta^2\sum_{x\in D}\sum_y\|O(x,y)\|\le\frac12\beta^2|D|\omega$. Second term: $\det(P_{DD}B^{-1})=\det(\mathbf 1+B^{-1/2}OB^{-1/2})$; the symmetric operator $B^{-1/2}OB^{-1/2}$ has $n|D|$ eigenvalues, each of modulus at most $\omega/\gamma_0\le\frac12$ ($B\ge\gamma_0$ as a direct sum of principal blocks), and $|\log(1+t)|\le2|t|$ for $|t|\le\frac12$; so the term is at most $n|D|\omega/\gamma_0$. Third term: $P_{DD}^{-1}-B^{-1}=-P_{DD}^{-1}OB^{-1}$ has norm at most $\omega/\gamma_0^2$, and $|u_x|\le\sum_{v\in\Gamma}c_Pe^{-\delta_P|x-v|}\beta\le c_P\tilde c(\delta_P)\beta$, so $\|u\|^2\le|D|(c_P\tilde c(\delta_P)\beta)^2$. The three bounds add up to $\omega_{\mathrm{tot}}$.
\end{proof}

% r03_r1_cumul: the cumulant bound for the step Gaussian, part II: tails, the normalisation of the cutoff, cumulants of bounded variables, the cutoff in the exponent and in cumulants, the connected Wick expansion and the comparison of two Gaussian laws.

\begin{lemma}[tails]\label{rot:lem:R1tail}
Let $\nu=N(h,G)$ on $(\R^n)^X$ with $\|h\|_\infty\le\beta/2$ and $G\le\gamma_0^{-1}$. Then $\nu(\exists x\in X:\ |z_x|>\beta)\le\epsilon(X,\beta)\defeq2n|X|e^{-\gamma_0\beta^2/(8n)}$; since $\bar B(\beta)\subset\beta K$, also $\nu(1-\chi^X_\beta)\le\epsilon(X,\beta)$.
\end{lemma}
\begin{proof}
$|z_x|>\beta$ forces $|\zeta_x|>\beta/2$ for $\zeta\defeq z-h$, hence $|\zeta^a_x|>\beta/(2\sqrt n)$ for some component $a$; $\zeta^a_x\sim N(0,\sigma^2)$ with $\sigma^2\le\gamma_0^{-1}$, $\Pr(|N(0,\sigma^2)|>t)\le2e^{-t^2/(2\sigma^2)}$, and a union bound over the $n|X|$ components gives the claim.
\end{proof}

\begin{lemma}[the normalisation of the cutoff]\label{rot:lem:R1gci}
Let $\mu$ be a centred Gaussian on $(\R^n)^\Lambda$ all of whose component variances are at most $\gamma_0^{-1}$, and let $2e^{-\gamma_0\beta^2/(2n)}\le\frac12$. Then $\mu(\chi_\beta)\ge\exp\big(-4n|\Lambda|e^{-\gamma_0\beta^2/(2n)}\big)$.
\end{lemma}
\begin{proof}
$\{z_x\in\beta K\}\supset\{|z_x|\le\beta\}\supset\bigcap_{a\le n}\{|z^a_x|\le\beta/\sqrt n\}$. By \v{S}id\'ak's inequality (Z.~\v{S}id\'ak, J.\ Amer.\ Statist.\ Assoc.\ 62 (1967), 626--633; C.~G.~Khatri, Ann.\ Math.\ Statist.\ 38 (1967), 1853--1867): for a centred Gaussian vector $(\xi_1,\dots,\xi_m)$ and $t_1,\dots,t_m>0$, $\Pr(\bigcap_i\{|\xi_i|\le t_i\})\ge\prod_i\Pr(|\xi_i|\le t_i)$; this is the case of one slab against a product of slabs of the Gaussian correlation inequality, and it is applied here to the $n|\Lambda|$ components of the real centred vector $z$. Each factor is at least $1-2e^{-\gamma_0\beta^2/(2n)}$, and $\log(1-u)\ge-2u$ for $u\le\frac12$.
\end{proof}

\begin{lemma}[cumulants of bounded variables]\label{rot:lem:R1cum}
For any probability measure $\nu$ and real variables with $|X_j|\le M_j$, $|\kappa_\nu(X_1,\dots,X_k)|\le k!\,2^{k-1}\prod_jM_j$. If $|X|\le M$, then $\big|\log\int e^{X}d\nu-\sum_{k=1}^{\bar n}\kappa_\nu(X;k)/k!\big|\le2^{\bar n}M^{\bar n+1}\le(2M)^{\bar n+1}$.
\end{lemma}
\begin{proof}
In the M\"obius formula $\kappa_\nu(X_1,\dots,X_k)=\sum_\pi(-1)^{|\pi|-1}(|\pi|-1)!\prod_{B\in\pi}\E_\nu\prod_{j\in B}X_j$, over the set partitions $\pi$ of $\{1,\dots,k\}$, each product of moments is at most $\prod_jM_j$, and $\sum_\pi(|\pi|-1)!$ is at most the number of ordered set partitions, which is at most $k!\,2^{k-1}$ (an ordered partition is a permutation cut into consecutive blocks). For the second claim, $f(s)\defeq\log\int e^{sX}d\nu$ has $f^{(m)}(s)=\kappa_{\nu_s}(X;m)$ with $\nu_s\propto e^{sX}\nu$ a probability measure under which still $|X|\le M$; Taylor's formula at $s=0$ with the Lagrange remainder $f^{(\bar n+1)}(s_*)/(\bar n+1)!$ and the first claim give the bound.
\end{proof}

\begin{lemma}[the cutoff in the exponent]\label{rot:lem:R1ext}
Let $\nu$ be a probability measure, $\chi$ an indicator with $\nu(1-\chi)\le\epsilon$, $W$ measurable with $|W|\le M$, and $e^{2M}\epsilon\le\frac12$. Then \textup{(i)} $\int\chi e^Wd\nu\ge e^{-2e^{2M}\epsilon}\int e^Wd\nu$, and \textup{(ii)} $\int\chi e^Wd\nu\le\int e^{\chi W}d\nu\le e^{2e^M\epsilon}\int\chi e^Wd\nu$.
\end{lemma}
\begin{proof}
(i) $\int(1-\chi)e^W\le e^M\epsilon\le e^{2M}\epsilon\int e^W$, because $\int e^W\ge e^{-M}$; and $1-u\ge e^{-2u}$ for $u\le\frac12$. (ii) On $\{\chi=0\}$, $e^{\chi W}=1\ge0=\chi e^W$, so $\int e^{\chi W}=\int\chi e^W+\nu(1-\chi)$; and $\int\chi e^W\ge e^{-M}(1-\epsilon)\ge e^{-M}/2$, so $\nu(1-\chi)\le2\epsilon e^M\int\chi e^W$.
\end{proof}
The lemma is applied with $W=\Psi\chi$, for which $\chi e^W=\chi e^\Psi$.

\begin{lemma}[connected Wick expansion]\label{rot:lem:R1wick}
Let $\zeta$ be a centred Gaussian vector with covariance $G$ and let $F_j=\zeta^{\alpha_j}$, $j=1,\dots,k$, be monomials, each a multiset of legs. Then $\kappa(F_1,\dots,F_k)=\sum_\pi\prod_{\{u,v\}\in\pi}G(u,v)$, the sum over the pairings $\pi$ of all legs that connect the $k$ factors (the graph on $\{1,\dots,k\}$ with an edge for each pair joining legs of two different factors is connected; pairs inside a factor are allowed). For $N(h,G)$ the same holds after the shift $z=h+\zeta$, the constants dropping out of joint cumulants with $k\ge2$.
\end{lemma}
\begin{proof}
Isserlis' formula gives $\E\prod_{j\in B}F_j=\sum\prod G$ over all pairings of the legs of the factors in $B$. Grouping each pairing by the partition of $B$ into the connected components of its factor graph gives $\E\prod_{j\in B}F_j=\sum_\sigma\prod_{b\in\sigma}\mathsf C(b)$, with $\mathsf C(b)$ the sum over the connected pairings of the legs of $b$; the moment--cumulant relation $\E\prod_{j\in B}F_j=\sum_\sigma\prod_{b\in\sigma}\kappa(b)$ determines the $\kappa(b)$ from the moments recursively, so $\kappa=\mathsf C$.
\end{proof}

\begin{lemma}[anchored connected bound]\label{rot:lem:R1conn}
Let $\nu=N(h,G)$ on $(\R^n)^X$ with $\|G(x,y)\|\le g_0e^{-\theta|x-y|}$ and $\|h\|_\infty\le\beta$. Let $F_1,\dots,F_k$ ($k\ge2$) be polynomials of degree at most $D$ in the variables of $X$, sorted by supports as in (R3), with $\|F_j\|_{\theta'}\le a_j$ for some $\theta'\le\theta$, and let $S_1,S_2\subset\mathbb{L}$ be such that every monomial of $F_1$ has support meeting $S_1$ and every monomial of $F_2$ has support inside $S_2$. Then
\[
|\kappa_\nu(F_1,\dots,F_k)|\le|S_1\cap X|\,K_k(\theta')\,(1+\beta)^{kD}\prod_ja_j\,e^{-(\theta'/2)\dist(S_1,S_2)},\qquad K_k(\theta')\defeq(kD)^{kD/2}\max\{1,g_0\}^{kD/2}\,\tilde c(\theta'/2)^{k-1},
\]
and the same bound without the distance factor when no $S_2$ is given (the anchored bound). For $k=1$, $|\E_\nu F_1|\le|S_1\cap X|\,D^{D/2}\max\{1,g_0\}^{D/2}(1+\beta)^Da_1$. For centred $\nu$ the factor $(1+\beta)^{kD}$ is absent.
\end{lemma}
\begin{proof}
(1) Shift $z=h+\zeta$: a monomial $z^\alpha$ with support $Y$ becomes $\sum_{\alpha'\le\alpha}\binom{\alpha}{\alpha'}h^{\alpha-\alpha'}\zeta^{\alpha'}$, whose coefficients have modulus sum at most $\prod(1+|h|)^{\alpha}\le(1+\beta)^D$; the pieces are kept grouped under the original support $Y$ (their sites lie in $Y$), so the shifted $F'_j$ has $\|F'_j\|_{\theta'}\le(1+\beta)^Da_j$, and the constant pieces drop out for $k\ge2$. For $h=0$ there is no shift and no factor. (2) By \Cref{rot:lem:R1wick}, $|\kappa_\nu(F_1,\dots,F_k)|\le\sum_{Y_1,\dots,Y_k}\sum_{\text{monomials}}\prod|\text{coefficient}|\sum_{\text{connected }\pi}\prod\|G\|$; the factors have at most $kD$ legs in all, so for fixed monomials there are at most $(kD)^{kD/2}$ pairings. (3) For a connected $\pi$ fix a spanning tree of its factor graph, rooted at factor $1$, with $k-1$ tree pairs. Sum the root over the sites $a\in S_1\cap X$ and the supports $Y_1\ni a$: at most $|S_1\cap X|\,\|F'_1\|_{\theta'}$ times the rest. Each child factor $j$ is entered through a tree pair $(y,y')$ with $y$ in the parent's support and $y'\in Y_j$, which gives $\sum_{y'}g_0e^{-\theta|y-y'|}\sum_{Y_j\ni y'}|F'_{j,Y_j}|_1$; the pairs not in the tree are bounded by $g_0$ each. (4) Write $|F'_{j,Y}|_1=\big(|F'_{j,Y}|_1e^{\theta'\ell(Y)}\big)e^{-\theta'\ell(Y)}$ and $e^{-\theta|y-y'|}\le e^{-(\theta'/2)|y-y'|}e^{-(\theta'/2)|y-y'|}$. The tree path from factor $1$ to factor $2$ lifts to a path in $\mathbb{L}$ from $a\in S_1$ to a point of $Y_2\subset S_2$, made of tree pairs and of paths inside shortest connecting graphs of the traversed supports $Y_j$, each of length at most $\ell(Y_j)$ (\Cref{rot:lem:R1lat}(iii)); hence $e^{-(\theta'/2)(\sum_j\ell(Y_j)+\sum_{\text{tree}}|y-y'|)}\le e^{-(\theta'/2)\dist(S_1,S_2)}$. The remaining halves give $\tilde c(\theta'/2)$ per tree pair (\Cref{rot:lem:R1lat}(ii)) and $\|F'_j\|_{\theta'}$ per factor. For $k=1$ the same count applies without a tree.
\end{proof}

\begin{lemma}[cumulants with and without the cutoff]\label{rot:lem:R1chi}
Let $\nu=N(h,G)$ on $(\R^n)^X$ with $\|h\|_\infty\le\beta/2$, $G\le\gamma_0^{-1}$ and $\beta\ge2$; put $\chi\defeq\chi^X_\beta$ and $\epsilon\defeq\nu(1-\chi)$. Let $F_1,\dots,F_k$ ($k\le\bar n$) be polynomials in the variables of $X$ of degree at most $D$ with $\|F_j\|_0\le a_j$, and assume the moment condition
\begin{equation}\label{rot:eq:R1MC}
2^q\,|X|\,(2n)^{1/2}\big(n^q(2q-1)!!\big)^{1/2}\gamma_0^{-q/2}\,e^{-\gamma_0\beta^2/(4n)}\le(2\beta)^q\qquad(1\le q\le2\bar nD).
\end{equation}
Then $|\kappa_\nu(\chi F_1,\dots,\chi F_k)-\kappa_\nu(F_1,\dots,F_k)|\le k\cdot k!\cdot4^k\prod_j(a_j|X|)\,(2\beta)^{kD}\,\epsilon^{1/2}$.
\end{lemma}
\begin{proof}
Let $Z\defeq\max\{1,\max_{x\in X}|z_x|\}$ and $\zeta\defeq z-h$. If $\max_x|\zeta_x|\le\beta$ then $Z\le2\beta$, otherwise $Z\le2\max_x|\zeta_x|$; so $\E Z^q\le(2\beta)^q+2^q\sum_x\E\big[|\zeta_x|^q\mathbf 1\{|\zeta_x|>\beta\}\big]\le2(2\beta)^q$ by the Cauchy--Schwarz inequality, $\E|\zeta_x|^{2q}\le n^q(2q-1)!!\gamma_0^{-q}$ (H\"older's inequality over the $n$ components and the Gaussian moments), $\Pr(|\zeta_x|>\beta)\le2ne^{-\gamma_0\beta^2/(2n)}$ and \eqref{rot:eq:R1MC}. Since $|F_j|\le a_j|X|Z^D$, every product of moments obeys $|\E\prod_{j\in B}F_j|\le2\prod_{j\in B}(a_j|X|)(2\beta)^{|B|D}$, and $|\E[(1-\chi)\prod_{j\in B}F_j]|\le\epsilon^{1/2}(\E\prod_{j\in B}F_j^2)^{1/2}\le\sqrt2\,\epsilon^{1/2}\prod_{j\in B}(a_j|X|)(2\beta)^{|B|D}$. In the M\"obius formula ($\chi^{|B|}=\chi$ on each block) replace the blocks one at a time: each partition contributes at most $k\cdot\sqrt2\cdot2^{k-1}\epsilon^{1/2}\prod_j(a_j|X|)(2\beta)^{kD}$, and $\sum_\pi(|\pi|-1)!\le k!\,2^{k-1}$.
\end{proof}
The moment condition \eqref{rot:eq:R1MC} is monotone: once it holds for a pair $(|X|,\beta)$ it holds for smaller $|X|$ and larger $\beta$.

\begin{lemma}[comparing two Gaussian laws]\label{rot:lem:R1repl}
Let $\nu,\tilde\nu$ be Gaussian laws of $(z_x)_{x\in S}$, $S$ finite, with means $h,\tilde h$ of sup norm at most $\beta/2$ and covariances $G,\tilde G$ all of whose blocks have norm at most $\gamma_0^{-1}$, and let $\sup_x|h(x)-\tilde h(x)|\le\eta_1$, $\sup_{x,y}\|(G-\tilde G)(x,y)\|\le\eta_2$. For polynomials $F_1,\dots,F_k$ in the variables of $S$ of degree at most $D$ with $\|F_j\|_0\le a_j$,
\[
|\kappa_\nu(F_1,\dots,F_k)-\kappa_{\tilde\nu}(F_1,\dots,F_k)|\le k\cdot k!\cdot4^k\,(kD)^{kD+1}\,(1+\beta+\gamma_0^{-1})^{kD}\prod_j(a_j|S|)\,\max\{\eta_1,\eta_2\}.
\]
\end{lemma}
\begin{proof}
A moment of $L\le kD$ coordinate variables is, by Isserlis' formula for a Gaussian with mean, a sum over at most $L^L$ partial pairings of products of at most $L$ factors (covariance entries or means), each of modulus at most $1+\beta+\gamma_0^{-1}$ under both laws; replacing the factors one at a time, two corresponding products differ by at most $L\max\{\eta_1,\eta_2\}(1+\beta+\gamma_0^{-1})^L$. The coefficient sums contribute $\prod_j(a_j|S|)$, and the M\"obius formula with the telescoping of \Cref{rot:lem:R1chi} gives the bound.
\end{proof}

% r03_r1_geo: the cumulant bound for the step Gaussian, part III: the four shifted pavements on Z^3, on the cubic tori and on the skew tori, and the layers of a cell.

\begin{lemma}[pavements]\label{rot:lem:R1geo}
Let $b\ge b_{\mathrm{geo}}\defeq576$, $q_b\defeq\lceil b^2\rceil$, $w_b\defeq\lceil b^{3/2}\rceil$ and $s_b\defeq\lfloor q_b/4\rfloor$, so that $s_b\ge5w_b+2$, $q_b\ge6w_b$ and $3q_b\le4b^2$; put $v_b\defeq(4b^2)^3=64\,b^6$. On each $\mathbb{L}$ of the setting there are four partitions $\mathcal P^i$, $i=0,1,2,3$, of $\mathbb{L}$ into \emph{cells} $\Box$, with \emph{levels} $\lambda_i:\mathbb{L}\to[0,\infty]$, such that
\begin{enumerate}[label=\textup{(G\arabic*)},leftmargin=*]
\item $\lambda_i$ is $1$-Lipschitz;
\item if $x,y$ lie in different cells of $\mathcal P^i$, then $|x-y|\ge\lambda_i(x)+\lambda_i(y)$;
\item $\bigcap_{i=0}^3\{x:\ \lambda_i(x)<2.5w_b+1\}=\emptyset$;
\item every cell has at most $v_b$ sites; in particular, if some $\mathcal P^i$ consists of a single cell, then $|\mathbb{L}|\le v_b$.
\end{enumerate}
\end{lemma}
\begin{proof}
\emph{Numbers.} For $b\ge576$: $s_b\ge b^2/4-1\ge5b^{3/2}+7\ge5w_b+2$ (the middle inequality is $\sqrt b\ge20+32b^{-3/2}$, true for $\sqrt b\ge24$), $q_b\ge b^2\ge6b^{3/2}+6\ge6w_b$, and $3q_b\le3b^2+3\le4b^2$.

\emph{Coordinates.} Each $\mathbb{L}$ has three coordinate cycles $\mathbb I_a$, $a=1,2,3$, with a well-defined map $x\mapsto x_a\in\mathbb I_a$ and $N_a\defeq|\mathbb I_a|\in\N\cup\{\infty\}$: for $\Z^3$, $\mathbb I_a=\Z$; for $(\Z/N\Z)^3$, $\mathbb I_a=\Z/N\Z$; for $\Z^3/mG_\pm$, $\mathbb I_a\defeq\Z/N_a\Z$ with $N_a\Z\defeq\pi_a(mG_\pm)$ the projection of the period lattice on the coordinate $a$, that is $N_1=N_2=10m$ and $N_3=50m$ (the entries $40,\mp30,0$ and $\pm30,40,0$ have greatest common divisor $10$); the map is well defined because the period lattice $G\defeq mG_\pm$ is contained in $G'\defeq\bigoplus_aN_a\Z$. In every case $\dist_{\mathbb I_a}(x_a,y_a)\le|\tilde x_a-\tilde y_a|\le|\tilde x-\tilde y|=|x-y|$ for lifts realising the quotient distance. For the skew tori we record: $\det G=(40\cdot40+30\cdot30)\cdot50\cdot m^3=125000\,m^3=(50m)^3$, $[G':G]=125000/5000=25$, and $G\cap\{g_3=0\}=m\langle(40,\mp30,0),(\pm30,40,0)\rangle$ has index $2500m^2$ in $\Z^2\times\{0\}$.

\emph{Cuts.} A coordinate $a$ is \emph{cut} if $N_a\ge q_b$ (always for $\Z^3$). For a cut coordinate choose a set $B^0_a\subset\mathbb I_a+\frac12$ of cuts all of whose (cyclic) gaps lie in $[q_b,2q_b)$: for $\mathbb I_a=\Z$ take $q_b\Z+\frac12$; on a cycle take $p\defeq\lfloor N_a/q_b\rfloor\ge1$ cuts with gaps $\lfloor N_a/p\rfloor$ or $\lceil N_a/p\rceil$, which lie in $[q_b,2q_b-1]$. Put $B^i_a\defeq B^0_a+is_b$, $i=0,\dots,3$. Cuts of $B^i_a$ and of $B^{i'}_a$, $i\ne i'$, are at cyclic distance at least $s_b$: two shifts of the same cut differ by $|i-i'|s_b\in[s_b,3s_b]$, and $N_a-3s_b\ge q_b-3s_b\ge s_b$; shifts of different cuts differ by at least $q_b-3s_b\ge s_b$. For $x\in\mathbb{L}$ put $\lambda_i(x)\defeq\min_{a\ \mathrm{cut}}\dist_{\mathbb I_a}(x_a,B^i_a)$, and $\lambda_i\defeq+\infty$ if no coordinate is cut. The \emph{cut set} $C^i$ is the set of points $\xi$ of the continuum torus ($\R^3$, $\R^3/N\Z^3$ or $\R^3/G$) with $\xi_a\in B^i_a$ for some cut coordinate $a$, and the cells of $\mathcal P^i$ are the intersections of $\mathbb{L}$ with the connected components of the complement of $C^i$; no site lies on a cut, so the cells partition $\mathbb{L}$.

\emph{(G1)} is the Lipschitz property of $\dist_{\mathbb I_a}(\cdot,B^i_a)$ together with $\dist_{\mathbb I_a}(x_a,y_a)\le|x-y|$. \emph{(G3)}: if $\lambda_i(x)<2.5w_b+1$ for all four $i$, each minimum is attained at a cut coordinate, and since there are at most three, two indices $i\ne i'$ use the same coordinate $a$; then a cut of $B^i_a$ and a cut of $B^{i'}_a$ are at cyclic distance less than $2(2.5w_b+1)=5w_b+2\le s_b$, a contradiction. If no coordinate is cut the four sets are empty. \emph{(G2)}: lift $x,y$ to $\tilde x,\tilde y\in\R^3$ with $|\tilde x-\tilde y|=|x-y|$. The lift $\tilde C^i$ of the cut set is the union of the planes $\{\xi_a=c\}$ over the cut coordinates $a$ and the lifts $c$ of the cuts of $B^i_a$; it is $G'$-periodic, hence $G$-periodic, the components of $\R^3\setminus\tilde C^i$ are the products of open gap intervals in the cut coordinates and of $\R$ in the others, and the projection maps them onto the components of the complement of $C^i$. If the segment $[\tilde x,\tilde y]$ met none of the planes it would lie in one component, and $x,y$ would lie in one cell; so it crosses a plane $\{\xi_a=c\}$ with $c$ between $\tilde x_a$ and $\tilde y_a$, and $|x-y|\ge|\tilde x_a-\tilde y_a|=|\tilde x_a-c|+|c-\tilde y_a|\ge\dist_{\mathbb I_a}(x_a,B^i_a)+\dist_{\mathbb I_a}(y_a,B^i_a)\ge\lambda_i(x)+\lambda_i(y)$. \emph{(G4)}: a component $Q$ of $\R^3\setminus\tilde C^i$ projects onto its cell injectively modulo $G_Q\defeq\{g\in G:\ Q+g=Q\}=\{g\in G:\ g_a=0\text{ for every cut }a\}$: if $Q$ and $Q+g$ meet they coincide, $\tilde C^i$ being $G$-invariant. On $\Z^3$, and whenever all three coordinates are cut, $G_Q=0$ and the cell is a box with sides at most $2q_b-1$, hence has at most $(2q_b)^3$ sites. On a cubic torus the only other case is that no coordinate is cut, when $|\mathbb{L}|=N^3<q_b^3$. On $\Z^3/mG_\pm$, where $N_3=5N_1$, two further cases occur: only the third coordinate is cut ($10m<q_b\le50m$), the cell is the product of $\Z^2/(G\cap\{g_3=0\})$ with an interval of fewer than $2q_b$ integers and has fewer than $2500m^2\cdot2q_b=25\,(10m)^2\cdot2q_b<50\,q_b^3$ sites; or no coordinate is cut ($50m<q_b$) and $|\mathbb{L}|=\det G=(50m)^3<q_b^3$. With $q_b\le b^2+1\le\frac{257}{256}b^2$ the three counts are below $8.1\,b^6$, $50.6\,b^6$ and $1.02\,b^6$, hence below $v_b=64\,b^6$.
\end{proof}

\paragraph{Layers.} For a cell $\Box$ of $\mathcal P^i$ write $\lambda\defeq\lambda_i$ and define, all intersected with $\Box\cap\Lambda$:
\[
\Gamma_1(\Box)\defeq\{\lambda<w_b\},\quad\Gamma_2(\Box)\defeq\{w_b\le\lambda<2w_b\},\quad\Gamma_3(\Box)\defeq\{1.5w_b\le\lambda<2w_b\},\quad\Box'\defeq\{\lambda\ge2w_b\},\quad\Gamma_4(\Box)\defeq\{2w_b\le\lambda<2.5w_b\},
\]
$\mathrm{Deep}(\Box)\defeq\{\lambda\ge2.5w_b\}$ and $D_\Box\defeq\Gamma_2(\Box)\cup\Box'=\{\lambda\ge w_b\}$. By (G1) and (G2), two points of one cell at levels $\lambda(x)\ge\lambda(y)$ are at distance at least $\lambda(x)-\lambda(y)$, and two points of different cells at distance at least $\lambda(x)+\lambda(y)$. Hence: the sets $D_\Box$ of distinct cells are at distance at least $2w_b$; $\dist(D_\Box,\Gamma_1(\tilde\Box))\ge w_b$ and $\dist(\mathrm{Deep}(\Box),\tilde\Box\cap\{\lambda\ge1.5w_b\})\ge4w_b$ for $\tilde\Box\ne\Box$; and $\dist(\mathrm{Deep}(\Box),\Lambda\setminus\Box')\ge\frac12w_b$, $\dist(\Box\cap\{\lambda\ge1.5w_b\},\Lambda\setminus D_\Box)\ge\frac12w_b$, $\dist(\mathrm{Deep}(\Box),\Lambda\setminus(\Box\cap\{\lambda\ge1.5w_b\}))\ge w_b$, $\dist(\Gamma_4(\Box),\Box\cap\{\lambda<1.5w_b\})\ge\frac12w_b$, $\dist(\Box',\Gamma_2(\Box)\setminus\Gamma_3(\Box))\ge\frac12w_b$, $\dist(\Gamma_1(\Box),\Box')\ge w_b$.

% r03_r1_stage: the cumulant bound for the step Gaussian, part IV: one pavement stage (the sorting of the monomials, the stage lemma and the bookkeeping lemma).

\paragraph{One stage.} Fix $i$, the pavement $\mathcal P^i$ with its layers, a set $J'\subset J$ (write $V'\defeq V_{J'}$) and a threshold $\beta\in[b_{\min},b_{\max}]$, where $b_{\min}\defeq\gamma_K^4b$ and $b_{\max}\defeq\gamma_K^{-4}b$. Put $\Phi(\beta;R)\defeq\int\chi_\beta e^{V_R}d\mu_P$, $\chi^\Box_\beta\defeq\chi^{D_\Box}_\beta$, $\mathcal B'\defeq\{\Box\in\mathcal P^i:\ \Box\cap J'\ne\emptyset\}$ (the \emph{$J$-cells}), $N'\defeq|\mathcal B'|\le|J'|$, $D\defeq\bigcup_{\Box\in\mathcal B'}D_\Box$, $\Gamma\defeq\Lambda\setminus D$ and $J''\defeq J'\cap\{\lambda_i<2.5w_b\}$. The monomials of $V'$, all of which have support in $J'$, are sorted as follows.
\begin{itemize}[leftmargin=*]
\item $V_\Gamma$: support in $\Gamma$. For $\Box\in\mathcal B'$, $\Psi_\Box$: support in $D_\Box$, or support in $\Gamma_1(\Box)\cup\Gamma_2(\Box)$ meeting both. $R_0\defeq V'-V_\Gamma-\sum_\Box\Psi_\Box$.
\item $\Psi_1$: support in $\Box\cap\{\lambda\ge1.5w_b\}$ and meeting $\Box'$. $\Psi_2$: support in $\Gamma_1(\Box)\cup\Gamma_2(\Box)$ and meeting $\Gamma_2(\Box)$. $\Psi_3\defeq\Psi_\Box-\Psi_1-\Psi_2$: the monomials of $\Psi_\Box$ meeting $\Box'$ and $\Gamma_2(\Box)\setminus\Gamma_3(\Box)$.
\item $\Psi_1'$: the monomials of $\Psi_1$ with support in $\Gamma_3(\Box)\cup\Gamma_4(\Box)$ (they meet $\Gamma_4(\Box)$). $\Psi_1''\defeq\Psi_1-\Psi_1'$: the monomials of $\Psi_1$ meeting $\mathrm{Deep}(\Box)$. $W_\Box\defeq\Psi_1'+\Psi_2$ and $R_1\defeq V_{J''}-V_\Gamma-\sum_\Box W_\Box$.
\end{itemize}
$\Psi_2$ depends on $z_{\Gamma_1(\Box)}$ as well as on $z_{D_\Box}$. Under the conditional laws below it is read as a polynomial in $z_{D_\Box}$ whose coefficients depend on $z_{\Gamma_1(\Box)}$: every such monomial keeps at least one variable in $D_\Box$, its support shrinks to $Y\cap D_\Box$ (so its length does not increase, \Cref{rot:lem:R1lat}(iii)), and on $\{\|z_\Gamma\|_\infty\le c_K\beta\}$ its coefficient sums grow by at most the factor $(c_K\beta)^D$. The other pieces involve $z_{D_\Box}$ only.

\emph{Length facts.} $V_\Gamma$, $\Psi_1'$ and $\Psi_2$ consist of monomials of $V_{J''}$ (their supports lie in $J'\cap\{\lambda<w_b\}$, in $\Gamma_3\cup\Gamma_4$ and in $\{\lambda<2w_b\}$ respectively; a site of $J'$ in a cell $\Box$ with $\lambda\ge w_b$ lies in $D_\Box$, since then $\Box\in\mathcal B'$), and the three classes are disjoint, so $R_1$ is a sub-sum of the monomials of $V_{J''}$. Every monomial of $R_0$ meets some $D_\Box$, $\Box\in\mathcal B'$, and has $\ell\ge w_b$: if it lies in $\Box$ it meets $\Gamma_1(\Box)$ (its support is not inside $D_\Box$) and $\Box'$ (its support is not inside $\Gamma_1\cup\Gamma_2$); otherwise it meets a second cell, in $D_{\tilde\Box}$ or in $\Gamma_1(\tilde\Box)$. Every monomial of $\Psi_3$ has $\ell\ge\frac12w_b$. Every monomial of $R_1$ meets $\Gamma_2(\Box)\cup\Gamma_4(\Box)$ for some $\Box\in\mathcal B'$ and has $\ell\ge\frac12w_b$: it meets some $D_\Box$ inside $\{\lambda<2.5w_b\}$; if it lies in $\Box$ and meets $\Gamma_4(\Box)$, it is not inside $\Gamma_3\cup\Gamma_4$ (it is not in $\Psi_1'$), so it meets $\{\lambda<1.5w_b\}$; if it lies in $\Box$ and misses $\Gamma_4(\Box)$, its support is inside $\Gamma_1\cup\Gamma_2$ and meets $\Gamma_2$, so it belongs to $\Psi_2$, which is excluded; if it meets a second cell, $\ell\ge w_b$. On $\{\|z\|_\infty\le c_K\beta\}$, with $\beta\le b_{\max}$ and $|D_\Box|\le v_b$, \Cref{rot:lem:R1lat}(iv) gives
\[
|R_0|,\ |R_1|\le N'v_b\mathsf A(c_Kb_{\max})^De^{-\varsigma w_b/2},\qquad|\Psi_\Box|,\ |W_\Box|\le M\defeq\mathsf Av_b(c_Kb_{\max})^D,\qquad|\Psi_3|\le Me^{-\varsigma w_b/2},
\]
and $M=\mathsf m_*\,\mathsf Ab^{D+6}$ with $\mathsf m_*\defeq64\,c_K^D\gamma_K^{-4D}$, so $M\le\mathsf m_*$ under $\mathsf Ab^{D+6}\le1$.

\paragraph{Thresholds.} Put $\epsilon_\beta\defeq2nv_be^{-\gamma_0\beta^2/(8n)}$, with $b_{\min}$, $b_{\max}$ and $\mathsf m_*$ as above. The threshold of \Cref{rot:thm:R1} is $b_*\defeq\max\{576,\,b_{\mathrm{dec}},\,b_{\mathrm{tail}},\,b_{\mathrm{mc}},\,b_{\mathrm{gci}}\}$, where $b_{\mathrm{dec}}$, $b_{\mathrm{tail}}$, $b_{\mathrm{mc}}$ and $b_{\mathrm{gci}}$ are the least numbers beyond which, respectively, $c_Pc_de^{-\delta_Pb^{3/2}}\le\gamma_0/2$; $2nv_be^{2\mathsf m_*}e^{-\gamma_0\gamma_K^8b^2/(8n)}\le\frac12$ and $\gamma_K^4b\ge2$; the moment condition \eqref{rot:eq:R1MC} holds with $|X|=v_b$ and $\beta=b_{\min}$; and $2e^{-\gamma_0\gamma_K^8b^2/(2n)}\le\frac12$. Each of these inequalities, once true, stays true for larger $b$, and all of them depend only on $(n,\bar n,D,\gamma_0,c_P,\delta_P,c_K)$.

\begin{lemma}[one stage]\label{rot:lem:R1stg}
Let $b\ge b_*$ and $\beta\in[b_{\min},b_{\max}]$, with $\beta/\gamma_K\le b_{\max}$ in \textup{(U)}. Then
\begin{align*}
\textup{(L)}\qquad&\Phi(\beta;J')\ge\Phi(\gamma_K\beta;J'')\,\exp\Big\{\sum_{k=1}^{\bar n}\frac1{k!}\big[\kappa_{\mu_P}(V_{J'};k)-\kappa_{\mu_P}(V_{J''};k)\big]-N'E_*\Big\},\\
\textup{(U)}\qquad&\Phi(\beta;J')\le\Phi(\beta/\gamma_K;J'')\,\exp\Big\{\sum_{k=1}^{\bar n}\frac1{k!}\big[\kappa_{\mu_P}(V_{J'};k)-\kappa_{\mu_P}(V_{J''};k)\big]+N'E_*\Big\},
\end{align*}
with the per-cell error $E_*=E_*(b,\mathsf A)\defeq e_R+e_{\mathrm{dec}}+e_{\mathrm{ext}}+e_{\mathrm{cum}}+e_3+e_{\mathrm{mix}}+e_{\mathrm{rep}}+e_{\mathrm{bk}}$ defined in the proof and in \Cref{rot:lem:R1bk}. All its terms except $e_{\mathrm{cum}}=2(2M)^{\bar n+1}$ are increasing in $\mathsf A$ and carry one of the factors $e^{-\varsigma w_b/2}$, $e^{-\delta_Pw_b}$, $e^{-\bar\mu w_b/4}$, $e^{-\mu_1w_b/8}$, $\epsilon(v_b,b_{\min})^{1/2}$.
\end{lemma}
\begin{proof}
Recall $\epsilon_\beta=\epsilon(v_b,\beta)=2nv_be^{-\gamma_0\beta^2/(8n)}$, which decreases in $\beta$.

\emph{Lower bound.} (L1) $V'=V_\Gamma+\sum_\Box\Psi_\Box+R_0$, so on $\supp\chi_\beta$, $e^{V'}\ge e^{-|R_0|_{\max}}e^{V_\Gamma+\sum_\Box\Psi_\Box}$. (L2) Shrink the cutoff on $\Gamma$: $\chi_\beta\ge\chi^\Gamma_{\gamma_K\beta}\chi^D_\beta$, as $\gamma_K\beta K\subset\beta K$. Disintegrate $\mu_P$ into its marginal on $z_\Gamma$ and the conditional law of $z_D$. The $z_D$-integrand $\prod_\Box\chi^\Box_\beta e^{\Psi_\Box}$ is non-negative and supported in $\{\|z_D\|_\infty\le c_K\beta\}$, and $\|z_\Gamma\|_\infty\le c_K\gamma_K\beta\le c_K\beta$ on $\supp\chi^\Gamma_{\gamma_K\beta}$; \Cref{rot:lem:R1dec} with the pieces $D_\Box$, $\Box\in\mathcal B'$ (separation $2w_b$, so $\omega=c_Pc_de^{-\delta_Pw_b}\le\gamma_0/2$ because $b\ge b_*$), gives
\[
\int\mu_P(dz_D\,|\,z_\Gamma)\prod_\Box\chi^\Box_\beta e^{\Psi_\Box}\ge e^{-\omega_{\mathrm{tot}}}\prod_{\Box\in\mathcal B'}\int\nu_\Box(dz_{D_\Box})\,\chi^\Box_\beta e^{\Psi_\Box},\qquad\nu_\Box\defeq N(h_\Box,G_{D_\Box}),\quad h_\Box\defeq-G_{D_\Box}P_{D_\Box\Gamma}z_\Gamma,
\]
with $\omega_{\mathrm{tot}}\le N'v_bc_Pc_de^{-\delta_Pw_b}\big[(c_K\beta)^2/2+n/\gamma_0+(c_P\tilde c(\delta_P)c_K\beta)^2/(2\gamma_0^2)\big]$; only the $J$-cells are decoupled, which is why the error is proportional to $N'$. (L3) By \Cref{rot:lem:R1sch}, $\|h_\Box\|_\infty\le\beta/2$ and $G_{D_\Box}\le\gamma_0^{-1}$; by \Cref{rot:lem:R1tail}, $\nu_\Box(1-\chi^\Box_\beta)\le\epsilon_\beta$; by \Cref{rot:lem:R1ext}(i) with $W\defeq\Psi_\Box\chi^\Box_\beta$, $|W|\le M$ (and $e^{2M}\epsilon_\beta\le\frac12$ because $b\ge b_*$), $\int\nu_\Box\chi^\Box_\beta e^{\Psi_\Box}\ge e^{-2e^{2M}\epsilon_\beta}\int\nu_\Box e^{W}$. (L4) By \Cref{rot:lem:R1cum}, $\log\int\nu_\Box e^{W}\ge\sum_{k\le\bar n}\kappa_{\nu_\Box}(W;k)/k!-(2M)^{\bar n+1}$. (L5) Expand $W=\Psi_1'\chi+\Psi_1''\chi+\Psi_2\chi+\Psi_3\chi$ ($\chi=\chi^\Box_\beta$) multilinearly: $\kappa_{\nu_\Box}(W;k)$ is the sum over the at most $4^k$ words of length $k$ in these four letters of their joint cumulants. Words containing $\Psi_3\chi$: by \Cref{rot:lem:R1cum} each is at most $k!\,2^{k-1}M^ke^{-\varsigma w_b/2}$, so they contribute at most $e_3\defeq\sum_{k\le\bar n}(8M)^ke^{-\varsigma w_b/2}$. (L6) Words with at least one $\Psi_1''\chi$ and one $\Psi_2\chi$ (so $k\ge2$): \Cref{rot:lem:R1chi} removes the cutoffs (with $X=D_\Box$, $|X|\le v_b$, $a_j\le\mathsf A(c_K\beta)^D$, and the moment condition holding because $b\ge b_*$), then \Cref{rot:lem:R1conn} with $S_1\defeq\mathrm{Deep}(\Box)$ (met by every monomial of $\Psi_1''$), $S_2\defeq\Box\cap\{\lambda<2w_b\}$ (containing the supports of $\Psi_2$ read in $z_{D_\Box}$), $\dist(S_1,S_2)\ge\frac12w_b$, $\theta=\mu_1$ (\Cref{rot:lem:R1ct}) and $\theta'=\bar\mu$ ($\|\cdot\|_{\bar\mu}\le\|\cdot\|_\varsigma$) bounds each by $v_bK_k(\bar\mu)(1+\beta)^{kD}(\mathsf A(c_K\beta)^D)^ke^{-\bar\mu w_b/4}$; in total these words contribute at most
\[
e_{\mathrm{mix}}\defeq\sum_{k=2}^{\bar n}\frac{4^k}{k!}\Big[k\,k!\,4^k\big(\mathsf A(c_Kb_{\max})^Dv_b\big)^k(2b_{\max})^{kD}\epsilon_{b_{\min}}^{1/2}+v_bK_k(\bar\mu)(1+b_{\max})^{kD}\big(\mathsf A(c_Kb_{\max})^D\big)^ke^{-\bar\mu w_b/4}\Big].
\]
(L7) Words with at least one $\Psi_1''\chi$ and neither $\Psi_2\chi$ nor $\Psi_3\chi$ (all factors $\Psi_1'\chi$ or $\Psi_1''\chi$, with supports in $S\defeq\Box\cap\{\lambda\ge1.5w_b\}$): \Cref{rot:lem:R1chi} removes the cutoffs, and \Cref{rot:lem:R1repl} compares the marginal of $\nu_\Box$ on $S$ (mean $h_\Box|_S$, covariance $(G_{D_\Box})_{SS}$) with that of $\mu_P$ (centred, covariance $C_{SS}$): since $\dist(S,\Lambda\setminus D_\Box)\ge\frac12w_b$, \Cref{rot:lem:R1sch} gives $\eta_1\le\mathsf M_1c_K\gamma_K\beta e^{-\mu_1w_b/4}$ and $\eta_2\le c_{\mathrm D}e^{-\mu_1w_b/8}$. This replaces these words by the $z_\Gamma$-independent quantity
\[
c_\Box\defeq\sum_{k=1}^{\bar n}\frac1{k!}\big[\kappa_{\mu_P}(\Psi_1'+\Psi_1'';k)-\kappa_{\mu_P}(\Psi_1';k)\big]
\]
with an error at most
\[
e_{\mathrm{rep}}\defeq\sum_{k=1}^{\bar n}\frac{2^k}{k!}\,k\,k!\,4^k\big(\mathsf A(c_Kb_{\max})^Dv_b\big)^k\Big[(2b_{\max})^{kD}\epsilon_{b_{\min}}^{1/2}+(kD)^{kD+1}(1+b_{\max}+\gamma_0^{-1})^{kD}\big(\mathsf M_1c_Kb_{\max}+c_{\mathrm D}\big)e^{-\mu_1w_b/8}\Big].
\]
What is left of $\sum_k\kappa_{\nu_\Box}(W;k)/k!$ are the words in $\Psi_1'\chi$ and $\Psi_2\chi$ alone, that is $\sum_k\kappa_{\nu_\Box}(W_\Box\chi;k)/k!$. (L8) By \Cref{rot:lem:R1cum} backwards and $e^{\chi W_\Box}\ge\chi e^{W_\Box}$, $\sum_k\kappa_{\nu_\Box}(W_\Box\chi;k)/k!\ge\log\int\nu_\Box\chi^\Box_\beta e^{W_\Box}-(2M)^{\bar n+1}$. (L9) $\prod_\Box\int\nu_\Box\chi^\Box_\beta e^{W_\Box}=\int\nu_{\mathrm{loc}}\prod_\Box\chi^\Box_\beta e^{W_\Box}\ge e^{-\omega_{\mathrm{tot}}}\int\mu_P(dz_D\,|\,z_\Gamma)\prod_\Box\chi^\Box_\beta e^{W_\Box}$ by \Cref{rot:lem:R1dec} backwards (the integrand is non-negative and supported in the required set). The $c_\Box$ do not depend on $z_\Gamma$, so they come out of the $z_\Gamma$ integral; with $\chi^D_\beta\ge\chi^D_{\gamma_K\beta}$ and $V_{J''}=V_\Gamma+\sum_\Box W_\Box+R_1$,
\[
\int d\mu_P\,\chi^\Gamma_{\gamma_K\beta}\chi^D_\beta\,e^{V_\Gamma+\sum_\Box W_\Box}\ge e^{-|R_1|_{\max}}\int d\mu_P\,\chi_{\gamma_K\beta}\,e^{V_{J''}}=e^{-|R_1|_{\max}}\,\Phi(\gamma_K\beta;J'').
\]
(L10) By \Cref{rot:lem:R1bk}, $\sum_{\Box\in\mathcal B'}c_\Box\ge\sum_{k\le\bar n}\frac1{k!}[\kappa_{\mu_P}(V_{J'};k)-\kappa_{\mu_P}(V_{J''};k)]-N'e_{\mathrm{bk}}$. Collecting the errors, with $e_R\defeq2v_b\mathsf A(c_Kb_{\max})^De^{-\varsigma w_b/2}$ (for $R_0$ and $R_1$), $e_{\mathrm{dec}}\defeq2v_bc_Pc_de^{-\delta_Pw_b}\big[(c_Kb_{\max})^2/2+n/\gamma_0+(c_P\tilde c(\delta_P)c_Kb_{\max})^2/(2\gamma_0^2)\big]$ (two uses of \Cref{rot:lem:R1dec}), $e_{\mathrm{ext}}\defeq4e^{2M}\epsilon_{b_{\min}}$ and $e_{\mathrm{cum}}\defeq2(2M)^{\bar n+1}$, proves (L).

\emph{Upper bound.} The same steps with the following changes. (U2) Enlarge the cutoff on $D$: $\chi_\beta\le\chi^\Gamma_\beta\chi^D_{\beta'}$ with $\beta'\defeq\beta/\gamma_K$; the $z_D$-integrand is non-negative and supported in $\{\|z_D\|_\infty\le c_K\beta'\}$, and $\|z_\Gamma\|_\infty\le c_K\beta=c_K\gamma_K\beta'$, so \Cref{rot:lem:R1sch} gives $\|h_\Box\|_\infty\le\beta'/2$ and the upper half of \Cref{rot:lem:R1dec} applies with the radius $c_K\beta'$. (U3) $\int\nu_\Box\chi^\Box_{\beta'}e^{\Psi_\Box}\le\int\nu_\Box e^{\Psi_\Box\chi^\Box_{\beta'}}$. (U4)--(U7) as (L4)--(L7) with $\beta'$ in place of $\beta$, using the upper half of \Cref{rot:lem:R1cum}. (U8) $\sum_k\kappa_{\nu_\Box}(W_\Box\chi;k)/k!\le\log\int\nu_\Box e^{W_\Box\chi}+(2M)^{\bar n+1}\le\log\int\nu_\Box\chi^\Box_{\beta'}e^{W_\Box}+2e^{M}\epsilon_{\beta'}+(2M)^{\bar n+1}$ by \Cref{rot:lem:R1ext}(ii). (U9) The upper half of \Cref{rot:lem:R1dec} backwards, then $\chi^\Gamma_\beta\chi^D_{\beta'}\le\chi_{\beta'}$ and $e^{V_\Gamma+\sum_\Box W_\Box}\le e^{|R_1|_{\max}}e^{V_{J''}}$. (U10) The upper half of \Cref{rot:lem:R1bk}. Since all thresholds used lie in $[b_{\min},b_{\max}]$, every error is again bounded by $E_*$.
\end{proof}

\begin{lemma}[bookkeeping]\label{rot:lem:R1bk}
$\big|\sum_{\Box\in\mathcal B'}c_\Box-\sum_{k=1}^{\bar n}\frac1{k!}\big[\kappa_{\mu_P}(V_{J'};k)-\kappa_{\mu_P}(V_{J''};k)\big]\big|\le N'e_{\mathrm{bk}}$ with $e_{\mathrm{bk}}\defeq\sum_{k=1}^{\bar n}\frac{6^k}{k!}\,v_bK_k(\bar\mu/2)\,\mathsf A^ke^{-\bar\mu w_b/4}$.
\end{lemma}
\begin{proof}
Here the law is the centred $\mu_P$, so \Cref{rot:lem:R1conn} applies with the covariance decay $g_0e^{-\mu_1|x-y|}$ of \Cref{rot:lem:R1ct} and without the factor $(1+\beta)^{kD}$. (a) $V_{J'}=V_{J''}+\sum_{\Box\in\mathcal B'}X_\Box+R_2$ with $X_\Box\defeq\Psi_1''$ of the cell $\Box$: a monomial of $V_{J'}$ not inside $J''$ has a site of $J'$ at level at least $2.5w_b$, so it meets $\mathrm{Deep}(\Box)$ for a $J$-cell $\Box$; if its support lies in $\Box\cap\{\lambda\ge1.5w_b\}$ it belongs to $\Psi_1$ and not to $\Psi_1'$, that is to $X_\Box$; otherwise it has $\ell\ge\dist(\mathrm{Deep}(\Box),\Lambda\setminus(\Box\cap\{\lambda\ge1.5w_b\}))\ge w_b$ and is put into $R_2$. (b) Expand $\kappa_{\mu_P}(V_{J''}+X+R_2;k)$, $X\defeq\sum_\Box X_\Box$, into words. Words with at least one $R_2$: \Cref{rot:lem:R1conn} anchored at $S_1\defeq\bigcup_{\Box\in\mathcal B'}\mathrm{Deep}(\Box)$ (met by every monomial of $R_2$; $|S_1\cap\Lambda|\le N'v_b$) with $\theta'=\bar\mu/2$ and $\|R_2\|_{\bar\mu/2}\le\mathsf Ae^{-\bar\mu w_b/2}$ (\Cref{rot:lem:R1lat}(iv)), the other factors having $\|\cdot\|_{\bar\mu/2}\le\mathsf A$: at most $N'v_bK_k(\bar\mu/2)\mathsf A^ke^{-\bar\mu w_b/2}$ each. Words with some $X$ and no $R_2$: sort them by the cell $\Box$ of the first $X$ factor and write every other $X$ factor as $X_\Box+X_{\ne\Box}$. Words containing an $X_{\ne\Box}$: \Cref{rot:lem:R1conn} with $S_1=\mathrm{Deep}(\Box)$, $S_2\defeq\bigcup_{\tilde\Box\ne\Box}(\tilde\Box\cap\{\lambda\ge1.5w_b\})$, $\dist(S_1,S_2)\ge4w_b$ and $\theta'=\bar\mu$: at most $v_bK_k(\bar\mu)\mathsf A^ke^{-2\bar\mu w_b}$ each. Words with only $X_\Box$ and $V_{J''}$ factors: for this cell write $V_{J''}=\Psi_1'+Y^{\mathrm{long}}+Y^{\mathrm{far}}$, where $Y^{\mathrm{long}}$ consists of the monomials of $V_{J''}$ meeting $\Gamma_4(\Box)$ that are not in $\Psi_1'$ (they are not inside $\Gamma_3\cup\Gamma_4$, so they meet $\{\lambda<1.5w_b\}$ or another cell, and have $\ell\ge\frac12w_b$) and $Y^{\mathrm{far}}$ of the rest (supports in $J''\setminus\Gamma_4(\Box)$, at distance at least $\frac12w_b$ from $\mathrm{Deep}(\Box)$). The words in $X_\Box$ and $\Psi_1'$ alone with at least one $X_\Box$ add up to exactly $\kappa_{\mu_P}(\Psi_1'+\Psi_1'';k)-\kappa_{\mu_P}(\Psi_1';k)$, which is $k!$ times the $k$-th term of $c_\Box$. Words with a $Y^{\mathrm{far}}$: \Cref{rot:lem:R1conn} with $S_1=\mathrm{Deep}(\Box)$, $S_2=J''\setminus\Gamma_4(\Box)$, $\theta'=\bar\mu$: at most $v_bK_k(\bar\mu)\mathsf A^ke^{-\bar\mu w_b/4}$ each. Words with a $Y^{\mathrm{long}}$ and no $Y^{\mathrm{far}}$: anchored at $\mathrm{Deep}(\Box)$ with $\theta'=\bar\mu/2$ and $\|Y^{\mathrm{long}}\|_{\bar\mu/2}\le\mathsf Ae^{-\bar\mu w_b/4}$: at most $v_bK_k(\bar\mu/2)\mathsf A^ke^{-\bar\mu w_b/4}$ each. Every slot of a word takes one of six roles ($\Psi_1'$, $Y^{\mathrm{far}}$, $Y^{\mathrm{long}}$, $X_\Box$, $X_{\ne\Box}$, $R_2$), so there are at most $6^k$ words for each $k$; each word outside $c_\Box$ is bounded, per cell, by $v_bK_k(\bar\mu/2)\mathsf A^ke^{-\bar\mu w_b/4}$ (as $K_k(\bar\mu)\le K_k(\bar\mu/2)$ and all rates are at least $\bar\mu w_b/4$), and the words made of $V_{J''}$ alone add up to $\kappa_{\mu_P}(V_{J''};k)$.
\end{proof}


\begin{theorem}[the cumulant bound for the step Gaussian]\label{rot:thm:R1}
In the setting \textup{(R1)}--\textup{(R4)} above (any $\mathbb{L}$ among $\Z^3$, the cubic tori and the skew tori, any finite $\Lambda\subset\mathbb{L}$, the Euclidean or the two-component cutoff) with $b_*$ the threshold fixed before \Cref{rot:lem:R1stg}, there are $S<\infty$ and $\rho_3>0$, all three depending only on $(n,\bar n,D,\varsigma,\gamma_0,c_P,\delta_P,c_K)$, such that for $b\ge b_*$ and $\mathsf A\,b^{D+6}\le1$
\begin{align*}
\log\int\chi_b\,e^{V}d\mu_P&\le\sum_{k=1}^{\bar n}\frac{\kappa_{\mu_P}(V;k)}{k!}+S|J|\big[(\mathsf A\,b^{D+6})^{\bar n+1}+e^{-\rho_3b^{3/2}}\big],\\
\log\int\chi_b\,e^{V}d\mu_P&\ge\sum_{k=1}^{\bar n}\frac{\kappa_{\mu_P}(V;k)}{k!}-S|J|\big[(\mathsf A\,b^{D+6})^{\bar n+1}+e^{-\rho_3b^{3/2}}\big]-4n|\Lambda|\,e^{-\gamma_0\gamma_K^{8}b^2/(2n)},
\end{align*}
where $\kappa_{\mu_P}(V;k)$ are the truncated expectations of $V$ under $\mu_P$ without cutoff and $\gamma_K=\min\{\frac12,(2c_K\mathsf M_1)^{-1}\}$ is the constant of \eqref{rot:eq:R1const}. The constants do not depend on $\Lambda$, on $J$, on the torus, or on $P$ and $V$ otherwise, and no translation invariance is used. For the step Gaussians of both flows, at every step $j$ below the stopping scale, on $T^{(j)}$, on a principal block of it or on a Neumann sub-domain that is a union of $j$-blocks, in the $X$-frame and in the $Y$-frame steps alike, the hypotheses hold with $\gamma_0$ Ba{\l}aban's coercivity constant \cite[(2.33)]{Ba1}, $\delta_P$ the decay rate of the step precision divided by $\sqrt3$, and $c_P$ a constant of $(a,M)$.
\end{theorem}
\begin{proof}
\emph{The hypotheses for the step Gaussians.} The lower bound $\Delta^{(j)}+a5^{-2}Q^*Q\ge\gamma_0$ is \cite[(2.33)]{Ba1} (for the hybrid flow by \Cref{rot:prop:BD}(a)), with constants independent of $A$, $j$ and the domain; this is (R1), which needs no upper bound. The kernel bound (R2) is the exponential decay of $(M_j)^{-1}$ by Combes--Thomas conjugation (for the hybrid flow, the kernel bound of \Cref{rot:prop:BD}(a)) (whose first-order hopping part is anti-Hermitian and whose remainder has norm $\le3d\delta^2$ uniformly in the spacing), with $\dist(B(u),B(v))\ge|u-v|-\sqrt3$, and the block term $a5^{-2}Q^*Q$ has range $\sqrt3\cdot5$; only the off-diagonal entries are constrained, so a diagonal term (such as the shell mass of the shell model) leaves the constants unchanged; Ba{\l}aban's bounds are stated in the sup metric, and $|x-y|\le\sqrt3\,d^*(x,y)$ converts them to the quotient Euclidean distance at the rate divided by $\sqrt3$; principal blocks and Neumann sub-domains inherit (R1) and (R2) as sub-matrices. The interaction at step $j$ is a real polynomial of fixed degree in the fluctuation variables (the analytic remainder vertices of order $\bar n+1$ are taken out of the exponent on the support of the cutoff as $e^{\pm C\ell_j^{7-\beta}|\Lambda|}$, as \cite[(2.31)]{Ba2} does), with $\mathsf A=\|V\|_\varsigma\le s_0\,O(1)\ell_j^{\beta_0}$ by \Cref{rot:cor:OF} at $\varsigma$ half of Ba{\l}aban's coefficient decay rate: a coefficient bound $O(1)\ell_j^{\beta_0}e^{-\delta_0d(x_1,\dots,x_p)}$, with $d$ the length of the shortest graph connecting the points \cite[(3.58)]{Ba1}, gives $\|V\|_{\delta_0/2}\le s_0O(1)\ell_j^{\beta_0}$ with $s_0\defeq\sum_{p\le D}p\,n^p\,\tilde c(\delta_0/(2p))^{p-1}$, because $d(x_1,\dots,x_p)\ge(p-1)^{-1}\sum_{i\ge2}|x_i-x_1|$; this is (R3). The cutoff is Ba{\l}aban's, a product over sites of $\mathbf 1\{|A'(x)|\le b\}\mathbf 1\{|\phi'(x)|\le b\}$ \cite[(3.43), (3.50)]{Ba1}, which is (R4) with $c_K=\sqrt2$. Finally $\mathsf Ab^{D+6}\le1$ and $b\ge b_*$ hold at $b=p_1(\ell_j)=b_1(1+\log\ell_j^{-1})^{p_1}$ for $\ell_j$ below Ba{\l}aban's final spacing, since $\ell_j^{\beta_0}$ beats every power of $1+\log\ell_j^{-1}$; the exponent $p_1\ge1$ is a choice within Ba{\l}aban's freedom, under which $e^{-\rho_3b^{3/2}}$ and $e^{-\gamma_0\gamma_K^8b^2/(2n)}$ are $O(\ell_j^{\kappa})$ for every $\kappa$.

% r03_r1_proof: the cumulant bound for the step Gaussian, part V: the thresholds, the small-volume case and the assembly over the four pavements. Input inside the proof of Theorem rot:thm:R1 in sec/r03_fill_K5a.tex.
\emph{The thresholds.} The threshold $b_*$ is the one fixed before \Cref{rot:lem:R1stg}; it depends only on $(n,\bar n,D,\gamma_0,c_P,\delta_P,c_K)$.

\emph{Small volume.} If some pavement consists of a single cell, then $|\Lambda|\le v_b$ by (G4) and no pavement is needed: with $\nu=\mu_P$ (centred), $\epsilon\defeq\mu_P(1-\chi_b)\le\epsilon_b$ by \Cref{rot:lem:R1tail} and $M\defeq\mathsf A|\Lambda|(c_Kb)^D\le\mathsf m_*$, \Cref{rot:lem:R1ext} moves the cutoff into the exponent at a cost of at most $2e^{2M}\epsilon$, \Cref{rot:lem:R1cum} expands to order $\bar n$ at a cost of $(2M)^{\bar n+1}$, and \Cref{rot:lem:R1chi} replaces $\kappa_{\mu_P}(V\chi_b;k)$ by $\kappa_{\mu_P}(V;k)$ at a cost of $k\,k!\,4^k(\mathsf A|\Lambda|)^k(2b)^{kD}\epsilon^{1/2}/k!$ for each $k$; the total is at most $E_*\le|J|E_*$ in both directions. (If $J=\emptyset$ then $V=0$ and the statement reduces to $0\ge\log\mu_P(\chi_b)\ge-4n|\Lambda|e^{-\gamma_0b^2/(2n)}$, which is \Cref{rot:lem:R1gci}.)

\emph{Assembly.} Otherwise iterate \Cref{rot:lem:R1stg} over $i=0,1,2,3$ with $J_0\defeq J$, $J_{i+1}\defeq J_i\cap\{\lambda_i<2.5w_b\}$ (the set $J''$ of stage $i$), and the thresholds $\beta_i\defeq\gamma_K^ib$ for the lower bound and $\beta_i\defeq\gamma_K^{-i}b$ for the upper bound, all of which lie in $[b_{\min},b_{\max}]$ together with $\beta_3/\gamma_K$. By (G3), $J_4\subset J\cap\bigcap_i\{\lambda_i<2.5w_b\}=\emptyset$, so $V_{J_4}=0$ and the cumulant increments of the four stages telescope to $\sum_{k\le\bar n}\kappa_{\mu_P}(V_J;k)/k!$; the numbers of $J$-cells add up to $\sum_iN_i\le4|J|$. The terminal factor is $\Phi(\gamma_K^4b;\emptyset)=\mu_P(\chi_{b_{\min}})\ge\exp\big(-4n|\Lambda|e^{-\gamma_0\gamma_K^8b^2/(2n)}\big)$ in the lower bound, by \Cref{rot:lem:R1gci} ($\mu_P$ is centred with component variances at most $\|C\|\le\gamma_0^{-1}$), and $\Phi(\gamma_K^{-4}b;\emptyset)=\mu_P(\chi_{b_{\max}})\le1$ in the upper bound. Hence
\[
\log\int\chi_be^Vd\mu_P-\sum_{k=1}^{\bar n}\frac{\kappa_{\mu_P}(V;k)}{k!}\ \le\ 4|J|E_*(b,\mathsf A),\qquad\log\int\chi_be^Vd\mu_P-\sum_{k=1}^{\bar n}\frac{\kappa_{\mu_P}(V;k)}{k!}\ \ge\ -4|J|E_*(b,\mathsf A)-4n|\Lambda|e^{-\gamma_0\gamma_K^8b^2/(2n)}.
\]

\emph{The constants.} $e_{\mathrm{cum}}=2(2M)^{\bar n+1}=2(2\mathsf m_*)^{\bar n+1}(\mathsf Ab^{D+6})^{\bar n+1}$ exactly. Every other term of $E_*$ is increasing in $\mathsf A$, so under $\mathsf Ab^{D+6}\le1$ it is at most its value at $\mathsf A=b^{-D-6}$; call the sum of these values $Q(b)$. It is a finite sum of polynomials in $b$ times one of $e^{-\varsigma w_b/2}$, $e^{-\delta_Pw_b}$, $e^{-\bar\mu w_b/4}$, $e^{-\mu_1w_b/8}$, $e^{-2\bar\mu w_b}$ and $\epsilon_{b_{\min}}^{1/2}$, each of which is at most a polynomial in $b$ times $e^{-\bar\mu b^{3/2}/8}$ (since $w_b\ge b^{3/2}$, $\bar\mu\le\mu_1\le\delta_P/2$ and $\bar\mu\le\varsigma$). With $\rho_3\defeq\bar\mu/16$ and
\[
S\defeq4\max\Big\{2(2\mathsf m_*)^{\bar n+1},\ \sup_{b\ge b_*}e^{\rho_3b^{3/2}}Q(b)\Big\}<\infty
\]
one has $4|J|E_*(b,\mathsf A)\le S|J|\big[(\mathsf Ab^{D+6})^{\bar n+1}+e^{-\rho_3b^{3/2}}\big]$, which is the statement. $S$, $b_*$, $\rho_3$ and $\gamma_K$ depend only on $(n,\bar n,D,\varsigma,\gamma_0,c_P,\delta_P,c_K)$; nothing depends on $\Lambda$, on $J$, on $\mathbb{L}$ or its period lattice, on $P$ and $V$ beyond the stated bounds, or on any translation invariance, of which none is used; no lower bound on $c_P$ and no upper bound on $P$ enters.

\end{proof}
